% Leçon 35 — Vocabulaire et compréhension de texte : internet — Chinois 1ère A — Cours signé OnBuch+
% Programme officiel MINESEC. Compiler avec : tectonic ZHA35-vocabulaire-internet.tex
\newcommand{\DOCMATIERE}{Chinois}
\newcommand{\DOCNIVEAU}{1ère A}
\newcommand{\DOCMODULE}{Module 5 : Mass médias et télécommunication}
\newcommand{\DOCLECON}{ZHA35}
\newcommand{\DOCTITRE}{Leçon 35 — Vocabulaire et compréhension de texte : internet}
% OnBuch+ — préambule « cours signés » ENRICHI (compilé avec tectonic / XeLaTeX)
% Système visuel « Pop » d'OnBuch+ : fond crème (jamais blanc), bordures encre
% épaisses, ombres « dures » décalées, titres Archivo Black en capitales.
% Version MATHÉMATIQUES : graphes (pgfplots/tikz), boîtes pédagogiques
% (définition, propriété, méthode, attention, le savais-tu, exercice résolu,
% exercice, TP, bilan), page de garde et signature OnBuch+ ancrée en bas.
%
% Macros attendues dans main.tex AVANT \input{preamble} :
%   \DOCMATIERE  \DOCNIVEAU  \DOCMODULE  \DOCLECON (numéro)  \DOCTITRE
\documentclass[11pt]{article}

\usepackage{fontspec}
\usepackage[french]{babel} % français = langue principale ; le chinois via \zh{..} / textebox (xeCJK)
\usepackage{xeCJK}
\usepackage{pifont} % \ding{51} \ding{55} utilisés par les modèles
\usepackage[a4paper,margin=2cm,top=3.4cm,bottom=2.8cm,headheight=1.4cm,footskip=1.3cm]{geometry}
\usepackage{amsmath,amssymb,amsfonts}
\usepackage{mathtools} % \xmapsto, \coloneqq... souvent utilisés par les modèles
\usepackage{cancel} % simplifications barrées (bilans de forces)
% \abs{z} (module, valeur absolue) et \norm{u} (norme), utilisés par les modèles
\providecommand{\abs}{}\let\abs\relax\DeclarePairedDelimiter{\abs}{\lvert}{\rvert}
\providecommand{\norm}{}\let\norm\relax\DeclarePairedDelimiter{\norm}{\lVert}{\rVert}
\usepackage[table]{xcolor} % [table] : fournit aussi \rowcolors (alternance de lignes)
\usepackage{pagecolor}
\usepackage{tikz}
\usetikzlibrary{arrows.meta,positioning,calc,shapes.geometric,shapes.misc,shapes.arrows,decorations.pathmorphing,decorations.pathreplacing,decorations.markings,patterns,shadows,mindmap,trees,fit,backgrounds,matrix,intersections,angles,quotes}
\usepackage{pgfplots}
\usepackage[version=4]{mhchem} % équations nucléaires \ce{^{235}_{92}U}
\usepackage[european,siunitx]{circuitikz} % schémas électriques
\pgfplotsset{compat=1.18}
\usepgfplotslibrary{fillbetween,groupplots} % aires entre courbes (calcul intégral)
\usepackage{enumitem}
\usepackage{fancyhdr}
% [most] charge toutes les bibliothèques tcolorbox (listings, minted, raster,
% documentation...) et gonfle inutilement la mémoire TeX sur un document long
% (~300 boîtes) : on ne charge que ce qui est réellement utilisé.
\usepackage[skins,breakable]{tcolorbox}
\usepackage{graphicx}
\usepackage{colortbl}
\usepackage{booktabs}
\usepackage{tabularx}
\usepackage{diagbox} % case d'en-tête diagonale des tableaux à double entrée
% Colonnes extensibles centrée / gauche / droite (C, L, R), souvent utilisées par les modèles
\newcolumntype{C}{>{\centering\arraybackslash}X}
\newcolumntype{L}{>{\raggedright\arraybackslash}X}
\newcolumntype{R}{>{\raggedleft\arraybackslash}X}
\usepackage{array}
\usepackage{multicol}
\usepackage{multirow}
\usepackage{siunitx}
% parse-numbers=false : accepte aussi \SI{2,5 \times 10^{-3}}{...}
\sisetup{locale=FR,output-decimal-marker={,},group-minimum-digits=4,parse-numbers=false}
\DeclareSIUnit\ohm{\ensuremath{\symup{\Omega}}} % U+2126 absent de la police
\DeclareSIUnit\division{div} % réglages d'oscilloscope (ms/div, V/div)
\DeclareSIUnit\tr{tr} % tours (vitesse de rotation, ex. \SI{1500}{\tr\per\minute})
\DeclareSIUnit\molar{M} % concentration molaire (ex. \SI{3}{\molar}, \SI{1}{\micro\molar})
\DeclareSIUnit\an{an} % année (le modèle confond parfois avec \year, un compteur TeX) — ex. \SI{5}{\kilo\meter\per\an}
\DeclareSIUnit\FCFA{FCFA} % franc CFA (ex. \SI{500}{\FCFA\per\kilo\gram})
\DeclareSIUnit\degreeBrix{\textdegree Brix} % teneur en sucre d'un sirop/jus (réfractométrie)
\DeclareSIUnit\month{mois} % le modèle utilise parfois \month, un compteur TeX, comme unité de durée
\usepackage{qrcode}
% \degree : commande fréquemment utilisée par les modèles pour un angle ou une
% température en dehors de \SI{}{} (non fournie par les paquets chargés).
\providecommand{\degree}{^{\circ}}
% \qed (fin de démonstration), utilisé par les modèles sans amsthm
\providecommand{\qed}{\hfill$\square$}
% \barre : double barre des valeurs interdites dans un tableau de variations (tkz-tab)
\providecommand{\barre}{\ensuremath{\|}}
% environnement proof (amsthm non chargé) utilisé par les modèles
\ifdefined\proof\else\newenvironment{proof}[1][Démonstration]{\par\noindent\textit{#1.}\ }{\qed\par}\fi
\usepackage{lastpage}
\usepackage{hyperref}
\hypersetup{hidelinks}

% ============================================================
% Palette « Pop » OnBuch+ — mêmes hex que la classe P (plus_ui.dart)
% ============================================================
\definecolor{poporange}{HTML}{F59321}
\definecolor{popdark}{HTML}{E07A0C}
\definecolor{popcream}{HTML}{FFF6E8}
\definecolor{popink}{HTML}{1C1714}
\definecolor{poppurple}{HTML}{7A5AE0}
\definecolor{popgreen}{HTML}{2E9E6A}
\definecolor{poppink}{HTML}{E0457B}
\definecolor{popblue}{HTML}{2F7DE1}
\definecolor{popgold}{HTML}{C98A12}
\definecolor{popmuted}{HTML}{8A7F76}
\definecolor{popline}{HTML}{EADFD2}
\definecolor{popgray}{HTML}{8A7F76}
\colorlet{popgrey}{popgray}
% Alias tolérants (le modèle nomme parfois les couleurs différemment)
\colorlet{popwhite}{white}
\colorlet{popblack}{popink}
\colorlet{popred}{poppink}
\colorlet{popyellow}{popgold}
% Teintes claires (fonds de tableaux/encadrés) — le modèle nomme parfois
% les couleurs avec un suffixe L (ex. popblueL) sans les avoir définies.
\colorlet{poporangeL}{poporange!20}
\colorlet{popdarkL}{popdark!20}
\colorlet{poppurpleL}{poppurple!20}
\colorlet{popgreenL}{popgreen!20}
\colorlet{poppinkL}{poppink!20}
\colorlet{popblueL}{popblue!20}
\colorlet{popgoldL}{popgold!20}
\colorlet{popredL}{popred!20}
\colorlet{popgrayL}{popgray!20}
% Teintes claires pour fonds de tableaux / zones
\colorlet{popblueL}{popblue!10!white}
\colorlet{poporangeL}{poporange!14!white}
\colorlet{popgreenL}{popgreen!12!white}
\colorlet{poppurpleL}{poppurple!12!white}
\colorlet{poppinkL}{poppink!12!white}
\colorlet{popgoldL}{popgold!14!white}

\pagecolor{popcream}

% ============================================================
% Typographie
% ============================================================
\defaultfontfeatures{Ligatures=TeX}
\setmainfont{PlusJakartaSans-Medium.ttf}[Path=../../../fonts/,BoldFont=PlusJakartaSans-Bold.ttf]
% Symboles, API et emojis absents de la police principale : repli sur DejaVu Sans (caractères actifs)
\newfontface\symfont{DejaVuSans.ttf}[Path=../../../fonts/]
\makeatletter
\newcommand\symactive[1]{\catcode#1=13 \begingroup\lccode`\~=#1 \lowercase{\endgroup\def~}{{\symfont\char#1}}}
\newcommand\symblank[1]{\catcode#1=13 \begingroup\lccode`\~=#1 \lowercase{\endgroup\def~}{}}
\newcommand\symhyph[1]{\catcode#1=13 \begingroup\lccode`\~=#1 \lowercase{\endgroup\def~}{\mbox{-}}}
\newcount\symc
\newcommand\symrange[2]{\symc=#1 \@whilenum\symc<#2 \do{\expandafter\symactive\expandafter{\the\symc}\advance\symc by 1 }}
\newcommand\symblankrange[2]{\symc=#1 \@whilenum\symc<#2 \do{\expandafter\symblank\expandafter{\the\symc}\advance\symc by 1 }}
\makeatother
\symrange{"0250}{"02FF}\symactive{"2713}\symactive{"2714}\symactive{"2717}\symactive{"2718}\symactive{"2610}\symactive{"2611}\symactive{"2612}
\symrange{"2600}{"27BF}
\catcode"2705=13 \begingroup\lccode`\~="2705 \lowercase{\endgroup\def~}{{\symfont\char"2713}}\symblank{"26BD}\symblank{"26A1}
\symrange{"2460}{"257F}\symactive{"223C}\symactive{"2248}\symactive{"2260}
\symactive{"01F8}\symactive{"01F9}\symactive{"1E3E}\symactive{"1E3F}
\symhyph{"2011}\symhyph{"2E1A}\symblankrange{"1F300}{"1FAFF}\symblank{"FE0F}
% Caractères chinois : WenQuanYi Zen Hei (sous-ensemble GB2312 + pinyin), gras/italique simulés
\setCJKmainfont{WenQuanYiZenHei-subset.ttf}[Path=../../../fonts/,AutoFakeBold=3,AutoFakeSlant=0.2]
\xeCJKsetup{CJKspace=false}
% Pinyin : voyelles à caron (ǎ ǐ ǒ ǔ ǖ ǘ ǚ ǜ) absentes de la police principale -> caractères actifs
\newfontface\pyfont{WenQuanYiZenHei-subset.ttf}[Path=../../../fonts/]
\newcommand\pyactive[1]{\catcode#1=13 \begingroup\lccode`\~=#1 \lowercase{\endgroup\def~}{{\pyfont\char#1}}}
\pyactive{"01CD}\pyactive{"01CE}\pyactive{"01CF}\pyactive{"01D0}\pyactive{"01D1}\pyactive{"01D2}\pyactive{"01D3}\pyactive{"01D4}\pyactive{"01D5}\pyactive{"01D6}\pyactive{"01D7}\pyactive{"01D8}\pyactive{"01D9}\pyactive{"01DA}\pyactive{"01DB}\pyactive{"01DC}
% Maths en sans-serif (Fira Math) avec chiffres et lettres droites de Plus
% Jakarta, pour que les formules se fondent dans le texte.
\usepackage[math-style=ISO,bold-style=ISO]{unicode-math}
\setmathfont{FiraMath-Regular.otf}[Path=../../../fonts/]
\setmathfont{PlusJakartaSans-Medium.ttf}[Path=../../../fonts/,range={up/num,up/{Latin,latin},\mathup}]
\usepackage{icomma} % virgule décimale sans espace en mode math
% Caractères mathématiques Unicode absents des polices de texte (Plus Jakarta,
% Archivo Black) : rendus par leur équivalent mathématique, en texte comme en
% formule — sinon ils s'affichent en carré vide (ex. « Arithmétique dans ℤ »).
\usepackage{newunicodechar}
\newunicodechar{ℝ}{\ensuremath{\mathbb{R}}}
\newunicodechar{ℤ}{\ensuremath{\mathbb{Z}}}
\newunicodechar{ℂ}{\ensuremath{\mathbb{C}}}
\newunicodechar{ℕ}{\ensuremath{\mathbb{N}}}
\newunicodechar{ℚ}{\ensuremath{\mathbb{Q}}}
\newunicodechar{𝔻}{\ensuremath{\mathbb{D}}}
\newunicodechar{α}{\ensuremath{\alpha}}
\newunicodechar{β}{\ensuremath{\beta}}
\newunicodechar{γ}{\ensuremath{\gamma}}
\newunicodechar{δ}{\ensuremath{\delta}}
\newunicodechar{ε}{\ensuremath{\varepsilon}}
\newunicodechar{θ}{\ensuremath{\theta}}
\newunicodechar{λ}{\ensuremath{\lambda}}
\newunicodechar{μ}{\ensuremath{\mu}}
\newunicodechar{σ}{\ensuremath{\sigma}}
\newunicodechar{τ}{\ensuremath{\tau}}
\newunicodechar{φ}{\ensuremath{\varphi}}
\newunicodechar{ψ}{\ensuremath{\psi}}
\newunicodechar{ω}{\ensuremath{\omega}}
\newunicodechar{Σ}{\ensuremath{\Sigma}}
% majuscules produites par \MakeUppercase dans les titres (α-aminés -> Α)
\newunicodechar{Α}{\ensuremath{\alpha}}
\newunicodechar{Β}{\ensuremath{\beta}}
\newunicodechar{Γ}{\ensuremath{\Gamma}}
\newunicodechar{Δ}{\ensuremath{\Delta}}
\newunicodechar{Ε}{\ensuremath{\varepsilon}}
\newunicodechar{Θ}{\ensuremath{\Theta}}
\newunicodechar{Λ}{\ensuremath{\Lambda}}
\newunicodechar{Μ}{\ensuremath{\mu}}
\newunicodechar{Τ}{\ensuremath{\tau}}
\newunicodechar{Φ}{\ensuremath{\Phi}}
\newunicodechar{Ψ}{\ensuremath{\Psi}}
\newunicodechar{Ω}{\ensuremath{\Omega}}
\newunicodechar{↦}{\ensuremath{\mapsto}}
\newunicodechar{∈}{\ensuremath{\in}}
\newunicodechar{∉}{\ensuremath{\notin}}
\newunicodechar{∧}{\ensuremath{\wedge}}
\newunicodechar{⇔}{\ensuremath{\Leftrightarrow}}
\newunicodechar{⟺}{\ensuremath{\Longleftrightarrow}}
\newunicodechar{⇒}{\ensuremath{\Rightarrow}}
\newunicodechar{⟹}{\ensuremath{\Longrightarrow}}
\newunicodechar{∘}{\ensuremath{\circ}}
\newunicodechar{∪}{\ensuremath{\cup}}
\newunicodechar{∩}{\ensuremath{\cap}}
\newunicodechar{≡}{\ensuremath{\equiv}}
\newunicodechar{⊂}{\ensuremath{\subset}}
\newunicodechar{⊥}{\ensuremath{\perp}}
\newunicodechar{∃}{\ensuremath{\exists}}
\newunicodechar{∀}{\ensuremath{\forall}}
\newunicodechar{∛}{\ensuremath{\sqrt[3]{\,}}}
\newunicodechar{∜}{\ensuremath{\sqrt[4]{\,}}}
\newunicodechar{⃗}{\ensuremath{{}^{\to}}}
\newunicodechar{ⁿ}{\ifmmode^{n}\else\textsuperscript{\ensuremath{n}}\fi}
\newunicodechar{ˣ}{\ifmmode^{x}\else\textsuperscript{\ensuremath{x}}\fi}
\newunicodechar{ᵇ}{\ifmmode^{b}\else\textsuperscript{\ensuremath{b}}\fi}
\newunicodechar{ʳ}{\ifmmode^{r}\else\textsuperscript{\ensuremath{r}}\fi}
\newunicodechar{ᵘ}{\ifmmode^{u}\else\textsuperscript{\ensuremath{u}}\fi}
\newunicodechar{ʸ}{\ifmmode^{y}\else\textsuperscript{\ensuremath{y}}\fi}
\newunicodechar{ᵃ}{\ifmmode^{a}\else\textsuperscript{\ensuremath{a}}\fi}
\newunicodechar{ᵉ}{\ifmmode^{e}\else\textsuperscript{\ensuremath{e}}\fi}
\newunicodechar{ᵏ}{\ifmmode^{k}\else\textsuperscript{\ensuremath{k}}\fi}
\newunicodechar{ᵖ}{\ifmmode^{p}\else\textsuperscript{\ensuremath{p}}\fi}
\newunicodechar{ᴺ}{\ifmmode^{N}\else\textsuperscript{\ensuremath{N}}\fi}
\newunicodechar{ᶜ}{\ifmmode^{c}\else\textsuperscript{\ensuremath{c}}\fi}
\newunicodechar{ʰ}{\ifmmode^{h}\else\textsuperscript{\ensuremath{h}}\fi}
\newunicodechar{ᵠ}{\ifmmode^{q}\else\textsuperscript{\ensuremath{q}}\fi}
\newunicodechar{ₙ}{\ifmmode_{n}\else\textsubscript{\ensuremath{n}}\fi}
\newunicodechar{ₐ}{\ifmmode_{a}\else\textsubscript{\ensuremath{a}}\fi}
\newunicodechar{ᵢ}{\ifmmode_{i}\else\textsubscript{\ensuremath{i}}\fi}
\newunicodechar{ᵣ}{\ifmmode_{r}\else\textsubscript{\ensuremath{r}}\fi}
\newunicodechar{ₖ}{\ifmmode_{k}\else\textsubscript{\ensuremath{k}}\fi}
\newunicodechar{ᵦ}{\ifmmode_{\beta}\else\textsubscript{\ensuremath{\beta}}\fi}
\newunicodechar{ₓ}{\ifmmode_{x}\else\textsubscript{\ensuremath{x}}\fi}
\newfontfamily\popdisplay{ArchivoBlack-Regular.ttf}[Path=../../../fonts/]
\newfontfamily\poplogo{SpaceGrotesk-Bold.ttf}[Path=../../../fonts/]
\newfontfamily\popsemi{PlusJakartaSans-SemiBold.ttf}[Path=../../../fonts/]
\newfontfamily\popxbold{PlusJakartaSans-ExtraBold.ttf}[Path=../../../fonts/]
\newfontfamily\popxboldls{PlusJakartaSans-ExtraBold.ttf}[Path=../../../fonts/,LetterSpace=3.0]

\setlist[enumerate]{leftmargin=1.7em,itemsep=5pt,topsep=4pt}
\setlist[itemize]{leftmargin=1.7em,itemsep=4pt,topsep=4pt,label={\color{poporange}\textbullet}}
\setlength{\parindent}{0pt}
\setlength{\parskip}{5pt}
\renewcommand{\arraystretch}{1.25}

% pgfplots : style Pop par défaut
\pgfplotsset{
  every axis/.append style={axis line style={popink,line width=1pt},tick style={popink},
    grid style={popline,line width=0.5pt},label style={font=\small},tick label style={font=\footnotesize}},
  popcurve/.style={poporange,line width=1.8pt,no markers,smooth},
}
% Même style disponible dans un \draw TikZ simple (hors pgfplots)
\tikzset{popcurve/.style={poporange,line width=1.8pt}}
\tikzset{popgrid/.style={popline,very thin}}
\colorlet{popcurve}{poporange} % le nom du style est parfois utilisé comme couleur

% ============================================================
% Pill / squiggle
% ============================================================
\newcommand{\poppill}[3]{%
  \tikz[baseline=-0.55ex]\node[draw=popink,line width=1.2pt,rounded corners=2.6pt,%
    fill=#2,inner xsep=6.5pt,inner ysep=3pt]{%
    \popxboldls\fontsize{7}{7}\selectfont\color{#3}\MakeUppercase{#1}};%
}
\newcommand{\popsquiggle}[1]{%
  \tikz[baseline=-0.4ex]{
    \draw[#1,line width=2.6pt,line cap=round]
      plot [smooth] coordinates {(0,0.09) (0.34,0.01) (0.68,0.21) (1.02,0.01) (1.36,0.10)};
  }%
}
% Mot-clé surligné (marqueur orange sous le texte)
\newcommand{\cle}[1]{{\popxbold\color{popdark}#1}}

% ============================================================
% En-tête / pied
% ============================================================
\pagestyle{fancy}
\fancyhf{}
\renewcommand{\headrulewidth}{0pt}
\renewcommand{\footrulewidth}{0pt}
\lhead{\raisebox{-0.15ex}{\poplogo\fontsize{13}{13}\selectfont\color{popink}OnBuch\color{poporange}+}}
\rhead{\poppill{\DOCMATIERE\ \textbullet\ \DOCNIVEAU\ \textbullet\ Leçon \DOCLECON}{white}{popink}}
\lfoot{\popxbold\fontsize{7.6}{7.6}\selectfont\color{popmuted}\MakeUppercase{Cours signé OnBuch+}\ \textbullet\ {\popsemi\DOCTITRE}}
\rfoot{\tikz[baseline=-0.6ex]\node[rounded corners=8.5pt,fill=popink,minimum height=17pt,minimum width=17pt,inner xsep=5pt,inner sep=0]{\popxbold\fontsize{8}{8}\selectfont\color{popcream}\thepage\,/\,\pageref*{LastPage}};}

% ============================================================
% Titres
% ============================================================
\newcommand{\chaptitre}[2]{%
  \begin{tcolorbox}[enhanced,colback=poporange,colframe=popink,boxrule=2.6pt,arc=11pt,
      drop shadow=popink,left=15pt,right=15pt,top=12pt,bottom=11pt,before skip=3pt,after skip=18pt]
    \poppill{#2}{popink}{popcream}\par\vspace{9pt}
    {\raggedright\hyphenpenalty=10000\exhyphenpenalty=10000\popdisplay\fontsize{22}{24}\selectfont\color{popcream}\MakeUppercase{#1}\par}
  \end{tcolorbox}
}
% Section numérotée : pastille orange avec le numéro + titre Archivo
\newcounter{popsec}
\newcommand{\coursec}[1]{%
  \stepcounter{popsec}\setcounter{popsub}{0}%
  \par\vspace{16pt}\noindent
  \begin{minipage}{\linewidth}
  \tikz[baseline=-0.7ex]\node[draw=popink,line width=1.6pt,fill=poporange,rounded corners=4pt,minimum size=22pt,inner sep=0]{\popdisplay\fontsize{12}{12}\selectfont\color{popcream}\thepopsec};%
  \hspace{8pt}{\popdisplay\fontsize{15}{17}\selectfont\color{popink}\MakeUppercase{#1}}\par
  \vspace{4pt}\noindent{\color{poporange}\rule{36pt}{3.6pt}}
  \end{minipage}\par\nopagebreak\vspace{8pt}
}
\newcounter{popsub}
\newcommand{\courssub}[1]{%
  \stepcounter{popsub}%
  \par\vspace{9pt}\noindent{\popxbold\fontsize{11.5}{13}\selectfont\color{popink}{\color{poporange}\thepopsec.\thepopsub}\hspace{6pt}#1}\par\nopagebreak\vspace{4pt}
}

% ============================================================
% Boîtes pédagogiques — même recette : fond blanc, bordure épaisse colorée,
% ombre dure encre, badge plein. Titre optionnel [#1].
% ============================================================
\tcbset{popbase/.style={breakable,enhanced,colback=white,boxrule=2.3pt,arc=9pt,
  shadow={3pt}{-3pt}{0pt}{popink},left=14pt,right=14pt,top=11pt,bottom=11pt,before skip=11pt,after skip=15pt,
  fontupper=\popsemi\fontsize{10.6}{14.5}\selectfont\color{popink},
  fontlower=\popsemi\fontsize{10.6}{14.5}\selectfont\color{popink}}}
% Coche dessinée (le glyphe ✓ n'existe pas dans les polices du document)
\newcommand{\popcheck}[1]{\tikz[baseline=-0.5ex]\draw[#1,line width=1.6pt,line cap=round,line join=round](-0.13,0.0)--(-0.04,-0.09)--(0.13,0.1);}
\providecommand{\coche}{\popcheck{popgreen}}
\providecommand{\resultat}[1]{\par\smallskip\noindent\fcolorbox{popgreen}{popgreenL}{\parbox{\dimexpr\linewidth-2\fboxsep-2\fboxrule\relax}{\textbf{Résultat.} #1}}\par\smallskip}
\newcommand{\popboxhead}[3]{\poppill{#1}{#2}{white}\ifx\relax#3\relax\else\hspace{6pt}{\popxbold\fontsize{10.6}{12}\selectfont\color{popink}#3}\fi\par\vspace{7pt}}

\newtcolorbox{aretenir}[1][]{popbase,colframe=popgreen,before upper={\popboxhead{À retenir}{popgreen}{#1}}}
% Texte en chinois (document de lecture, dialogue, extrait) : cadre
% d'étude. Un titre optionnel (source, auteur) s'affiche dans l'en-tête.
\newtcolorbox{textebox}[1][]{popbase,colframe=popblue,before upper={\popboxhead{文本 · Texte}{popblue}{#1}},after upper={}}
% Marques ✓ / ✗ (forme correcte / fautive) souvent utilisées par les modèles
\providecommand{\cross}{\textcolor{poppink}{\ensuremath{\times}}}
\providecommand{\xmark}{\cross}\providecommand{\crossmark}{\cross}\providecommand{\cmark}{\ensuremath{\checkmark}}
% Ligne de réponse / justification à compléter (inventée par les modèles)
\providecommand{\justification}[1]{\par\noindent\textit{Justification~:} \dotfill\par}
% \textipa (package tipa non chargé) : la police ne couvre pas l'API -> simple passage
\providecommand{\textipa}[1]{#1}
% Soulignement (ulem non chargé) : \uline, \ul
\providecommand{\uline}[1]{\underline{#1}}\providecommand{\ul}[1]{\underline{#1}}
% \glossaire{\item ...} inventée par les modèles : liste de glossaire
\providecommand{\glossaire}[1]{\begin{itemize}#1\end{itemize}}
% \alert (beamer) inventée par les modèles : texte mis en évidence
\providecommand{\alert}[1]{\textbf{\textcolor{poppink}{#1}}}
% variantes ASCII d'ordinaux écrites en macro
\providecommand{\iere}{\textsuperscript{re}}\providecommand{\ieme}{\textsuperscript{e}}\providecommand{\ere}{\textsuperscript{re}}\providecommand{\eme}{\textsuperscript{e}}\providecommand{\er}{\textsuperscript{er}}\providecommand{\ie}{\textsuperscript{e}}
% Ordinaux français écrits en macro par les modèles : 1\ière, 3\ième
\newcommand{\ière}{\textsuperscript{re}}\newcommand{\ième}{\textsuperscript{e}}
% \exemple (inventée par les modèles) : étiquette « Exemple : »
\providecommand{\exemple}{\textbf{Exemple~:}\ }
% \square utilisable aussi hors mode math (cases à cocher)
\AtBeginDocument{\let\squareMath\square\renewcommand{\square}{\ensuremath{\squareMath}}}
% Mot ou phrase en chinois à l'intérieur d'un paragraphe français
\newcommand{\zh}[1]{\ifmmode\text{#1}\else{#1}\fi}
\let\eng\zh \let\esp\zh \let\all\zh % alias tolérés
% flèches d'intonation inventées par les modèles
\providecommand{\textsearrow}{}\renewcommand{\textsearrow}{\ensuremath{\searrow}}\providecommand{\textnearrow}{}\renewcommand{\textnearrow}{\ensuremath{\nearrow}}
% \fr{..} : mot français (inventé par les modèles)
\providecommand{\fr}[1]{\textit{#1}}
% zhbox : encadré de texte chinois inventé par les modèles
\newenvironment{zhbox}{\begin{quote}}{\end{quote}}
% \courssubsub : titre de niveau 3 inventé par les modèles
\providecommand{\courssubsub}[1]{\par\medskip\noindent\textbf{#1}\par\smallskip}
% \zhbf : texte chinois en gras inventé par les modèles
\providecommand{\zhbf}[1]{\textbf{#1}}
% \vs : « versus » inventé par les modèles
\providecommand{\vs}{\textit{vs}\ }
% \blank : case à compléter inventée par les modèles
\providecommand{\blank}[1][]{\underline{\hspace{1.6em}}}
\newtcolorbox{exemplebox}[1][]{popbase,colframe=poporange,before upper={\popboxhead{Exemple}{poporange}{#1}}}
\newtcolorbox{definition}[1][]{popbase,colframe=popblue,before upper={\popboxhead{Définition}{popblue}{#1}}}
\newtcolorbox{propriete}[1][]{popbase,colframe=poppurple,before upper={\popboxhead{Propriété}{poppurple}{#1}}}
\newtcolorbox{methode}[1][]{popbase,colframe=popgold,before upper={\popboxhead{Méthode}{popgold}{#1}}}
\newtcolorbox{attention}[1][]{popbase,colframe=poppink,before upper={\popboxhead{Attention}{poppink}{#1}}}
\newtcolorbox{experience}[1][]{popbase,colframe=popblue,colback=popblueL,before upper={\popboxhead{Expérience}{popblue}{#1}}}
% « Le savais-tu ? » : carte violette pleine (contexte, culture, Cameroun)
\newtcolorbox{savaistu}[1][]{popbase,colback=poppurple,colframe=popink,
  fontupper=\popsemi\fontsize{10.4}{14.2}\selectfont\color{white},
  before upper={\poppill{Le savais-tu\,?}{white}{poppurple}\ifx\relax#1\relax\else\hspace{6pt}{\popxbold\color{white}#1}\fi\par\vspace{7pt}}}
% Exercice résolu : énoncé en haut, \tcblower puis la solution sur fond crème
\newcounter{popexo}\newcounter{popexr}
\newtcolorbox{exoresolu}[1][]{popbase,colframe=poporange,colbacklower=popcream,
  segmentation style={popink,line width=1.2pt,dashed},
  before upper={\stepcounter{popexr}\popboxhead{Exercice résolu \thepopexr}{poporange}{#1}},
  before lower={\poppill{Solution}{popink}{popcream}\par\vspace{6pt}}}
% Exercice (non résolu en place) — la correction est dans la partie Corrigés
\newtcolorbox{exercice}[1][]{popbase,colframe=popink,
  before upper={\stepcounter{popexo}\popboxhead{Exercice \thepopexo}{popink}{#1}}}
% Corrigé
\newtcolorbox{corrige}[1][]{popbase,colframe=popgreen,colback=popgreenL,
  before upper={\popboxhead{Corrigé}{popgreen}{#1}}}
% Objectifs / prérequis / bilan
\newtcolorbox{objectifs}{popbase,colframe=popink,colback=poporangeL,
  before upper={\popboxhead{Objectifs de la leçon}{popink}{}}}
\newtcolorbox{prerequis}{popbase,colframe=popink,colback=popblueL,
  before upper={\popboxhead{Prérequis}{popblue}{}}}
\newtcolorbox{bilan}{popbase,colframe=popink,colback=white,boxrule=2.8pt,
  before upper={\popboxhead{Fiche bilan}{poporange}{L'essentiel à connaître pour l'examen}}}
% Figure encadrée avec légende
\newtcolorbox{popfigure}[1][]{enhanced,breakable=false,colback=white,colframe=popline,boxrule=1.4pt,arc=8pt,
  left=10pt,right=10pt,top=10pt,bottom=8pt,before skip=10pt,after skip=12pt,halign=center,
  lower separated=false,halign lower=center,
  fontlower=\popsemi\fontsize{9}{11.5}\selectfont\color{popmuted},
  #1}
\newcommand{\legende}[1]{\par\vspace{6pt}{\popsemi\fontsize{9}{11.5}\selectfont\color{popmuted}#1}}

% ============================================================
% Signature OnBuch+
% ============================================================
\newcommand{\poplogoinline}[1]{{\poplogo\fontsize{#1}{#1}\selectfont\color{popink}OnBuch\color{poporange}+}}
\newcommand{\signatureonbuch}{%
  \begin{tcolorbox}[enhanced,colback=popink,colframe=popink,boxrule=2.2pt,arc=10pt,
      drop shadow=poporange,left=14pt,right=14pt,top=10pt,bottom=10pt,before skip=0pt,after skip=0pt]
    \tikz[baseline=-0.7ex]\node[circle,fill=popgreen,draw=popcream,line width=1.3pt,minimum size=22pt,inner sep=0]{\popcheck{white}};%
    \hspace{9pt}\begin{minipage}[c]{0.62\linewidth}
      {\popxbold\fontsize{12}{13}\selectfont\color{popcream}Signé\ \textbullet\ {\poplogo OnBuch\color{poporange}+}}\par\vspace{2pt}
      {\raggedright\popsemi\fontsize{8}{10.5}\selectfont\color{popline}\MakeUppercase{Équipe pédagogique OnBuch+}\\\MakeUppercase{Conforme au programme officiel MINESEC}\par}
    \end{minipage}\hfill
    {\poplogo\fontsize{22}{22}\selectfont\color{popcream}OnBuch\color{poporange}+}
  \end{tcolorbox}%
}

% ============================================================
% Page de garde
% #1 titre · #2 matière · #3 niveau · #4 module · #5 n° leçon · #6 durée ·
% #7 description / objectifs (texte ou itemize)
% ============================================================
\newcommand{\pagegarde}[7]{%
  \thispagestyle{empty}
  \newgeometry{a4paper,margin=1.7cm,top=1.6cm,bottom=1.4cm}%
  \begin{tikzpicture}[remember picture,overlay]
    \fill[poporange!18] ([xshift=-3.2cm,yshift=-3.6cm]current page.north east) circle (4.6cm);
    \fill[poppurple!14] ([xshift=2.4cm,yshift=7.6cm]current page.south west) circle (3.8cm);
  \end{tikzpicture}%
  {\poplogo\fontsize{30}{30}\selectfont\color{popink}OnBuch\color{poporange}+}\hfill
  \raisebox{0.6ex}{\poppill{Cours signé \textbullet\ Premium}{popink}{popcream}}\par
  \vspace{1pt}\popsquiggle{poppurple}\par
  \vspace{8pt}
  \begin{tcolorbox}[enhanced,colback=poporange,colframe=popink,boxrule=2.8pt,arc=13pt,
      drop shadow=popink,left=18pt,right=18pt,top=12pt,bottom=13pt,after skip=10pt]
    \poppill{#2 \textbullet\ #3}{popink}{popcream}\hspace{5pt}\poppill{#4}{white}{popink}\par\vspace{8pt}
    {\popdisplay\fontsize{14}{14}\selectfont\color{popink}LEÇON #5}\par\vspace{4pt}
    {\raggedright\hyphenpenalty=10000\exhyphenpenalty=10000\popdisplay\fontsize{25}{27}\selectfont\color{popcream}\MakeUppercase{#1}\par}
    \vspace{8pt}
    {\popxbold\fontsize{9.5}{9.5}\selectfont\color{popink}\MakeUppercase{Durée conseillée : #6}}
  \end{tcolorbox}
  \begin{tcolorbox}[enhanced,colback=white,colframe=popink,boxrule=2.2pt,arc=10pt,
      drop shadow=popink,left=16pt,right=16pt,top=10pt,bottom=10pt,after skip=8pt]
    \poppill{Dans cette leçon}{poppurple}{white}\par\vspace{5pt}
    \popsemi\fontsize{9.3}{12.6}\selectfont\color{popink}#7
  \end{tcolorbox}
  \noindent\begin{minipage}[t]{0.487\linewidth}
    \begin{tcolorbox}[enhanced,colback=popgreen,colframe=popink,boxrule=2.2pt,arc=10pt,
        drop shadow=popink,left=11pt,right=11pt,top=10pt,bottom=10pt,height=3.1cm,valign=center]
      \tcbox[colback=white,colframe=popink,boxrule=1.3pt,arc=5pt,boxsep=0pt,nobeforeafter,
        left=3pt,right=3pt,top=3pt,bottom=3pt,valign=center]{\qrcode[height=1.9cm]{https://play.google.com/store/apps/details?id=com.luvvix.onbuch&hl=fr}}%
      \hspace{8pt}\begin{minipage}[c]{0.5\linewidth}
        {\popdisplay\fontsize{11}{12.5}\selectfont\color{white}\MakeUppercase{Télécharge OnBuch}}\par\vspace{4pt}
        {\raggedright\popsemi\fontsize{8.4}{11}\selectfont\color{white}Gratuit \textbullet\ Google Play\\Tuteur IA, annales, concours, résultats\par}
      \end{minipage}
    \end{tcolorbox}
  \end{minipage}\hfill
  \begin{minipage}[t]{0.487\linewidth}
    \begin{tcolorbox}[enhanced,colback=poppurple,colframe=popink,boxrule=2.2pt,arc=10pt,
        drop shadow=popink,left=11pt,right=11pt,top=10pt,bottom=10pt,height=3.1cm,valign=center]
      \tikz[baseline=-0.7ex]\node[circle,fill=white,draw=popink,line width=1.3pt,minimum size=1.6cm,inner sep=0]{\popdisplay\fontsize{24}{24}\selectfont\color{poppurple}+};%
      \hspace{8pt}\begin{minipage}[c]{0.6\linewidth}
        {\popdisplay\fontsize{11}{12.5}\selectfont\color{white}\MakeUppercase{Passe à OnBuch+}}\par\vspace{4pt}
        {\raggedright\popsemi\fontsize{8.4}{11}\selectfont\color{white}Tous les cours signés, Tuteur IA illimité, fiches PDF — onglet OnBuch+ de l'app\par}
      \end{minipage}
    \end{tcolorbox}
  \end{minipage}
  \vspace{8pt}
  \popcontenu
  \vfill
  \signatureonbuch
  \restoregeometry
  \newpage
}

% Rangée « Ce cours contient » (page de garde)
\newcommand{\poptile}[3]{% #1 couleur #2 gros texte #3 légende
  \begin{tcolorbox}[enhanced,colback=#1,colframe=popink,boxrule=2pt,arc=9pt,shadow={2.5pt}{-2.5pt}{0pt}{popink},
      left=6pt,right=6pt,top=9pt,bottom=9pt,halign=center,valign=center,height=1.75cm,width=\linewidth,nobeforeafter]
    {\popdisplay\fontsize{12.5}{13}\selectfont\color{white}\MakeUppercase{#2}}\par\vspace{3pt}
    {\centering\popsemi\fontsize{7.8}{9.5}\selectfont\color{white}#3\par}
  \end{tcolorbox}}
\newcommand{\popcontenu}{%
  \par\noindent{\popxboldls\fontsize{7.5}{7.5}\selectfont\color{popmuted}CE COURS SIGNÉ CONTIENT}\par\vspace{7pt}
  \noindent\begin{minipage}[t]{0.235\linewidth}\poptile{popblue}{Cours}{complet, illustré, pas à pas}\end{minipage}\hfill
  \begin{minipage}[t]{0.235\linewidth}\poptile{poporange}{Méthodes}{et exercices résolus}\end{minipage}\hfill
  \begin{minipage}[t]{0.235\linewidth}\poptile{poppink}{Exercices}{type examen + corrigés}\end{minipage}\hfill
  \begin{minipage}[t]{0.235\linewidth}\poptile{popgreen}{Fiche bilan}{l'essentiel à retenir}\end{minipage}\par}

% Sommaire
\newtcolorbox{sommaire}{enhanced,colback=white,colframe=popink,boxrule=2.2pt,arc=10pt,
  drop shadow=popink,left=18pt,right=18pt,top=13pt,bottom=13pt,before skip=6pt,after skip=20pt,
  before upper={\poppill{Sommaire}{poppurple}{white}\par\vspace{9pt}\popsemi\fontsize{10.6}{14.5}\selectfont\color{popink}}}

% Signature de fin de cours — poussée tout en bas de la dernière page
\newcommand{\findecours}{%
  \par\vfill\nopagebreak
  \begin{center}\popsquiggle{poporange}\end{center}\vspace{4pt}
  \signatureonbuch
}
\AtBeginDocument{\def\llbracket{\mathopen{[\mkern-3.5mu[}}\def\rrbracket{\mathclose{]\mkern-3.5mu]}}} % intervalles d'entiers (glyphe ⟦ absent de la police)
% Glyphes absents de Fira Math : redéfinis à partir de glyphes disponibles
\AtBeginDocument{%
  \def\setminus{\mathbin{\backslash}}\let\smallsetminus\setminus
  \def\vdots{\mathord{\vbox{\baselineskip3.2pt\lineskiplimit0pt\kern4pt\hbox{.}\hbox{.}\hbox{.}}}}%
  \def\ddots{\mathinner{\mkern1mu\raise6.5pt\hbox{.}\mkern2mu\raise3.6pt\hbox{.}\mkern2mu\raise0.7pt\hbox{.}\mkern1mu}}%
  \def\triangle{\mathord{\scalebox{0.9}{$\symup{\Delta}$}}}%
  \def\nabla{\mathord{\raisebox{\depth}{\scalebox{1}[-1]{$\symup{\Delta}$}}}}%
  \def\checkmark{\popcheck{popdark}}%
  \def\star{\mathbin{\ast}}\def\bigstar{\mathord{\ast}}%
  \def\diamond{\mathbin{\scalebox{0.6}{$\Diamond$}}}%
  \def\bigcup{\mathop{\mathchoice{\raisebox{-0.3em}{\scalebox{1.7}{$\cup$}}}{\scalebox{1.25}{$\cup$}}{\cup}{\cup}}\displaylimits}%
  \def\bigcap{\mathop{\mathchoice{\raisebox{-0.3em}{\scalebox{1.7}{$\cap$}}}{\scalebox{1.25}{$\cap$}}{\cap}{\cap}}\displaylimits}%
  \def\lesseqqgtr{\mathrel{\lessgtr}}\def\bigoplus{\mathop{\oplus}\displaylimits}%
}
\providecommand{\ssi}{si et seulement si} % abréviation fréquente
\AtBeginDocument{\providecommand{\wideparen}{\overparen}} % arc de cercle

%% FIGFIX-BEGIN v5 — correctifs globaux des figures (docs/onbuch-generation/tools/figfix)
\usepackage{adjustbox}\usepackage{etoolbox}\tcbuselibrary{raster}
% toute figure TikZ/pgfplots plus large que la ligne est réduite à la largeur disponible
\BeforeBeginEnvironment{tikzpicture}{\begin{adjustbox}{max width=\linewidth}}
\AfterEndEnvironment{tikzpicture}{\end{adjustbox}}
% légendes pgfplots lisibles par-dessus les courbes
\pgfplotsset{every axis/.append style={legend style={fill=white,fill opacity=0.92,text opacity=1,draw=black!15,inner sep=3pt}}}
% étiquettes d'arêtes (syntaxe edge["..."]) avec fond blanc
\tikzset{every edge quotes/.append style={fill=white,inner sep=1.2pt}}
%% FIGFIX-END
\begin{document}

\pagegarde{Leçon 35 — Vocabulaire et compréhension de texte : internet}{Chinois}{1ère A}{Module 5 : Mass médias et télécommunication}{ZHA35}{2 h}{Cette leçon propose la lecture et la compréhension d'un court texte chinois sur internet et ses usages quotidiens. Elle permet d'acquérir le vocabulaire essentiel du numérique (caractères, pinyin, tons), d'étudier les radicaux et l'ordre des traits de caractères clés, et de s'entraîner à répondre à des questions de compréhension.
\vspace{4pt}
\begin{multicols}{2}\small
\begin{itemize}
\item À la fin de cette leçon, je sais lire à voix haute un court texte chinois sur internet en respectant le pinyin et les tons.
\item Je sais reconnaître, prononcer et traduire au moins 15 mots de vocabulaire liés à internet.
\item Je sais employer correctement ce vocabulaire dans des phrases simples.
\item Je sais identifier les radicaux (clés) et tracer l'ordre des traits des caractères 网, 络, 信, 电 et 脑.
\item Je sais répondre par écrit à des questions de compréhension portant sur le texte.
\item Je sais utiliser les collocations fréquentes du thème (上网, 发邮件, 查资料...).
\end{itemize}\end{multicols}}

\begin{sommaire}
\begin{multicols}{2}
\begin{enumerate}
\item Mise en route et découverte du thème
\item Texte support : lecture et écoute
\item Glossaire du thème
\item Clés (radicaux) et ordre des traits
\item Compréhension du texte
\item Collocations et production guidée
\item Activité d'intégration
\item Méthodes et exercices résolus
\item Exercices
\item Corrigés des exercices
\item Fiche bilan
\end{enumerate}
\end{multicols}
\end{sommaire}

\begin{objectifs}
\begin{itemize}
\item À la fin de cette leçon, je sais lire à voix haute un court texte chinois sur internet en respectant le pinyin et les tons.
\item Je sais reconnaître, prononcer et traduire au moins 15 mots de vocabulaire liés à internet.
\item Je sais employer correctement ce vocabulaire dans des phrases simples.
\item Je sais identifier les radicaux (clés) et tracer l'ordre des traits des caractères 网, 络, 信, 电 et 脑.
\item Je sais répondre par écrit à des questions de compréhension portant sur le texte.
\item Je sais utiliser les collocations fréquentes du thème (上网, 发邮件, 查资料...).
\item Je sais rédiger un court paragraphe guidé sur mes usages d'internet.
\end{itemize}
\end{objectifs}

\begin{prerequis}
\begin{itemize}
\item Vocabulaire des médias et de la télécommunication vu dans les leçons précédentes du Module 5 (téléphone, télévision, radio).
\item Maîtrise du pinyin et des quatre tons (Seconde).
\item Structures de base : phrase sujet-verbe-complément, 是, 有, 可以.
\item Notions sur les radicaux et l'ordre des traits (Seconde).
\item Nombres et expressions de temps pour parler de fréquence d'utilisation.
\end{itemize}
\end{prerequis}

\newpage

\coursec{Mise en route et découverte du thème}

\courssub{Situation d'accroche et problématique}

À Yaoundé comme à Douala, dès la sortie des classes, les élèves sortent leur téléphone : WhatsApp pour discuter avec les amis, Facebook pour suivre l'actualité, TikTok pour regarder des vidéos. Tout cela est possible grâce au forfait internet MTN ou Orange que les parents achètent chaque semaine. Au quartier, le cybercafé ne désemplit pas : on y imprime des documents, on y remplit des formulaires en ligne, on y passe des appels vidéo.

De l'autre côté du monde, en Chine, plus d'un milliard de personnes sont connectées chaque jour. Mais surprise : elles n'utilisent ni WhatsApp ni Google ! Elles discutent, paient leurs courses et commandent des taxis avec une seule application, \zh{微信} (Wēixìn, WeChat), et cherchent des informations sur \zh{百度} (Bǎidù, Baidu).

\textbf{Problématique :} comment dit-on « surfer sur internet », « envoyer un e-mail » ou « mot de passe » en chinois ? Comment un Chinois parle-t-il de sa vie numérique quotidienne ? Cette leçon vous ouvre les portes du monde numérique chinois : à la fin de la séance, vous saurez nommer les outils, décrire vos habitudes en ligne et comprendre un court texte sur internet en Chine et au Cameroun.

\begin{savaistu}[Un milliard d'internautes]
La Chine compte plus d'un milliard d'internautes : c'est le plus grand nombre d'utilisateurs d'internet au monde. L'application la plus utilisée est \zh{微信} (Wēixìn, WeChat) : elle sert à discuter avec ses amis, à payer au marché, à commander un repas ou un taxi, et même à prendre rendez-vous chez le médecin. Au Cameroun, on paie souvent avec Mobile Money ; en Chine, on paie presque tout avec le téléphone grâce à WeChat ou Alipay. Un Chinois peut sortir de chez lui sans argent liquide, mais jamais sans son téléphone !
\end{savaistu}

\courssub{Brainstorming : quels mots associez-vous à « internet » ?}

Avant d'ouvrir le dictionnaire, activons nos connaissances. Travail en deux temps : d'abord en français, puis tentative de traduction en chinois.

\textbf{Étape 1 — En français.} Fermez les yeux et pensez à « internet ». Quels mots vous viennent à l'esprit ? Voici les réponses typiques d'une classe de Première à Yaoundé :

\begin{itemize}
\item le téléphone, l'ordinateur, l'écran ;
\item le mot de passe, le forfait, la connexion, le wifi ;
\item envoyer un message, un e-mail, une photo ;
\item surfer, chercher, télécharger, discuter en ligne.
\end{itemize}

\textbf{Étape 2 — Tentative en chinois.} Certains mots sont déjà connus des leçons précédentes : \zh{手机} (shǒujī, téléphone portable), \zh{电脑} (diànnǎo, ordinateur). D'autres sont nouveaux : c'est justement le cœur de cette leçon. Observons d'abord la carte mentale du thème.

\begin{popfigure}
\begin{center}\tcbox[colback=popblue,colframe=popblue,coltext=white,arc=9pt,boxsep=4pt,left=6pt,right=6pt]{\bfseries\small 互联网 hùliánwǎng internet}\end{center}
\vspace{-2pt}
\begin{tcbraster}[raster columns=3,raster equal height=rows,raster column skip=6pt,raster row skip=6pt,enhanced,blankest]
\begin{tcolorbox}[enhanced,colback=popblue,colframe=popblue,coltext=white,boxrule=0.9pt,arc=3pt,left=4pt,right=4pt,top=3pt,bottom=3pt,boxsep=1pt,halign=left]\bfseries\scriptsize 上网 shàngwǎng surfer en ligne\end{tcolorbox}
\begin{tcolorbox}[enhanced,colback=poporange,colframe=poporange,coltext=white,boxrule=0.9pt,arc=3pt,left=4pt,right=4pt,top=3pt,bottom=3pt,boxsep=1pt,halign=left]\bfseries\scriptsize 电脑 diànnǎo ordinateur\end{tcolorbox}
\begin{tcolorbox}[enhanced,colback=popgreen,colframe=popgreen,coltext=white,boxrule=0.9pt,arc=3pt,left=4pt,right=4pt,top=3pt,bottom=3pt,boxsep=1pt,halign=left]\bfseries\scriptsize 手机 shǒujī téléphone\end{tcolorbox}
\begin{tcolorbox}[enhanced,colback=poppurple,colframe=poppurple,coltext=white,boxrule=0.9pt,arc=3pt,left=4pt,right=4pt,top=3pt,bottom=3pt,boxsep=1pt,halign=left]\bfseries\scriptsize 邮件 yóujiàn e-mail\end{tcolorbox}
\begin{tcolorbox}[enhanced,colback=popgold,colframe=popgold,coltext=white,boxrule=0.9pt,arc=3pt,left=4pt,right=4pt,top=3pt,bottom=3pt,boxsep=1pt,halign=left]\bfseries\scriptsize 信息 xìnxī information\end{tcolorbox}
\end{tcbraster}

\legende{Carte mentale du vocabulaire d'internet : le mot central 互联网 et ses cinq premières branches.}
\end{popfigure}

Observez bien : quatre de ces mots contiennent des caractères que vous connaissez déjà (\zh{电} « électricité » dans \zh{电脑} et \zh{电话}, \zh{手} « main » dans \zh{手机}). Le chinois construit son vocabulaire moderne avec des briques anciennes : c'est une excellente nouvelle pour votre mémoire !

\courssub{Le mot central : 互联网 (hùliánwǎng) et sa forme courte 网 (wǎng)}

\begin{definition}[互联网 et 网]
\zh{互联网} (hùliánwǎng) est le nom complet d'« internet » : \zh{互} (hù) signifie « mutuellement », \zh{联} (lián) « relier, connecter », \zh{网} (wǎng) « filet, réseau ». Internet est donc littéralement « le réseau qui relie mutuellement ». Dans la langue courante, on utilise presque toujours la forme courte \zh{网} (wǎng) : « aller sur internet » se dit \zh{上网} (shàngwǎng), littéralement « monter sur le filet ».
\end{definition}

\begin{attention}[网, un « filet » qui relie le monde]
\zh{网} (wǎng) signifie « filet » à l'origine : le filet du pêcheur, le filet du footballeur. Internet est imaginé comme un immense filet dont les nœuds relient le monde entier. Ne traduisez donc jamais \zh{网} par « toile » : « la toile » en français se dit aussi \zh{网络} (wǎngluò, le réseau). Autre piège : \zh{上网} se prononce shàngwǎng avec deux quatrièmes tons bien marqués, pas « shang-wang » à la française !
\end{attention}

Voici les premiers mots à découvrir aujourd'hui, organisés autour de ce mot central :

\begin{center}
\begin{tabularx}{\linewidth}{|X|X|X|X|}
\hline
\rowcolor{popblueL} Caractères & Pinyin & Nature & Traduction \\
\hline
互联网 & hùliánwǎng & nom & internet \\
\hline
网 & wǎng & nom & filet, réseau, internet (forme courte) \\
\hline
上网 & shàngwǎng & verbe & aller sur internet, surfer \\
\hline
网络 & wǎngluò & nom & réseau, la toile \\
\hline
网友 & wǎngyǒu & nom & ami rencontré sur internet, internaute ami \\
\hline
网站 & wǎngzhàn & nom & site web \\
\hline
\end{tabularx}
\end{center}

Remarquez la logique de construction : \zh{网} + \zh{友} (yǒu, ami) donne \zh{网友} « ami du net » ; \zh{网} + \zh{站} (zhàn, station) donne \zh{网站} « station du net », c'est-à-dire un site web. Comparer avec le français : nous disons « internaute », « site », « surfer » ; le chinois dit « ami du filet », « station du filet », « monter sur le filet ». Les images sont différentes, mais l'idée est la même.

\begin{methode}[Mémoriser un mot nouveau en cinq points]
Pour retenir durablement un mot chinois, fichez-le toujours ainsi :
\begin{enumerate}
\item \textbf{Caractère} : écrivez-le trois fois en respectant l'ordre des traits ;
\item \textbf{Pinyin} : notez-le avec les lettres, sans les tons d'abord ;
\item \textbf{Ton} : ajoutez la marque du ton sur la bonne voyelle et prononcez à voix haute ;
\item \textbf{Traduction} : une traduction précise, pas approximative ;
\item \textbf{Phrase d'exemple} : une phrase courte et personnelle.
\end{enumerate}
Exemple de fiche réussie : \zh{上网} — shàngwǎng — 4\textsuperscript{e} + 3\textsuperscript{e} ton — « aller sur internet » — \zh{我每天晚上上网。} (Wǒ měi tiān wǎnshang shàngwǎng. « Je vais sur internet tous les soirs. »)
\end{methode}

\courssub{Observation d'images : deux mondes connectés}

L'enseignant projette deux images. Décrivez-les d'abord en français, puis essayez deux ou trois mots chinois.

\textbf{Image 1 — Un cybercafé à Yaoundé.} Dans une petite salle du quartier Essos, dix ordinateurs sont allumés. Un élève en uniforme imprime ses notes, une jeune femme remplit un formulaire d'inscription à l'université, un commerçant envoie un e-mail à un fournisseur chinois. Sur le mur, une affiche : « Wifi gratuit ». Mots à repérer : \zh{电脑} (diànnǎo, ordinateur), \zh{上网} (shàngwǎng, surfer), \zh{发邮件} (fā yóujiàn, envoyer un e-mail).

\textbf{Image 2 — Un jeune Chinois sur WeChat.} Dans le métro de Shanghai, Li Ming tient son téléphone. Il discute avec ses amis sur \zh{微信} (Wēixìn), paie son ticket avec son téléphone et regarde une courte vidéo. Mots à repérer : \zh{手机} (shǒujī, téléphone), \zh{信息} (xìnxī, message, information), \zh{密码} (mìmǎ, mot de passe).

\begin{exemplebox}[Premières phrases du thème]
\zh{我每天都上网。} — \emph{Wǒ měi tiān dōu shàngwǎng.} — « Je vais sur internet tous les jours. »\\[2pt]
\zh{他用手机发邮件。} — \emph{Tā yòng shǒujī fā yóujiàn.} — « Il envoie un e-mail avec son téléphone. »\\[2pt]
\zh{你的密码是什么？} — \emph{Nǐ de mìmǎ shì shénme ?} — « Quel est ton mot de passe ? »\\[2pt]
\zh{我在网上看法语信息。} — \emph{Wǒ zài wǎngshang kàn Fǎyǔ xìnxī.} — « Je lis des informations en français sur internet. »
\end{exemplebox}

Notez la structure de la dernière phrase : \zh{在} + lieu + verbe, comme dans \zh{我在学校学习} « j'étudie à l'école ». Ici, le « lieu » est virtuel : \zh{网上} (wǎngshang, « sur le net »). Le chinois traite internet comme un endroit où l'on se trouve — exactement comme le français « sur internet ».

\courssub{Annonce des objectifs de la leçon}

À l'issue de cette leçon de deux heures, vous serez capable de :

\begin{enumerate}
\item reconnaître, prononcer (pinyin et tons) et écrire une quinzaine de mots du thème « internet » ;
\item lire et comprendre un court texte chinois sur internet en Chine et au Cameroun ;
\item répondre à des questions de compréhension en phrases complètes ;
\item employer les collocations essentielles : \zh{上网}, \zh{发邮件}, \zh{发信息}, \zh{用电脑} ;
\item dire en chinois trois de vos habitudes numériques quotidiennes.
\end{enumerate}

\coursec{Texte support : lecture et écoute}

\courssub{Première écoute du texte support (sans support écrit)}

Fermez vos cahiers. L'enseignant lit le texte deux fois, à vitesse normale. Votre tâche à la première écoute : identifier le thème général et les mots que vous reconnaissez déjà. À la deuxième écoute : notez en français les deux ou trois idées principales. Voici le texte que vous entendez (il sera étudié en détail dans la section suivante) :

\begin{textebox}[互联网和我们的生活 — Hùliánwǎng hé wǒmen de shēnghuó]
现在，互联网对我们的生活非常重要。在中国，有十亿多人上网。他们用手机和电脑看新闻、发信息、买东西。很多中国人用微信：他们用微信聊天、付钱，也用微信叫出租车。\\
\emph{Xiànzài, hùliánwǎng duì wǒmen de shēnghuó fēicháng zhòngyào. Zài Zhōngguó, yǒu shí yì duō rén shàngwǎng. Tāmen yòng shǒujī hé diànnǎo kàn xīnwén, fā xìnxī, mǎi dōngxi. Hěn duō Zhōngguórén yòng Wēixìn : tāmen yòng Wēixìn liáotiān, fùqián, yě yòng Wēixìn jiào chūzūchē.}\\
« Aujourd'hui, internet est très important pour notre vie. En Chine, plus d'un milliard de personnes vont sur internet. Avec leur téléphone et leur ordinateur, ils lisent les nouvelles, envoient des messages et achètent des choses. Beaucoup de Chinois utilisent WeChat : ils s'en servent pour discuter, payer, et aussi pour appeler un taxi. »\\[4pt]
在喀麦隆，上网的人也越来越多。雅温得和杜阿拉的年轻人喜欢用脸书和WhatsApp。他们上网看法语和英语的信息，也看足球比赛。很多学生在网上学习：他们看老师的课，做练习，还给朋友发邮件。\\
\emph{Zài Kāmàilóng, shàngwǎng de rén yě yuè lái yuè duō. Yǎwēndé hé Dù'ālā de niánqīngrén xǐhuan yòng liǎnshū hé WhatsApp. Tāmen shàngwǎng kàn Fǎyǔ hé Yīngyǔ de xìnxī, yě kàn zúqiú bǐsài. Hěn duō xuésheng zài wǎngshang xuéxí : tāmen kàn lǎoshī de kè, zuò liànxí, hái gěi péngyou fā yóujiàn.}\\
« Au Cameroun, les gens qui vont sur internet sont aussi de plus en plus nombreux. Les jeunes de Yaoundé et de Douala aiment utiliser Facebook et WhatsApp. Ils vont en ligne pour lire des informations en français et en anglais, et aussi pour regarder les matchs de football. Beaucoup d'élèves étudient sur internet : ils regardent les cours du professeur, font des exercices et envoient aussi des e-mails à leurs amis. »\\[4pt]
互联网让世界变小了。喀麦隆的学生可以在网上认识中国朋友，中国学生也可以了解喀麦隆。但是，上网的时候，我们要小心：不要把密码告诉别人，也不要看太长时间的视频。互联网是好老师，也是好朋友，可是我们要好好地用它。\\
\emph{Hùliánwǎng ràng shìjiè biàn xiǎo le. Kāmàilóng de xuésheng kěyǐ zài wǎngshang rènshi Zhōngguó péngyou, Zhōngguó xuésheng yě kěyǐ liǎojiě Kāmàilóng. Dànshì, shàngwǎng de shíhou, wǒmen yào xiǎoxīn : bú yào bǎ mìmǎ gàosu biérén, yě bú yào kàn tài cháng shíjiān de shìpín. Hùliánwǎng shì hǎo lǎoshī, yě shì hǎo péngyou, kěshì wǒmen yào hǎohāo de yòng tā.}\\
« Internet a rendu le monde plus petit. Les élèves camerounais peuvent se faire des amis chinois en ligne, et les élèves chinois peuvent aussi découvrir le Cameroun. Mais quand nous surfons, nous devons être prudents : il ne faut pas donner son mot de passe aux autres, ni regarder des vidéos trop longtemps. Internet est un bon professeur et un bon ami, mais il faut bien l'utiliser. »
\end{textebox}

\textbf{Petit glossaire d'écoute} (pour lever les blocages immédiats) :

\begin{center}
\begin{tabularx}{\linewidth}{|X|X|X|X|}
\hline
\rowcolor{popblueL} Caractères & Pinyin & Nature & Traduction \\
\hline
重要 & zhòngyào & adjectif & important \\
\hline
新闻 & xīnwén & nom & nouvelles, actualités \\
\hline
买东西 & mǎi dōngxi & expression verbale & acheter des choses, faire des achats \\
\hline
聊天 & liáotiān & verbe & bavarder, discuter \\
\hline
付钱 & fùqián & verbe & payer \\
\hline
越来越 & yuè lái yuè & structure & de plus en plus \\
\hline
比赛 & bǐsài & nom & match, compétition \\
\hline
认识 & rènshi & verbe & connaître, faire connaissance \\
\hline
了解 & liǎojiě & verbe & comprendre, découvrir \\
\hline
小心 & xiǎoxīn & adjectif/verbe & prudent, faire attention \\
\hline
密码 & mìmǎ & nom & mot de passe \\
\hline
告诉 & gàosu & verbe & dire, raconter à quelqu'un \\
\hline
视频 & shìpín & nom & vidéo \\
\hline
\end{tabularx}
\end{center}

\begin{exoresolu}[Première écoute — vérification globale]
\textbf{Énoncé.} Après la première écoute (sans texte sous les yeux), répondez en français : (1) De quoi parle ce texte ? (2) Quels deux pays sont mentionnés ? (3) Citez deux activités que les gens font sur internet selon le texte.
\tcblower
\textbf{Solution détaillée.}
\begin{enumerate}
\item Le texte parle d'internet (\zh{互联网}) et de son importance dans la vie quotidienne (\zh{对我们的生活非常重要}).
\item Les deux pays sont la Chine (\zh{中国}) et le Cameroun (\zh{喀麦隆}).
\item Activités possibles : lire les nouvelles (\zh{看新闻}), envoyer des messages (\zh{发信息}), acheter des choses (\zh{买东西}), payer avec WeChat (\zh{付钱}), regarder des matchs de football (\zh{看足球比赛}), étudier en ligne (\zh{在网上学习}).
\end{enumerate}
\textbf{Astuce de méthode :} à la première écoute, ne cherchez pas à tout comprendre. Repérez les mots déjà connus (\zh{手机}, \zh{电脑}, \zh{上网}) et les noms propres (\zh{中国}, \zh{喀麦隆}, \zh{雅温得}) : ils donnent le squelette du texte.
\end{exoresolu}

\begin{attention}[Pièges de francophones à l'écoute]
Trois difficultés classiques à surveiller dès maintenant :
\begin{itemize}
\item \textbf{Les tons proches} : \zh{买} (mǎi, acheter, 3\textsuperscript{e} ton) et \zh{卖} (mài, vendre, 4\textsuperscript{e} ton) ne diffèrent que par le ton. Dans le texte, c'est bien \zh{买东西} « acheter ».
\item \textbf{Les mots à deux syllabes} : ne découpez pas \zh{信息} en « xìn + xī » ; écoutez le mot entier xìnxī comme un seul bloc, comme « informa-tion » en français.
\item \textbf{Les noms propres transcrits} : \zh{喀麦隆} (Kāmàilóng) et \zh{雅温得} (Yǎwēndé) sont des transcriptions phonétiques de « Cameroun » et « Yaoundé ». Ne cherchez pas à traduire chaque caractère : \zh{隆} ne veut pas dire « dragon » ici !
\end{itemize}
\end{attention}

\courssub{Lecture guidée et analyse linguistique}

Ouvrez le texte. Lisez-le à voix basse en suivant du doigt. Repérez les structures grammaticales déjà vues : \zh{有} (il y a), \zh{用} (utiliser), \zh{在} + lieu + verbe, \zh{也} (aussi), \zh{还} (en plus), \zh{把} (structure de disposition), \zh{了} (changement d'état / aspect accompli).

\begin{propriete}[Structure \zh{在} + lieu + Verbe : action située]
Pour dire où se déroule une action, le chinois place le complément de lieu \textbf{avant} le verbe, introduit par \zh{在} (zài).
\zh{Sujet + 在 + Lieu + Verbe (+ Complément)}.
Exemples du texte : \zh{在网上学习} (étudier sur internet), \zh{在喀麦隆} (au Cameroun — ici \zh{在} introduit le cadre spatial de la phrase).
\end{propriete}

\begin{propriete}[Structure \zh{把} (bǎ) : action sur objet défini]
\zh{把} (bǎ) place l'objet \textbf{avant} le verbe pour insister sur ce qu'il lui arrive. L'objet doit être connu/défini.
\zh{Sujet + 把 + Objet + Verbe + Résultat/Complément}.
Exemple du texte : \zh{不要把密码告诉别人} (Ne dis pas le mot de passe à d'autres / Ne donne pas le mot de passe à d'autres). \zh{密码} est l'objet manipulé.
\end{propriete}

\begin{propriete}[Structure \zh{越来越} + Adjectif/Verbe : « de plus en plus »]
\zh{越来越} (yuè lái yuè) précède un adjectif ou un verbe d'état pour exprimer une progression.
Exemple du texte : \zh{上网的人也越来越多} (Les gens qui vont sur internet sont de plus en plus nombreux).
\end{propriete}

\begin{exemplebox}[Analyse de phrases clés du texte]
\zh{他们用手机和电脑看新闻、发信息、买东西。}\\
\emph{Tāmen yòng shǒujī hé diànnǎo kàn xīnwén, fā xìnxī, mǎi dōngxi.}\\
Structure : Sujet + \zh{用} Outil + Verbe 1 + Objet 1, Verbe 2 + Objet 2, Verbe 3 + Objet 3. (Série d'actions avec le même sujet et outil).\\[4pt]
\zh{互联网让世界变小了。}\\
\emph{Hùliánwǎng ràng shìjiè biàn xiǎo le.}\\
Structure : Sujet + \zh{让} (ràng, faire/laisser) + Objet + Verbe/Adjectif + \zh{了} (changement d'état). \zh{让} exprime la causalité.\\[4pt]
\zh{我们要好好地用它。}\\
\emph{Wǒmen yào hǎohāo de yòng tā.}\\
\zh{好好} (hǎohāo, bien/soigneusement) + \zh{地} (de, marqueur d'adverbe) + Verbe. Règle : adjectif bisyllabique + \zh{地} + verbe.
\end{exemplebox}

\coursec{Glossaire du thème : vocabulaire principal}

\courssub{Tableau récapitulatif (leçons + texte support)}

Ce tableau regroupe le vocabulaire actif à maîtriser pour l'évaluation. Les mots marqués d'un astérisque (*) sont du niveau HSK 1-2 (révision), les autres sont HSK 3-4 (nouveaux).

\begin{center}
\begin{tabularx}{\linewidth}{|X|X|X|X|}
\hline
\rowcolor{popblueL} Caractères & Pinyin & Nature & Traduction \\
\hline
互联网 & hùliánwǎng & nom & internet \\
\hline
网 & wǎng & nom & filet, réseau, internet (forme courte) \\
\hline
上网 & shàngwǎng & verbe & aller sur internet, surfer \\
\hline
网络 & wǎngluò & nom & réseau, la toile \\
\hline
网站 & wǎngzhàn & nom & site web \\
\hline
网友 & wǎngyǒu & nom & ami du net, internaute \\
\hline
手机* & shǒujī & nom & téléphone portable \\
\hline
电脑* & diànnǎo & nom & ordinateur \\
\hline
邮件* & yóujiàn & nom & e-mail, courriel \\
\hline
发邮件* & fā yóujiàn & loc. verbale & envoyer un e-mail \\
\hline
信息 & xìnxī & nom & information, message, données \\
\hline
发信息 & fā xìnxī & loc. verbale & envoyer un message (SMS, WeChat) \\
\hline
密码 & mìmǎ & nom & mot de passe, code secret \\
\hline
新闻 & xīnwén & nom & nouvelles, actualités \\
\hline
看新闻 & kàn xīnwén & loc. verbale & lire/regarder les nouvelles \\
\hline
买东西 & mǎi dōngxi & loc. verbale & faire des achats, acheter des choses \\
\hline
聊天 & liáotiān & verbe & bavarder, discuter (en ligne) \\
\hline
付钱 & fùqián & verbe & payer \\
\hline
视频 & shìpín & nom & vidéo \\
\hline
微信 & Wēixìn & nom propre & WeChat \\
\hline
脸书 & liǎnshū & nom propre & Facebook \\
\hline
雅温得 & Yǎwēndé & nom propre & Yaoundé \\
\hline
杜阿拉 & Dù'ālā & nom propre & Douala \\
\hline
喀麦隆* & Kāmàilóng & nom propre & Cameroun \\
\hline
中国* & Zhōngguó & nom propre & Chine \\
\hline
越来越 & yuè lái yuè & adv. & de plus en plus \\
\hline
认识 & rènshi & verbe & connaître, faire connaissance \\
\hline
了解 & liǎojiě & verbe & comprendre, se renseigner sur \\
\hline
小心 & xiǎoxīn & adj./verbe & prudent / faire attention \\
\hline
告诉 & gàosu & verbe & dire à, informer \\
\hline
别人 & biérén & pronom & les autres, d'autres personnes \\
\hline
\end{tabularx}
\end{center}

\courssub{Collocations indispensables (Verb + Objet)}

En chinois, les verbes d'action sur internet s'utilisent presque toujours avec leur objet. Mémorisez ces blocs inséparables :

\begin{center}
\begin{tabularx}{\linewidth}{|X|X|X|}
\hline
\rowcolor{popblueL} Collocation (Caractères + Pinyin) & Structure & Traduction \\
\hline
\zh{上网} (shàngwǎng) & Verbe (VO) & Surfer / Aller sur internet \\
\hline
\zh{发邮件} (fā yóujiàn) & Verbe + Objet & Envoyer un e-mail \\
\hline
\zh{发信息} (fā xìnxī) & Verbe + Objet & Envoyer un message \\
\hline
\zh{看新闻} (kàn xīnwén) & Verbe + Objet & Lire les nouvelles \\
\hline
\zh{看视频} (kàn shìpín) & Verbe + Objet & Regarder des vidéos \\
\hline
\zh{买东西} (mǎi dōngxi) & Verbe + Objet & Faire des achats \\
\hline
\zh{聊天} (liáotiān) & Verbe (VO) & Discuter / Bavarder \\
\hline
\zh{付钱} (fùqián) & Verbe + Objet & Payer \\
\hline
\zh{输密码} (shūrù mìmǎ) & Verbe + Objet & Entrer un mot de passe \\
\hline
\zh{登录} (dēnglù) & Verbe & Se connecter / Se logger \\
\hline
\end{tabularx}
\end{center}

\begin{attention}[Erreurs fréquentes : \zh{发} vs \zh{送}, \zh{看} vs \zh{读}]
\begin{itemize}
\item \zh{发} (fā) = envoyer, émettre (focus sur l'action d'expédition). \zh{送} (sòng) = livrer, offrir (focus sur la réception). On dit \zh{发邮件}, \zh{发信息}, \zh{发微信}.
\item \zh{看} (kàn) = regarder, lire (visuel, général). \zh{读} (dú) = lire à haute voix, étudier un texte. Pour les nouvelles ou vidéos : \zh{看新闻}, \zh{看视频}. Pour un livre ou un article long : \zh{读书} (dúshū), \zh{读文章} (dú wénzhāng).
\item \zh{上网} est un verbe intransitif (pas d'objet). On ne dit pas \zh{上网互联网}. On dit \zh{上网} ou \zh{在网上} + verbe.
\end{itemize}
\end{attention}

\coursec{Clés (radicaux) et ordre des traits}

\courssub{Analyse morphologique : les briques du vocabulaire numérique}

Comprendre les radicaux (clés) permet de deviner le sens et de mémoriser l'écriture. Voici l'analyse de 8 caractères clés du thème.

\begin{center}
\begin{tabularx}{\linewidth}{|X|X|X|X|X|}
\hline
\rowcolor{popblueL} Caractère & Radical (Clé) & Sens du radical & Nombre de traits & Ordre des traits (description) \\
\hline
\zh{网} & \zh{纟} (sī, soie) & Filet, tissu, toile & 6 & 1. \zh{纟} (3 traits : point, pli, horizontal) 2. \zh{亡} (wáng, 3 traits : horizontaux, verticale) \\
\hline
\zh{联} & \zh{耳} (ěr, oreille) & Relier (oreilles qui entendent ensemble) & 12 & 1. \zh{耳} (6 traits) 2. \zh{关} (guān, 6 traits : haut en bas, cadre puis dedans) \\
\hline
\zh{信} & \zh{亻} (rén, homme) & Parole de l'homme, confiance & 9 & 1. \zh{亻} (2 traits) 2. \zh{言} (yán, 7 traits : point, horizontaux, verticale, point final) \\
\hline
\zh{息} & \zh{心} (xīn, cœur) & Souffle, information, repos & 10 & 1. \zh{自} (zì, 6 traits : cadre, traits intérieurs) 2. \zh{心} (3 traits : points, crochet) \\
\hline
\zh{密} & \zh{宀} (mián, toit) & Secret, caché sous le toit & 11 & 1. \zh{宀} (3 traits) 2. \zh{必} (bì, 7 traits : haut en bas, cœur au milieu) \\
\hline
\zh{码} & \zh{石} (shí, pierre) & Code, chiffre (pierre pour compter) & 10 & 1. \zh{石} (5 traits) 2. \zh{马} (mǎ, 3 traits : horizontaux, verticale, points) — \textbf{attention :} \zh{马} simplifié ici ! \\
\hline
\zh{站} & \zh{立} (lì, debout) & Station, se tenir debout & 10 & 1. \zh{立} (5 traits) 2. \zh{占} (zhàn, 5 traits : horizontal, verticale, traits croisés) \\
\hline
\zh{视} & \zh{见} (jiàn, voir) & Regarder, vision & 11 & 1. \zh{示} (shì, 5 traits : autel) 2. \zh{见} (4 traits : œil sur jambes) — \textbf{ordre :} \zh{示} puis \zh{见} \\
\hline
\end{tabularx}
\end{center}

\begin{methode}[Écrire \zh{密} (mìmǎ) et \zh{码} (mǎ) correctement]
\begin{enumerate}
\item \textbf{\zh{密} (mì)} : Commencez par le toit \zh{宀} (point, pli, horizontal). Ensuite \zh{必} : haut \zh{丷} (huit), puis \zh{心} (cœur) au centre, enfin le trait de base horizontal long qui traverse.
\item \textbf{\zh{码} (mǎ)} : Écrivez \zh{石} (pierre : horizontale, verticale, horizontale, crochet, point). Puis \zh{马} (cheval simplifié : horizontale, verticale, horizontale, trois points groupés). \textbf{Piège :} \zh{马} a 3 traits en simplifié (contre 10 en traditionnel) !
\end{enumerate}
\end{methode}

\begin{experience}[Atelier écriture : « Mon carnet de radicaux »]
Sur votre cahier d'exercices, créez une page « Radicaux du numérique ». Pour chaque radical ci-dessus (\zh{纟}, \zh{耳}, \zh{亻}, \zh{心}, \zh{宀}, \zh{石}, \zh{立}, \zh{见}), écrivez-le grand, notez son nom (\emph{sī, ěr, rén, xīn, mián, shí, lì, jiàn}) et trouvez 2 autres caractères du manuel qui le contiennent (ex: \zh{纟} dans \zh{红} hóng, \zh{绿} lǜ ; \zh{心} dans \zh{想} xiǎng, \zh{恩} ēn). Temps : 10 min.
\end{experience}

\coursec{Compréhension du texte}

\courssub{Compréhension détaillée (Questions fermées et ouvertes)}

Relisez le texte \emph{互联网和我们的生活} calmement. Répondez aux questions suivantes en phrases complètes en français (ou en chinois pour les questions 6 à 8).

\begin{exercice}[Questions de compréhension]
\begin{enumerate}
\item D'après le premier paragraphe, quel est le nombre approximatif d'internautes en Chine ? Citez le passage chinois.
\item Quels sont les trois appareils/outils mentionnés pour aller sur internet ? (Caractères + pinyin)
\item Listez quatre activités que font les Chinois sur WeChat d'après le texte.
\item Au Cameroun, quelles sont les deux applications préférées des jeunes ? Donnez leurs noms chinois.
\item Que font « beaucoup d'élèves » (\zh{很多学生}) sur internet au Cameroun ? Citez trois actions précises.
\item Traduisez en français : \zh{互联网让世界变小了。}
\item Traduisez en chinois : « Les élèves camerounais peuvent se faire des amis chinois en ligne. » (Utilisez le vocabulaire du texte).
\item Completez la phrase du texte : \zh{上网的时候，我们要小心：不要把密码\_\_\_\_\_\_别人，也不要看太长时间的\_\_\_\_\_\_。}
\item Pourquoi l'auteur dit-il qu'internet est « un bon professeur et un bon ami » (\zh{好老师，也是好朋友}) ? Justifiez par deux exemples tirés du texte (un pour la Chine, un pour le Cameroun).
\item \textbf{Expression personnelle :} Êtes-vous d'accord avec le dernier conseil (\zh{我们要好好地用它}) ? Donnez un exemple personnel d'une bonne et d'une mauvaise utilisation d'internet.
\end{enumerate}
\end{exercice}

\begin{corrige}[Exercice Compréhension détaillée]
\begin{enumerate}
\item Plus d'un milliard (\zh{十亿多人}). Passage : \zh{在中国，有十亿多人上网。}
\item Téléphone (\zh{手机} shǒujī), ordinateur (\zh{电脑} diànnǎo). Le texte dit « 用手机和电脑 ».
\item Discuter (\zh{聊天} liáotiān), payer (\zh{付钱} fùqián), appeler un taxi (\zh{叫出租车} jiào chūzūchē).
\item Facebook (\zh{脸书} liǎnshū) et WhatsApp.
\item Regarder les cours du prof (\zh{看老师的课}), faire des exercices (\zh{做练习}), envoyer des e-mails à des amis (\zh{给朋友发邮件}).
\item Internet a rendu le monde plus petit. / Internet fait que le monde devient plus petit.
\item \zh{喀麦隆的学生可以在网上认识中国朋友。} (Phrase exacte du texte).
\item \zh{告诉} (gàosu) ; \zh{视频} (shìpín).
\item Professeur : étude en ligne, cours, exercices (Cameroun/Chine). Ami : se faire des amis transfrontaliers, discuter, partager (WeChat en Chine, WhatsApp/FB au Cameroun).
\item Réponse libre attendue : Ex. bonne utilisation : réviser ses cours, parler à sa famille loin. Mauvaise : donner son mot de passe, passer 5h sur TikTok, croire de fausses infos.
\end{enumerate}
\end{corrige}

\coursec{Collocations et production guidée}

\courssub{Entraînement aux collocations : transformer et combiner}

\begin{methode}[Construire des phrases « Internet » en 3 étapes]
Pour produire des phrases correctes et naturelles sur ce thème :
\begin{enumerate}
\item \textbf{Choisir le sujet} : \zh{我} (je), \zh{我们} (nous), \zh{学生} (les élèves), \zh{爸爸} (papa)…
\item \textbf{Choisir l'outil/lieu} : \zh{用手机} (avec téléphone), \zh{用电脑} (avec ordi), \zh{在网上} (sur internet), \zh{用微信} (sur WeChat).
\item \textbf{Choisir l'action (Collocation V+O)} : \zh{上网}, \zh{发信息}, \zh{看新闻}, \zh{买东西}, \zh{聊天}, \zh{付钱}, \zh{看视频}.
\item \textbf{Assembler} : Sujet + (Outil/Lieu) + Action. Ajouter fréquence (\zh{每天}, \zh{经常}) ou temps (\zh{昨天}, \zh{现在}) si besoin.
\end{enumerate}
Exemple : \zh{我} + \zh{用手机} + \zh{发信息} $\rightarrow$ \zh{我用手机发信息。} / \zh{我每天用手机发信息。}
\end{methode}

\begin{exercice}[Production écrite guidée : Mon journal numérique]
Rédigez un court paragraphe (6-8 phrases) en chinois pour décrire vos habitudes internet. Utilisez \textbf{au moins 8 mots du tableau de vocabulaire} et \textbf{3 collocations} de la liste. Suivez ce plan :
\begin{enumerate}
\item Fréquence globale : \zh{我每天……} / \zh{我经常……} / \zh{我有时候……}
\item Outil principal : \zh{我用手机……} / \zh{我用电脑……}
\item 3 activités précises (collocations) : \zh{发信息}, \zh{看视频}, \zh{上网}…
\item Contexte lieu : \zh{在网上……} / \zh{在学校……} / \zh{在家……}
\item Application préférée : \zh{我喜欢用微信/脸书/WhatsApp……}
\item Prudence : \zh{我小心，不……} (négatif \zh{不} ou \zh{别} / \zh{不要})
\end{enumerate}
\textbf{Contrainte :} Écrivez les caractères. Le pinyin et la traduction française sont optionnels mais recommandés pour l'auto-correction.
\end{exercice}

\begin{corrige}[Exercice Production écrite — Modèles de réponses]
\textbf{Modèle A (Niveau HSK 2-3 standard) :}
\zh{我每天都上网。我用手机发信息，看视频，也看新闻。我在网上学习中文。我喜欢用微信聊天。我小心，不告诉别人密码。}\\
\emph{Wǒ měi tiān dōu shàngwǎng. Wǒ yòng shǒujī fā xìnxī, kàn shìpín, yě kàn xīnwén. Wǒ zài wǎngshang xuéxí Zhōngwén. Wǒ xǐhuan yòng Wēixìn liáotiān. Wǒ xiǎoxīn, bù gàosu biérén mìmǎ.}\\
« Je vais sur internet tous les jours. J'envoie des messages et regarde des vidéos et les nouvelles avec mon téléphone. J'étudie le chinois en ligne. J'aime discuter sur WeChat. Je suis prudent, je ne dis pas mon mot de passe aux autres. »

\textbf{Modèle B (Contexte camerounais, niveau HSK 3) :}
\zh{我经常用电脑上网。在网上，我发邮件给老师，也看足球比赛。我用WhatsApp和朋友聊天。我不买东西网上，因为我没有银行卡。上网的时候，我小心，不看坏信息。}\\
\emph{Wǒ jīngcháng yòng diànnǎo shàngwǎng. Zài wǎngshang, wǒ fā yóujiàn gěi lǎoshī, yě kàn zúqiú bǐsài. Wǒ yòng WhatsApp hé péngyou liáotiān. Wǒ bù mǎi dōngxi wǎngshang, yīnwèi wǒ méiyǒu yínhángkǎ. Shàngwǎng de shíhou, wǒ xiǎoxīn, bù kàn huài xìnxī.}\\
« Je vais souvent sur internet avec l'ordinateur. En ligne, j'envoie des e-mails au prof et je regarde les matchs de foot. Je discute avec mes amis sur WhatsApp. Je n'achète pas en ligne car je n'ai pas de carte bancaire. Quand je surfe, je fais attention, je ne regarde pas de mauvaises infos. »

\textbf{Points de vigilance pour l'auto-correction :}
\begin{itemize}
\item Ordre : Sujet + Temps + Outil/Lieu + Verbe + Objet.
\item \zh{在网上} placé avant le verbe (pas à la fin).
\item Négation : \zh{不} + Verbe (habitude) / \zh{没} + Verbe (passé/absence) / \zh{别/不要} + Verbe (impératif/conseil).
\item \zh{因为}…\zh{所以}… (optionnel, niveau +).
\end{itemize}
\end{corrige}

\courssub{Production orale : Exposé « Mon application préférée »}

\begin{experience}[Présentation orale (2 minutes par élève)]
Préparez un exposé de 2 minutes sur \textbf{une seule application} que vous utilisez (WeChat, WhatsApp, TikTok, Facebook, YouTube, Google Maps, application bancaire MTN/Orange Money…). Structure imposée :
\begin{enumerate}
\item Nom chinois + pinyin de l'appli (\zh{微信 Wēixìn}, \zh{抖音 Dǒuyīn} pour TikTok Chine, \zh{地图 Dìtú} pour Maps…).
\item Fonction principale : \zh{它可以……} (Tā kěyǐ… « Elle peut… ») + 3 verbes/collocations.
\item Votre fréquence : \zh{我每天/每周用……次。}
\item Un avantage / Un inconvénient : \zh{优点是…… 缺点是……} (Vocabulaire bonus : \zh{方便} biànliàn pratique, \zh{快} kuài rapide, \zh{危险} wéixiǎn dangereux, \zh{浪费时间} làngfèi shíjiān perte de temps).
\item Conseil : \zh{建议大家……} (Jiànyì dàjiā… « Je conseille à tous de… »).
\end{enumerate}
L'enseignant note : prononciation des tons (sur les mots-clés), fluidité, richesse du vocabulaire thème, structure logique.
\end{experience}

\coursec{Synthèse et évaluation}

\begin{aretenir}[Bilan de la leçon 35 — Vocabulaire \& Compréhension : Internet]
\begin{itemize}
\item \textbf{Vocabulaire clé (15+ mots) :} \zh{互联网, 网, 上网, 网络, 网站, 手机, 电脑, 邮件, 发邮件, 信息, 发信息, 密码, 新闻, 视频, 微信, 脸书, 越来越, 认识, 了解, 小心, 告诉, 别人}. Maîtrise : caractères, pinyin (tons), nature, traduction.
\item \textbf{Structures grammaticales ancrées :} \zh{在} + lieu virtuel (\zh{网上}) + V ; \zh{把} + Objet + V (disposition) ; \zh{越来越} + Adj/V (progression) ; \zh{让} + O + V/Adj + \zh{了} (causatif) ; \zh{好好地} + V (manière).
\item \textbf{Collocations V+O indispensables :} \zh{上网, 发邮件, 发信息, 看新闻, 看视频, 买东西, 聊天, 付钱, 输密码, 登录}. Ne séparez pas le verbe de son objet !
\item \textbf{Radicaux à reconnaître :} \zh{纟} (réseau), \zh{亻} (humain/confiance), \zh{心} (info/cœur), \zh{宀} (secret/toit), \zh{石} (code/pierre), \zh{立} (station/debout), \zh{见} (vision).
\item \textbf{Compétences validées :} Lire/écouter un texte authentique bilingue Chine-Cameroun ; répondre à des questions fermées/ouvertes ; produire un paragraphe écrit cohérent ; présenter une application à l'oral.
\item \textbf{Socioculturel :} Différence d'écosystème (WeChat/Alipay/Baidu vs WhatsApp/Facebook/Google) ; paiement mobile omniprésent en Chine ; usage éducatif croissant au Cameroun ; vigilance sécurité (mot de passe, temps d'écran) universelle.
\end{itemize}
\end{aretenir}

\begin{exercice}[Auto-évaluation finale — Cochez vos acquis]
\begin{itemize}
\item[ ] Je sais écrire de mémoire 10 mots du thème (caractères + pinyin + sens).
\item[ ] Je peux lire le texte support à voix haute sans hésiter sur les tons.
\item[ ] Je réussis l'exercice de compréhension détaillée (8/10 minimum).
\item[ ] J'ai produit mon paragraphe « Mon journal numérique » (corrigé par le prof/paire).
\item[ ] J'ai présenté mon application préférée à l'oral (2 min).
\item[ ] Je connais la différence entre \zh{上网}, \zh{网上}, \zh{网络}, \zh{网站}.
\item[ ] Je sais expliquer le radical de \zh{信}, \zh{密}, \zh{码}, \zh{视} à un camarade.
\end{itemize}
\textbf{Objectif suivant (Leçon 36 - Grammaire) :} Structures complexes de la phrase sur internet : \zh{除了……以外} (hormis), \zh{不仅……而且……} (non seulement… mais aussi), comparatif \zh{比} pour débattre « Téléphone vs Ordinateur ».
\end{exercice}

\coursec{Texte support : lecture et écoute}

Dans la section précédente, nous avons découvert le thème du jour à travers des images et des questions : \emph{internet}. Vous avez déjà rencontré quelques mots clés comme \zh{互联网} (hùliánwǎng, internet) et \zh{上网} (shàngwǎng, se connecter). Nous allons maintenant passer à l'étape centrale de la leçon : la lecture d'un court texte original. La méthode est toujours la même : d'abord on écoute et on répète, ensuite on lit soi-même, puis on observe les structures pour bien comprendre.

\courssub{Le texte support}

Voici le texte de la leçon. Il compte cinq phrases et suit une progression logique simple : l'importance d'internet, ses usages, les achats en ligne, puis une mise en garde.

\begin{textebox}[Texte 35 : 互联网和我们的生活 — Internet et notre vie]
现在，互联网对我们的生活很重要。\\
Xiànzài, hùliánwǎng duì wǒmen de shēnghuó hěn zhòngyào.\\
\emph{Aujourd'hui, internet est très important pour notre vie.}\\[4pt]
很多人每天上网。\\
Hěn duō rén měitiān shàngwǎng.\\
\emph{Beaucoup de gens se connectent chaque jour.}\\[4pt]
我们用电脑和手机查资料、发邮件、聊天。\\
Wǒmen yòng diànnǎo hé shǒujī chá zīliào, fā yóujiàn, liáotiān.\\
\emph{Nous utilisons l'ordinateur et le téléphone portable pour chercher des informations, envoyer des e-mails et discuter.}\\[4pt]
网上购物也很方便。\\
Wǎngshàng gòuwù yě hěn fāngbiàn.\\
\emph{Les achats en ligne sont aussi très pratiques.}\\[4pt]
但是，我们不能花太多时间在网上，要注意休息。\\
Dànshì, wǒmen bù néng huā tài duō shíjiān zài wǎngshàng, yào zhùyì xiūxi.\\
\emph{Mais nous ne devons pas passer trop de temps sur internet, il faut penser à se reposer.}
\end{textebox}

\begin{aretenir}[Traduction globale]
Aujourd'hui, internet est très important pour notre vie. Beaucoup de gens se connectent chaque jour. Nous utilisons l'ordinateur et le téléphone portable pour chercher des informations, envoyer des e-mails et discuter. Les achats en ligne sont aussi très pratiques. Mais nous ne devons pas passer trop de temps sur internet, il faut penser à se reposer.
\end{aretenir}

\courssub{Lecture en écho : écouter d'abord, répéter ensuite}

Avant de lire seul, il faut entendre le texte. La \cle{lecture en écho} est la méthode la plus efficace pour fixer les tons et le rythme de la phrase chinoise.

\begin{methode}[Comment travailler la lecture en écho]
\begin{enumerate}
\item \textbf{Écoute globale.} L'enseignant lit le texte entier une première fois, à vitesse normale. Les élèves écoutent \emph{sans lire}, les yeux fermés si possible, pour saisir la mélodie générale.
\item \textbf{Lecture phrase par phrase.} L'enseignant lit une phrase ; les élèves répètent en chœur, en imitant les tons et les pauses. On répète chaque phrase deux ou trois fois.
\item \textbf{Zoom sur les mots difficiles.} On isole les mots nouveaux (\zh{互联网}, \zh{资料}, \zh{方便}) et on les répète seuls, puis dans leur phrase.
\item \textbf{Lecture en chaîne.} Chaque élève lit une phrase à voix haute, à tour de rôle ; l'enseignant corrige immédiatement les tons.
\item \textbf{Lecture complète.} Un ou deux élèves volontaires lisent le texte entier sans interruption.
\end{enumerate}
\end{methode}

\begin{attention}[Les tons à surveiller dans ce texte]
Plusieurs mots du texte piègent les francophones :
\begin{itemize}
\item \zh{上网} shàngwǎng : 4\textsuperscript{e} ton + 3\textsuperscript{e} ton. Ne dites pas *shángwáng.
\item \zh{手机} shǒujī : 3\textsuperscript{e} ton + 1\textsuperscript{er} ton. Le 3\textsuperscript{e} ton descend puis remonte.
\item \zh{资料} zīliào : 1\textsuperscript{er} ton + 4\textsuperscript{e} ton, contraste fort entre le haut et le bas.
\item \zh{方便} fāngbiàn : 1\textsuperscript{er} ton + 4\textsuperscript{e} ton.
\item \zh{不能} bù néng : 4\textsuperscript{e} ton + 2\textsuperscript{e} ton ; ici \zh{不} reste au 4\textsuperscript{e} ton car \zh{能} n'est pas au 4\textsuperscript{e} ton.
\item \zh{要注意} yào zhùyì : deux 4\textsuperscript{e} tons se succèdent ; marquez bien la chute sur chacun.
\end{itemize}
Rappel : devant un 4\textsuperscript{e} ton, \zh{不} passe au 2\textsuperscript{e} ton (bú), par exemple \zh{不是} bú shì. Dans notre texte, \zh{不能} garde bù au 4\textsuperscript{e} ton.
\end{attention}

\courssub{Observer le texte : la structure en quatre idées}

Un bon lecteur repère d'abord le \emph{plan} du texte. Notre texte est construit en quatre idées qui s'enchaînent logiquement, signalées par des mots-charnières : \zh{现在} (aujourd'hui) ouvre le sujet, \zh{也} (aussi) ajoute une idée, \zh{但是} (mais) introduit le contraste final.

\begin{popfigure}
\begin{tikzpicture}[node distance=0.45cm, every node/.style={font=\small}]
\node[draw=popblue, fill=popblueL, rounded corners, align=center, minimum width=3.1cm, minimum height=1.5cm] (a) {\textbf{1. Importance}\\ \zh{互联网对我们的生活}\\ \zh{很重要。}};
\node[draw=popgreen, fill=popgreenL, rounded corners, align=center, minimum width=3.1cm, minimum height=1.5cm, right=of a] (b) {\textbf{2. Usages}\\ \zh{查资料、发邮件、}\\ \zh{聊天}};
\node[draw=poporange, fill=poporangeL, rounded corners, align=center, minimum width=3.1cm, minimum height=1.5cm, right=of b] (c) {\textbf{3. Achats en ligne}\\ \zh{网上购物}\\ \zh{也很方便。}};
\node[draw=poppink, fill=poppinkL, rounded corners, align=center, minimum width=3.1cm, minimum height=1.5cm, right=of c] (d) {\textbf{4. Mise en garde}\\ \zh{不能花太多时间，}\\ \zh{要注意休息。}};
\draw[-{Stealth}, thick, popdark] (a) -- (b);
\draw[-{Stealth}, thick, popdark] (b) -- (c) node[midway, above, font=\footnotesize]{\zh{也}};
\draw[-{Stealth}, thick, popdark] (c) -- (d) node[midway, above, font=\footnotesize]{\zh{但是}};
\end{tikzpicture}
\legende{Structure du texte en quatre idées : importance, usages, achats en ligne, mise en garde.}
\end{popfigure}

\courssub{Repérage : mots déjà connus et mots nouveaux}

Exercice de lecture active : soulignez dans le texte les mots que vous connaissez déjà, et entourez les mots nouveaux. Voici le bilan attendu.

\begin{center}
\begin{tabularx}{\linewidth}{|X|X|X|X|}
\hline
\rowcolor{popblueL} Caractères & Pinyin & Nature & Traduction \\
\hline
现在 & xiànzài & nom de temps & maintenant, aujourd'hui \\
\hline
生活 & shēnghuó & nom & la vie \\
\hline
很 & hěn & adverbe & très \\
\hline
多 & duō & adjectif & nombreux, beaucoup \\
\hline
每天 & měitiān & nom de temps & chaque jour \\
\hline
我们 & wǒmen & pronom & nous \\
\hline
用 & yòng & verbe & utiliser \\
\hline
和 & hé & conjonction & et \\
\hline
也 & yě & adverbe & aussi \\
\hline
但是 & dànshì & conjonction & mais \\
\hline
时间 & shíjiān & nom & le temps \\
\hline
\end{tabularx}
\end{center}

\begin{center}
\begin{tabularx}{\linewidth}{|X|X|X|X|}
\hline
\rowcolor{poporangeL} Caractères (mots nouveaux) & Pinyin & Nature & Traduction \\
\hline
互联网 & hùliánwǎng & nom & internet \\
\hline
对 & duì & préposition & pour, envers \\
\hline
重要 & zhòngyào & adjectif & important \\
\hline
上网 & shàngwǎng & verbe & se connecter à internet \\
\hline
电脑 & diànnǎo & nom & ordinateur \\
\hline
手机 & shǒujī & nom & téléphone portable \\
\hline
查资料 & chá zīliào & verbe + nom & chercher des informations \\
\hline
发邮件 & fā yóujiàn & verbe + nom & envoyer un e-mail \\
\hline
聊天 & liáotiān & verbe & bavarder, discuter \\
\hline
网上购物 & wǎngshàng gòuwù & locution & achats en ligne \\
\hline
方便 & fāngbiàn & adjectif & pratique \\
\hline
花 & huā & verbe & dépenser, passer (du temps) \\
\hline
注意 & zhùyì & verbe & faire attention à \\
\hline
休息 & xiūxi & verbe & se reposer \\
\hline
\end{tabularx}
\end{center}

\courssub{Grammaire en contexte : la préposition 对}

Observons la première phrase du texte : \zh{互联网对我们的生活很重要。} Que remarquez-vous ? En français, on dit « internet est important \emph{pour} notre vie » ; en chinois, le mot \zh{对} (duì) se place \emph{avant} le complément, et le tout se met \emph{avant} l'adjectif.

\begin{propriete}[La structure 对 + nom/pronom + 很 + adjectif]
\cle{Structure} : Sujet + 对 + complément + 很 + adjectif.\\
Elle signifie « être ... pour quelqu'un / quelque chose ».
\begin{itemize}
\item \zh{汉语对我很重要。} Hànyǔ duì wǒ hěn zhòngyào. — Le chinois est très important pour moi.
\item \zh{运动对健康很有好处。} Yùndòng duì jiànkāng hěn yǒu hǎochu. — Le sport est très bon pour la santé.
\item \zh{这本书对学生很有用。} Zhè běn shū duì xuésheng hěn yǒuyòng. — Ce livre est très utile pour les élèves.
\end{itemize}
\cle{Négation} : on place \zh{不} devant l'adjectif : \zh{这件事对他不重要。} Zhè jiàn shì duì tā bù zhòngyào. — Cela n'est pas important pour lui.\\
\cle{Attention à l'ordre} : le groupe 对 + complément se place toujours \emph{avant} l'adjectif ou le verbe, jamais après comme en français.
\end{propriete}

\begin{attention}[Erreur fréquente des francophones]
Ne traduisez pas mot à mot « internet est très important pour notre vie » en *\zh{互联网很重要对我们的生活}. En chinois, le complément introduit par \zh{对} doit précéder l'adjectif : \zh{互联网对我们的生活很重要。} Retenez le schéma : \textbf{Sujet + 对 + complément + adjectif}.
\end{attention}

\courssub{Grammaire en contexte : 上网, un verbe séparable}

Le mot \zh{上网} (shàngwǎng, se connecter) est formé d'un verbe \zh{上} (monter, ici : accéder) et d'un nom \zh{网} (filet, réseau). C'est un \cle{verbe séparable} : on peut insérer des éléments \emph{entre} les deux caractères.

\begin{propriete}[Les verbes séparables de type verbe + objet]
Avec les verbes comme \zh{上网}, \zh{聊天} (liáotiān, bavarder), \zh{休息} (xiūxi, se reposer), le complément de durée ou de fréquence s'insère \emph{entre} le verbe et son objet intégré :
\begin{itemize}
\item \zh{我上一次网。} Wǒ shàng yí cì wǎng. — Je me connecte une fois.
\item \zh{他每天上两个小时网。} Tā měitiān shàng liǎng ge xiǎoshí wǎng. — Il se connecte deux heures par jour.
\item \zh{我们聊了一会儿天。} Wǒmen liáo le yíhuìr tiān. — Nous avons bavardé un moment.
\item \zh{你休息一下。} Nǐ xiūxi yíxià. — Repose-toi un peu.
\end{itemize}
\end{propriete}

\begin{attention}[Piège de traduction]
On dit \zh{上一次网} (shàng yí cì wǎng, se connecter une fois), et non *\zh{上网一次} employé de façon rigide comme en français « aller sur internet une fois ». De même, préférez \zh{花很多时间在网上} ou \zh{上很长时间网} pour « passer beaucoup de temps sur internet ». Le français place le complément après le verbe ; le chinois l'insère au milieu du verbe séparable.
\end{attention}

\courssub{Grammaire en contexte : 用 + outil + verbe}

Troisième observation : \zh{我们用电脑和手机查资料。} Le verbe \zh{用} (yòng, utiliser) introduit l'outil, et le tout se place \emph{avant} le verbe principal.

\begin{propriete}[La structure 用 + outil + verbe]
\cle{Structure} : Sujet + 用 + outil + verbe (+ complément). « Utiliser ... pour faire ... ».
\begin{itemize}
\item \zh{我用手机查资料。} Wǒ yòng shǒujī chá zīliào. — J'utilise mon téléphone pour chercher des informations.
\item \zh{他用电脑发邮件。} Tā yòng diànnǎo fā yóujiàn. — Il utilise l'ordinateur pour envoyer des e-mails.
\item \zh{我们用刀吃饭吗？不，我们用筷子吃饭。} Wǒmen yòng dāo chīfàn ma? Bù, wǒmen yòng kuàizi chīfàn. — Mangeons-nous avec un couteau ? Non, nous mangeons avec des baguettes.
\item \zh{她用手机给朋友打电话。} Tā yòng shǒujī gěi péngyou dǎ diànhuà. — Elle téléphone à ses amis avec son portable.
\end{itemize}
Comparaison avec le français : le français dit « chercher des informations \emph{avec} le téléphone » (outil après le verbe) ; le chinois place l'outil introduit par \zh{用} \emph{avant} le verbe.
\end{propriete}

\begin{exemplebox}[Les verbes du texte en action]
\zh{我用手机查资料。} Wǒ yòng shǒujī chá zīliào. — J'utilise mon téléphone pour chercher des informations.\\
\zh{老师用电脑发邮件。} Lǎoshī yòng diànnǎo fā yóujiàn. — Le professeur envoie des e-mails avec l'ordinateur.\\
\zh{晚上，我和朋友在网上聊天。} Wǎnshang, wǒ hé péngyou zài wǎngshàng liáotiān. — Le soir, je discute avec mes amis sur internet.\\
\zh{妈妈在网上买东西，很方便。} Māma zài wǎngshàng mǎi dōngxi, hěn fāngbiàn. — Maman fait des achats en ligne, c'est très pratique.
\end{exemplebox}

\courssub{Tableau récapitulatif des structures du texte}

\begin{center}
\begin{tabularx}{\linewidth}{|X|X|X|}
\hline
\rowcolor{popblueL} Structure & Exemple (caractères + pinyin) & Traduction \\
\hline
Sujet + 对 + compl. + 很 + adj. & \zh{互联网对我们的生活很重要。} Hùliánwǎng duì wǒmen de shēnghuó hěn zhòngyào. & Internet est très important pour notre vie. \\
\hline
Sujet + 用 + outil + verbe & \zh{我们用电脑查资料。} Wǒmen yòng diànnǎo chá zīliào. & Nous utilisons l'ordinateur pour chercher des informations. \\
\hline
Verbe séparable + élément inséré & \zh{上一次网} shàng yí cì wǎng & se connecter une fois \\
\hline
Sujet + 也 + 很 + adj. & \zh{网上购物也很方便。} Wǎngshàng gòuwù yě hěn fāngbiàn. & Les achats en ligne sont aussi très pratiques. \\
\hline
不能 + verbe (interdiction) & \zh{我们不能花太多时间。} Wǒmen bù néng huā tài duō shíjiān. & Nous ne devons pas passer trop de temps. \\
\hline
要 + verbe (il faut) & \zh{要注意休息。} Yào zhùyì xiūxi. & Il faut penser à se reposer. \\
\hline
\end{tabularx}
\end{center}

\courssub{Lecture individuelle et entraînement}

\begin{exoresolu}[Lecture et réemploi]
\textbf{Énoncé.} Lisez à voix haute la phrase suivante, puis traduisez-la en français : \zh{很多人每天上网，但是他们不能花太多时间在网上。}
\tcblower
\textbf{Solution.} Pinyin : Hěn duō rén měitiān shàngwǎng, dànshì tāmen bù néng huā tài duō shíjiān zài wǎngshàng.\\
Traduction : « Beaucoup de gens se connectent chaque jour, mais ils ne doivent pas passer trop de temps sur internet. »\\
Points de prononciation : bien marquer le 4\textsuperscript{e} ton de \zh{上} dans \zh{上网} ; garder \zh{不} au 4\textsuperscript{e} ton devant \zh{能} (2\textsuperscript{e} ton) ; faire une courte pause après la virgule, devant \zh{但是}.
\end{exoresolu}

\begin{exercice}[À vous de lire]
\begin{enumerate}
\item Lisez à voix haute le texte support en entier, en respectant les tons et les pauses.
\item Soulignez dans le texte tous les mots de temps (\zh{现在}, \zh{每天}...) et dites où ils se placent dans la phrase.
\item Traduisez : \zh{手机对年轻人很重要。} Shǒujī duì niánqīngrén hěn zhòngyào.
\item Complétez avec \zh{用} : \zh{我＿＿电脑发邮件。}
\end{enumerate}
\end{exercice}

\begin{corrige}[Exercice « À vous de lire »]
\begin{enumerate}
\item Lecture orale : vérifier les tons de \zh{上网} (4+3), \zh{资料} (1+4), \zh{方便} (1+4), \zh{休息} (1 + ton léger).
\item Les mots de temps \zh{现在} et \zh{每天} se placent juste après le sujet (ou en tête de phrase) : \zh{现在，互联网...} ; \zh{很多人每天上网。} En chinois, le complément de temps précède toujours le verbe.
\item « Le téléphone portable est très important pour les jeunes. » (structure 对 + complément + 很 + adjectif).
\item \zh{我用电脑发邮件。} Wǒ yòng diànnǎo fā yóujiàn. — J'utilise l'ordinateur pour envoyer des e-mails.
\end{enumerate}
\end{corrige}

\begin{aretenir}[Ce qu'il faut retenir de cette section]
\begin{itemize}
\item Le texte suit quatre idées : importance d'internet → usages quotidiens → achats en ligne → mise en garde (\zh{但是}).
\item \zh{对 + complément + 很 + adjectif} : le complément en \zh{对} précède l'adjectif, contrairement au français.
\item \zh{用 + outil + verbe} : l'outil se place avant le verbe.
\item \zh{上网} est un verbe séparable : \zh{上一次网}, pas *\zh{上网一次}.
\item La lecture en écho (écouter, répéter phrase par phrase, puis lire seul) est la clé pour maîtriser tons et rythme.
\end{itemize}
\end{aretenir}

\coursec{Mise en route et découverte du thème}

Bienvenue dans cette leçon consacrée au vocabulaire de l'internet et du monde numérique. Ce thème, au cœur du Module 5 « Mass médias et télécommunication », est indispensable pour décrire votre quotidien d'élève connecté, que vous soyez à Yaoundé, Douala, Buea ou dans n'importe quelle ville du Cameroun. Nous allons découvrir comment le chinois nomme les appareils, les actions et les objets du monde en ligne, en nous appuyant sur un texte authentique mettant en scène un élève camerounais.

\begin{experience}[Brainstorming : le monde numérique autour de moi]
En classe entière, au tableau, dressez la liste des mots français que vous associez à « internet » (ex. : wifi, mot de passe, télécharger, réseau social, fake news...). Le professeur note les mots. Ensuite, essayez de deviner comment dire « ordinateur » (indice : cerveau électrique) et « téléphone portable » (indice : machine de la main) en chinois.
\end{experience}

\coursec{Texte support : lecture et écoute}

Lisez le texte ci-dessous. Écoutez-le plusieurs fois (fichier audio OnBuch+). Repérez les mots que vous connaissez déjà et soulignez les nouveaux. La traduction est fournie pour vérifier votre compréhension globale.

\begin{textebox}[我和互联网 (Wǒ hé hùliánwǎng) — Moi et internet]
\zh{我叫恩东，是雅温得一所中学的学生。我有一部手机，家里也有一台电脑。每天放学以后，我先做作业，然后上网。我用电脑查资料，学习汉语；我也用手机跟朋友聊天，给他们发信息和邮件。}\\
\emph{Wǒ jiào Ēndōng, shì Yǎwēndé yì suǒ zhōngxué de xuésheng. Wǒ yǒu yì bù shǒujī, jiā lǐ yě yǒu yì tái diànnǎo. Měitiān fàngxué yǐhòu, wǒ xiān zuò zuòyè, ránhòu shàngwǎng. Wǒ yòng diànnǎo chá zīliào, xuéxí Hànyǔ; wǒ yě yòng shǒujī gēn péngyou liáotiān, gěi tāmen fā xìnxī hé yóujiàn.}\\[3pt]
\zh{互联网很方便。我可以在网上看新闻、下载音乐。我的姐姐喜欢网上购物，她说网上购物又快又方便。但是上网的时候，我们要小心：不要把密码告诉别人，也不要上不好的网站。}\\
\emph{Hùliánwǎng hěn fāngbiàn. Wǒ kěyǐ zài wǎngshàng kàn xīnwén, xiàzài yīnyuè. Wǒ de jiějie xǐhuan wǎngshàng gòuwù, tā shuō wǎngshàng gòuwù yòu kuài yòu fāngbiàn. Dànshì shàngwǎng de shíhou, wǒmen yào xiǎoxīn: bù yào bǎ mìmǎ gàosu biérén, yě bù yào shàng bù hǎo de wǎngzhàn.}\\[3pt]
\zh{网络是好东西，但是我们要好好用它。}\\
\emph{Wǎngluò shì hǎo dōngxi, dànshì wǒmen yào hǎohǎo yòng tā.}\\[3pt]
« Je m'appelle Ndong, je suis élève dans un lycée de Yaoundé. J'ai un téléphone portable, et à la maison il y a aussi un ordinateur. Chaque jour après les cours, je fais d'abord mes devoirs, puis je me connecte. J'utilise l'ordinateur pour chercher des documents et apprendre le chinois ; j'utilise aussi mon téléphone pour bavarder avec mes amis et leur envoyer des messages et des e-mails. Internet est très pratique. Je peux lire les nouvelles en ligne et télécharger de la musique. Ma grande sœur aime les achats en ligne ; elle dit que c'est rapide et pratique. Mais quand nous nous connectons, nous devons être prudents : il ne faut pas donner son mot de passe aux autres, ni visiter de mauvais sites. Le réseau est une bonne chose, mais il faut bien l'utiliser. »
\end{textebox}

\coursec{Glossaire du thème}

\courssub{Observation : les mots du texte}

Reprenons quelques phrases du texte support et soulignons les mots nouveaux qui structurent le thème.

\begin{exemplebox}[Les mots nouveaux en contexte]
\zh{我每天晚上上网。} \emph{Wǒ měitiān wǎnshang shàngwǎng.} « Je me connecte à internet tous les soirs. »\\[2pt]
\zh{我用电脑查资料。} \emph{Wǒ yòng diànnǎo chá zīliào.} « J'utilise l'ordinateur pour chercher des documents. »\\[2pt]
\zh{姐姐喜欢在网上购物。} \emph{Jiějie xǐhuan zài wǎngshàng gòuwù.} « Ma grande sœur aime faire des achats en ligne. »\\[2pt]
\zh{请告诉我你的密码。} \emph{Qǐng gàosu wǒ nǐ de mìmǎ.} « Dis-moi ton mot de passe, s'il te plaît. »
\end{exemplebox}

Observez : le caractère \zh{网} (wǎng) revient partout : \zh{上网}, \zh{网络}, \zh{网站}, \zh{网上}. C'est le cœur de la famille lexicale d'aujourd'hui. \zh{网} signifie à l'origine « filet » (le filet du pêcheur) ; par extension, il désigne tout « réseau », et donc internet.

\courssub{Tableau de vocabulaire du thème}

\begin{center}
\begin{tabularx}{\linewidth}{|X|X|X|X|}
\hline
\rowcolor{popblueL} Caractères & Pinyin & Nature & Traduction \\
\hline
互联网 & hùliánwǎng & n. & internet \\
\hline
网 & wǎng & n. & filet, réseau \\
\hline
网络 & wǎngluò & n. & réseau \\
\hline
上网 & shàngwǎng & v. & se connecter à internet \\
\hline
电脑 & diànnǎo & n. & ordinateur \\
\hline
手机 & shǒujī & n. & téléphone portable \\
\hline
邮件 & yóujiàn & n. & e-mail, courrier \\
\hline
发邮件 & fā yóujiàn & loc. v. & envoyer un e-mail \\
\hline
信息 & xìnxī & n. & information, message \\
\hline
资料 & zīliào & n. & documents, données \\
\hline
聊天 & liáotiān & v. & bavarder, chatter \\
\hline
网上购物 & wǎngshàng gòuwù & loc. v. & achats en ligne \\
\hline
方便 & fāngbiàn & adj. & pratique, commode \\
\hline
密码 & mìmǎ & n. & mot de passe \\
\hline
网站 & wǎngzhàn & n. & site web \\
\hline
下载 & xiàzài & v. & télécharger \\
\hline
\end{tabularx}
\end{center}

\begin{popfigure}
\begin{tikzpicture}[
  every node/.style={font=\small},
  racine/.style={rectangle, rounded corners, draw=popblue, fill=popblueL, thick, align=center, minimum width=2.2cm},
  fam/.style={rectangle, rounded corners, draw=poporange, fill=poporangeL, thick, align=center},
  mot/.style={rectangle, draw=popline, fill=popcream, rounded corners, align=center, font=\footnotesize}
]
\node[racine, align=center] (wang){网 wǎng\\filet, réseau};
\node[fam, above left=1.1cm and 2.6cm of wang] (app) {Appareils};
\node[fam, above=1.1cm of wang] (act) {Actions};
\node[fam, above right=1.1cm and 2.6cm of wang] (obj) {Objets numériques};
\node[mot, left=0.35cm of app, align=center] (a1){电脑 diànnǎo\\手机 shǒujī};
\node[mot, right=0.35cm of act, align=center] (a2){上网 shàngwǎng\\下载 xiàzài\\聊天 liáotiān\\发邮件 fā yóujiàn};
\node[mot, right=0.35cm of obj, align=center] (a3){网站 wǎngzhàn\\邮件 yóujiàn\\密码 mìmǎ\\信息 xìnxī};
\draw[-{Stealth}, thick, popblue] (wang) -- (app);
\draw[-{Stealth}, thick, popblue] (wang) -- (act);
\draw[-{Stealth}, thick, popblue] (wang) -- (obj);
\draw[-{Stealth}, poporange] (app) -- (a1);
\draw[-{Stealth}, poporange] (act) -- (a2);
\draw[-{Stealth}, poporange] (obj) -- (a3);
\end{tikzpicture}
\legende{Carte mentale lexicale : la famille du caractère 网 (wǎng) classée en trois groupes.}
\end{popfigure}

\courssub{Étude des mots : formation et emploi}

\textbf{1. 互联网 hùliánwǎng : internet.}\par
Ce mot est construit de manière très logique : \zh{互} (hù, « mutuel ») + \zh{联} (lián, « relier ») + \zh{网} (wǎng, « réseau ») = « le réseau qui relie mutuellement ». C'est le nom officiel d'internet en chinois. Dans la conversation courante, on dit souvent simplement \zh{网} : \zh{网上} (wǎngshàng) signifie « sur internet, en ligne ».

\begin{exemplebox}[互联网 et 网上]
\zh{互联网很方便。} \emph{Hùliánwǎng hěn fāngbiàn.} « Internet est très pratique. »\\[2pt]
\zh{我在网上看新闻。} \emph{Wǒ zài wǎngshàng kàn xīnwén.} « Je lis les nouvelles sur internet. »
\end{exemplebox}

\begin{definition}[网络 wǎngluò : le réseau]
\zh{网络} (wǎngluò) est formé de \zh{网} (wǎng, « filet ») et de \zh{络} (luò, « relier, connecter ») : c'est l'ensemble des « filets » reliés entre eux, c'est-à-dire \cle{le réseau}. On l'emploie dans \zh{网络信息} (wǎngluò xìnxī, « information en ligne ») ou \zh{网络购物} (wǎngluò gòuwù, « achats en ligne »). Distinction : \zh{网络} désigne le réseau en général (intranet, réseau local, internet), \zh{互联网} désigne internet précisément (le réseau mondial).
\end{definition}

\textbf{2. 上网 shàngwǎng : se connecter à internet.}\par
C'est un \cle{verbe séparable} (离合词) : \zh{上} (shàng, « monter ») + \zh{网} (wǎng, « le filet/réseau »). Littéralement : « monter sur le filet ». Entre les deux caractères, on peut insérer un complément de durée ou un aspect. Structure : Sujet + 上 + (durée/aspect) + 网.

\begin{exemplebox}[上网 en phrases]
\zh{你什么时候上网？} \emph{Nǐ shénme shíhou shàngwǎng ?} « Quand te connectes-tu ? »\\[2pt]
\zh{我每天上两个小时网。} \emph{Wǒ měitiān shàng liǎng ge xiǎoshí wǎng.} « Je me connecte deux heures par jour. »\\[2pt]
\zh{他上了网就不想下线。} \emph{Tā shàng le wǎng jiù bù xiǎng xiàxiàn.} « Une fois connecté, il ne veut plus se déconnecter. »
\end{exemplebox}

\textbf{3. 电脑 diànnǎo et 手机 shǒujī : les appareils.}\par
\zh{电脑} (diànnǎo) : \zh{电} (diàn, « électricité ») + \zh{脑} (nǎo, « cerveau ») = « cerveau électrique » = ordinateur.\\
\zh{手机} (shǒujī) : \zh{手} (shǒu, « main ») + \zh{机} (jī, « machine ») = « machine de la main » = téléphone portable.\\
Classificateurs : \zh{一台电脑} (yì tái diànnǎo), \zh{一部手机} (yì bù shǒujī) ou \zh{一个手机} (yí ge shǒujī).

\begin{attention}[电脑 vs 电话 vs 电视]
Ne confondez pas ces trois mots commençant par \zh{电} (électricité) :\\
\zh{电脑} (diànnǎo) = ordinateur (cerveau électrique).\\
\zh{电话} (diànhuà) = téléphone fixe (parole électrique).\\
\zh{电视} (diànshì) = télévision (vision électrique).\\
\emph{Erreur fréquente :} \zh{我用电话学习} (faux). \emph{Correct :} \zh{我用电脑学习。} (J'apprends avec mon ordinateur).
\end{attention}

\textbf{4. 邮件 yóujiàn et 发邮件 fā yóujiàn : l'e-mail.}\par
\zh{邮} (yóu) = poste. \zh{件} (jiàn) = classificateur d'objets/documents. \zh{邮件} = courrier, e-mail. Verbe d'envoi : \zh{发} (fā, « émettre, envoyer »). Structure : Sujet + 给 + Destinataire + 发邮件.

\begin{exemplebox}[邮件 en phrases]
\zh{我给老师发邮件。} \emph{Wǒ gěi lǎoshī fā yóujiàn.} « J'envoie un e-mail au professeur. »\\[2pt]
\zh{你收到我的邮件了吗？} \emph{Nǐ shōudào wǒ de yóujiàn le ma ?} « As-tu reçu mon e-mail ? » (\zh{收到} shōudào = recevoir avec succès).
\end{exemplebox}

\textbf{5. 信息 xìnxī et 资料 zīliào : information et documents.}\par
\zh{信息} (xìnxī) : information, message (ex. : un SMS = \zh{短信} duǎnxìn, une info = \zh{信息}).\\
\zh{资料} (zīliào) : documents, données, matériaux de travail (plus formel, plus volumineux).\\
Collocations : \zh{发信息} (envoyer un message), \zh{查资料} (chercher des documents).

\begin{exemplebox}[信息 et 资料]
\zh{请给我发一条信息。} \emph{Qǐng gěi wǒ fā yì tiáo xìnxī.} « Envoie-moi un message, s'il te plaît. » (\zh{条} tiáo = classif. pour messages/textes longs).\\[2pt]
\zh{我在网上查资料。} \emph{Wǒ zài wǎngshàng chá zīliào.} « Je cherche des documents sur internet. »
\end{exemplebox}

\textbf{6. 聊天 liáotiān : bavarder, chatter.}\par
\zh{聊} (liáo, « bavarder ») + \zh{天} (tiān, « ciel ») : littéralement « parler du ciel » (parler de tout et de rien). Verbe séparable : \zh{聊了半个小时天}. Avec une personne : \zh{跟/和 + quelqu'un + 聊天}.

\begin{exemplebox}[聊天]
\zh{我喜欢跟朋友聊天。} \emph{Wǒ xǐhuan gēn péngyou liáotiān.} « J'aime bavarder avec mes amis. »\\[2pt]
\zh{他们每天晚上在网上聊天。} \emph{Tāmen měitiān wǎnshang zài wǎngshàng liáotiān.} « Ils chattent en ligne tous les soirs. »
\end{exemplebox}

\textbf{7. 网上购物 wǎngshàng gòuwù : les achats en ligne.}\par
\zh{网上} (wǎngshàng, « en ligne ») + \zh{购物} (gòuwù, « faire des achats » : \zh{购} gòu « acheter » + \zh{物} wù « chose/objet »). En Chine, le e-commerce (淘宝 Táobǎo, 京东 Jīngdōng) est omniprésent ; au Cameroun, Jumia et les boutiques Facebook/WhatsApp développent cet usage à Douala et Yaoundé.

\textbf{8. 方便 fāngbiàn : pratique, commode.}\par
Adjectif prédicatif : Sujet + 很/真/挺 + 方便. \zh{便} se lit \textbf{biàn} (4\textsuperscript{e} ton), jamais pián ici.

\begin{exemplebox}[方便]
\zh{用手机上网很方便。} \emph{Yòng shǒujī shàngwǎng hěn fāngbiàn.} « Se connecter avec un téléphone est très pratique. »\\[2pt]
\zh{网上购物又快又方便。} \emph{Wǎngshàng gòuwù yòu kuài yòu fāngbiàn.} « Les achats en ligne sont à la fois rapides et pratiques. »
\end{exemplebox}

\textbf{9. 密码 mìmǎ : le mot de passe.}\par
\zh{密} (mì, « secret ») + \zh{码} (mǎ, « code, chiffre ») = « code secret ». Indispensable pour la sécurité numérique (téléphone, email, Mobile Money, banque).

\begin{exemplebox}[密码]
\zh{请告诉我你的密码。} \emph{Qǐng gàosu wǒ nǐ de mìmǎ.} « Dis-moi ton mot de passe, s'il te plaît. »\\[2pt]
\zh{我忘了我的密码。} \emph{Wǒ wàng le wǒ de mìmǎ.} « J'ai oublié mon mot de passe. »
\end{exemplebox}

\textbf{10. 网站 wǎngzhàn : le site web.}\par
\zh{网} (wǎng, « réseau ») + \zh{站} (zhàn, « station, gare ») = « station sur le réseau ». Classificateur : \zh{个} (ge). Ex. : \zh{一个网站} (yí ge wǎngzhàn).

\textbf{11. 下载 xiàzài : télécharger.}\par
\zh{下} (xià, « descendre ») + \zh{载} (zài, « charger, transporter ») = « faire descendre en chargeant » = télécharger (du serveur vers l'appareil). Antonyme : \zh{上传} (shàngchuán, « téléverser/uploader »).

\begin{savaistu}[Le ton de 载 dans 下载]
Le caractère \zh{载} a deux prononciations principales : \textbf{zài} (4\textsuperscript{e} ton) pour « charger, transporter, contenir » (\zh{载重} zàizhòng « charge utile », \zh{载人} zàirén « transporter des personnes ») et \textbf{zǎi} (3\textsuperscript{e} ton) pour « année » (\zh{载} zǎi = an) ou dans certains termes techniques d'enregistrement. Dans le vocabulaire HSK standard et l'usage courant d'internet (\zh{下载} « télécharger », \zh{下载速度} « vitesse de téléchargement »), la prononciation normative est \textbf{xiàzài} (4\textsuperscript{e} ton sur zài). Retenez : \zh{下载 xiàzài}.
\end{savaistu}

\courssub{Clés (radicaux) et ordre des traits}

Pour écrire correctement les caractères du thème, identifions leur clé (radical) et l'ordre des traits principal (règles : haut→bas, gauche→droite, horizontal avant vertical, extérieur avant intérieur).

\begin{center}
\begin{tabularx}{\linewidth}{|X|X|X|X|}
\hline
\rowcolor{popblueL} Caractère & Clé (Radical) & Nb traits & Ordre des traits (description) \\
\hline
网 & 纟 (sī, soie) & 6 & 1. 纟 (フ フ 丶) 2. 亡 (一 ノ 丶 丶) \\
\hline
电 & 田 (tián, champ) & 5 & 1. 田 (丨 一 丨 一 一) 2. 乙 (乙) — *simplifié* \\
\hline
脑 & 月 (yuè ⺼, viande) & 10 & 1. 月 (ノ ㇆ 一 一) 2. 亠 (一 丨) 3. 冂 (㇆ 一 丨) 4. 几 (㇆ 丶) \\
\hline
手 & 手 (shǒu, main) & 4 & 1. 横 (一) 2. 竖钩 (丨亅) 3. 提 (㇀) 4. 捺 (㇏) \\
\hline
机 & 木 (mù, bois) & 6 & 1. 木 (一 丨 丿 丶) 2. 幂 (㇆ 丨 ㇆) — *partie droite* \\
\hline
密 & 宀 (mián, toit) & 11 & 1. 宀 (丶 丶 ㇆) 2. 必 (一 丨 丿 丶 ㇆ 丶 丶) — *partie basse* \\
\hline
码 & 石 (shí, pierre) & 10 & 1. 石 (一 丿 丨 一 一) 2. 马 (㇆ 一 丨 一 丶) — *partie droite* \\
\hline
站 & 立 (lì, debout) & 10 & 1. 立 (丨 一 丿 丶 丨) 2. 占 (㇆ 丨 一 丨 一) — *partie basse* \\
\hline
载 & 车 (chē, véhicule) & 10 & 1. 车 (一 ㇆ 丨 一 一) 2. 才 (丿 丨 丶) — *partie haute droite* \\
\hline
\end{tabularx}
\end{center}

\begin{methode}[Écrire un caractère composé (ex. 网)]
\begin{enumerate}
\item Identifiez la clé : \zh{纟} (à gauche). Écrivez-la d'abord (3 traits).
\item Écrivez la partie droite \zh{亡} (4 traits) : horizontal, trait oblique, point, point.
\item Respectez l'alignement : la clé gauche est souvent plus étroite ; la partie droite occupe l'espace restant.
\end{enumerate}
\end{methode}

\courssub{Compréhension du texte}

Répondez aux questions suivantes en français (ou en chinois si possible) d'après le texte \zh{我和互联网}.

\begin{exercice}[Questions de compréhension]
\begin{enumerate}
\item Comment s'appelle l'élève et dans quelle ville habite-t-il ?
\item Quels sont les deux appareils numériques qu'il possède ?
\item Quelle est sa routine après les cours ? (Deux étapes)
\item À quoi utilise-t-il l'ordinateur ? Et le téléphone ? (Deux usages distincts)
\item Que fait sa grande sœur sur internet ? Quel est son avis ?
\item Quels sont les deux conseils de prudence donnés à la fin du texte ?
\item Trouvez dans le texte l'équivalent chinois de : « rapidité », « mot de passe », « site web », « documents ».
\end{enumerate}
\end{exercice}

\begin{corrige}[Exercice Compréhension]
\begin{enumerate}
\item Il s'appelle Ndong (恩东) et habite à Yaoundé (雅温得).
\item Un téléphone portable (一部手机) et un ordinateur (一台电脑).
\item 1. Faire ses devoirs (做作业). 2. Se connecter à internet (上网).
\item Ordinateur : chercher des documents (查资料) et apprendre le chinois (学习汉语). Téléphone : bavarder avec des amis (跟朋友聊天), envoyer des messages (发信息) et des e-mails (发邮件).
\item Elle fait des achats en ligne (网上购物). Elle trouve cela rapide et pratique (又快又方便).
\item 1. Ne pas donner son mot de passe aux autres (不要把密码告诉别人). 2. Ne pas visiter de mauvais sites (不要上不好的网站).
\item Rapidité : 快 (kuài) ; Mot de passe : 密码 (mìmǎ) ; Site web : 网站 (wǎngzhàn) ; Documents : 资料 (zīliào).
\end{enumerate}
\end{corrige}

\courssub{Collocations et production guidée}

Maîtriser un mot, c'est savoir avec quels autres mots il s'associe (collocations). Le chinois fonctionne beaucoup par paires Verbe + Objet (VO).

\begin{center}
\begin{tabularx}{\linewidth}{|X|X|X|}
\hline
\rowcolor{popblueL} Verbe (V) & Objet (O) & Collocation VO (Pinyin + Traduction) \\
\hline
上 (shàng) & 网 (wǎng) & \zh{上网} shàngwǎng — se connecter \\
\hline
发 (fā) & 邮件 (yóujiàn) / 信息 (xìnxī) & \zh{发邮件} fā yóujiàn — envoyer un mail / \zh{发信息} fā xìnxī — envoyer un message \\
\hline
查 (chá) & 资料 (zīliào) & \zh{查资料} chá zīliào — chercher des documents \\
\hline
聊 (liáo) & 天 (tiān) & \zh{聊天} liáotiān — bavarder \\
\hline
下载 (xiàzài) & 音乐 (yīnyuè) / 文件 (wénjiàn) / 软件 (ruǎnjiàn) & \zh{下载音乐} xiàzài yīnyuè — télécharger de la musique \\
\hline
网上购物 (wǎngshàng gòuwù) & — & \zh{网上购物} wǎngshàng gòuwù — faire du shopping en ligne (locution verbale) \\
\hline
看 (kàn) & 新闻 (xīnwén) / 视频 (shìpín) & \zh{看新闻} kàn xīnwén — lire les nouvelles \\
\hline
告诉 (gàosu) & 密码 (mìmǎ) & \zh{告诉密码} gàosu mìmǎ — donner/communiquer son mot de passe \\
\hline
\end{tabularx}
\end{center}

\begin{exercice}[Production guidée : Mon journal numérique]
Rédigez un court paragraphe (6--8 phrases) en chinois pour décrire votre utilisation d'internet. Utilisez \textbf{au moins 10 mots} du tableau de vocabulaire (dont 3 verbes séparables) et respectez l'ordre des mots (Sujet + Temps + Lieu + Verbe + Objet).
\begin{itemize}
\item Modèle de début : \zh{我叫……，我是……的学生。我有……} 
\item Contraintes : Utilisez \zh{先……然后……} (d'abord... puis...), \zh{因为……所以……} (parce que... donc...), \zh{虽然……但是……} (bien que... mais...).
\end{itemize}
\end{exercice}

\begin{corrige}[Exercice Production guidée (Exemple de réponse)]
\zh{我叫保罗，是杜阿拉一所中学的学生。我有一部手机和一台电脑。因为功课多，所以我每天放学先做作业，然后上网。我用电脑查资料，学习法语和汉语。虽然网上购物很方便，但我不用手机买东西。我常跟朋友发信息，聊天。上网的时候，我小心保护密码，不上坏网站。互联网很好，我们要好好用它。}\\
\emph{Wǒ jiào Bǎoluó, shì Dùā'lā yì suǒ zhōngxué de xuésheng. Wǒ yǒu yì bù shǒujī hé yì tái diànnǎo. Yīnwèi gōngkè duō, suǒyǐ wǒ měitiān fàngxué xiān zuò zuòyè, ránhòu shàngwǎng. Wǒ yòng diànnǎo chá zīliào, xuéxí Fǎyǔ hé Hànyǔ. Suīrán wǎngshàng gòuwù hěn fāngbiàn, dàn wǒ bù yòng shǒujī mǎi dōngxi. Wǒ cháng gēn péngyou fā xìnxī, liáotiān. Shàngwǎng de shíhou, wǒ xiǎoxīn bǎohù mìmǎ, bù shàng huài wǎngzhàn. Hùliánwǎng hěn hǎo, wǒmen yào hǎohǎo yòng tā.}
\end{corrige}

\courssub{Méthode : retenir le vocabulaire par familles}

\begin{methode}[Classer les mots pour mieux les mémoriser]
\begin{enumerate}
\item \textbf{Repérez le caractère-clé commun.} Ici, c'est \zh{网} (wǎng). Tous les mots qui le contiennent appartiennent au monde du réseau : \zh{网络}, \zh{上网}, \zh{网站}, \zh{网上}.
\item \textbf{Classez les mots en trois familles} : les \emph{appareils} (\zh{电脑}, \zh{手机}), les \emph{actions} (\zh{上网}, \zh{发邮件}, \zh{聊天}, \zh{下载}, \zh{网上购物}, \zh{查资料}), les \emph{objets numériques} (\zh{网站}, \zh{邮件}, \zh{信息}, \zh{资料}, \zh{密码}).
\item \textbf{Apprenez chaque mot avec sa collocation} : on n'apprend pas \zh{邮件} seul, mais \zh{发邮件} ; pas \zh{资料} seul, mais \zh{查资料} ; pas \zh{天} seul, mais \zh{聊天}.
\item \textbf{Fabriquez une phrase personnelle} avec chaque mot, ancrée dans votre vie (Yaoundé, Douala, votre lycée) : \zh{我用手机给朋友发邮件。}
\item \textbf{Révisez en cachant le pinyin} : lisez les caractères, prononcez à voix haute avec les tons, puis vérifiez.
\end{enumerate}
\end{methode}

\courssub{Exercice résolu}

\begin{exoresolu}[Complétez avec le mot qui convient]
Complétez chaque phrase avec le mot juste parmi : \zh{密码}, \zh{上网}, \zh{资料}, \zh{手机}, \zh{网站}.\\
1. 我每天晚上用电脑……。\\
2. 请告诉我你的……。\\
3. 我在网上查……，学习汉语。\\
4. 我有一部新的……。\\
5. 这个……很好，有很多信息。
\tcblower
\textbf{Solution détaillée.}\\
1. \zh{上网} : « Je me connecte à internet tous les soirs avec mon ordinateur. » Le verbe de connexion est \zh{上网} (verbe séparable, ici non séparé).\\
2. \zh{密码} : « Dis-moi ton mot de passe, s'il te plaît. » Le verbe \zh{告诉} (dire à qqn) attend une information comme objet : le code secret.\\
3. \zh{资料} : « Je cherche des documents sur internet pour apprendre le chinois. » La collocation fixe est \zh{查资料} (chercher des documents/données).\\
4. \zh{手机} : Le classificateur \zh{部} (bù) est spécifique aux téléphones portables (et machines). \zh{一部手机}.\\
5. \zh{网站} : « Ce site web est très bien, il y a beaucoup d'informations. » \zh{这个} (ceci) + nom singulier. \zh{网站} est le seul nom de lieu numérique de la liste.
\end{exoresolu}

\begin{attention}[Pièges fréquents des francophones — Check-list avant l'évaluation]
\begin{itemize}
\item \textbf{Tons de 上网} : \zh{shàngwǎng} = 4\textsuperscript{e} ton + 3\textsuperscript{e} ton. Ne dites pas \emph{shāngwǎng} (1\textsuperscript{e} ton sur shang).
\item \textbf{Ne traduisez pas « surfer sur internet » mot à mot} : pas de verbe « surfer ». Le chinois dit \zh{上网} ou \zh{上网看……} (se connecter pour regarder...).
\item \textbf{Classificateurs obligatoires en écriture soignée} : \zh{一台电脑} (yì tái diànnǎo) ; \zh{一部手机} (yì bù shǒujī). Évitez \zh{个} pour l'ordinateur.
\item \textbf{Verbe d'envoi : 发 (fā), pas 送 (sòng)}. \zh{发邮件}, \zh{发信息}. \zh{送} implique un déplacement physique (livrer, offrir).
\item \textbf{方便 est un adjectif statif} : \zh{很方便} (très pratique). Jamais \zh{是方便} (\emph{shì fāngbiàn}). Pas d'article « un » devant.
\item \textbf{Verbes séparables (上网, 聊天)} : Ne mettez pas le complément d'objet *entre* les deux caractères si c'est l'objet même du verbe (ex. \zh{上网} n'a pas d'objet). Mais insérez la durée \emph{entre} les deux : \zh{上了两个小时网}.
\item \textbf{网 vs 网络 vs 互联网} : \zh{网} = filet/réseau (court, oral) ; \zh{网络} = réseau (technique, général) ; \zh{互联网} = Internet (formel, précis).
\end{itemize}
\end{attention}

Vous disposez maintenant de tout le lexique, des structures grammaticales clés (verbes séparables, collocations VO), de la compréhension du texte et des bases de l'écriture des caractères du thème « internet ». La leçon suivante portera sur les structures grammaticales pour exprimer la fréquence, la durée et le but de vos activités en ligne.

\coursec{Clés (radicaux) et ordre des traits}

Dans la section précédente, nous avons découvert le glossaire du thème « internet » : \zh{网络} (wǎngluò, « réseau »), \zh{信息} (xìnxī, « information »), \zh{电脑} (diànnǎo, « ordinateur »), etc. Vous savez maintenant reconnaître et prononcer ces mots. Mais un sinisant ne s'arrête pas là : il doit aussi savoir \emph{écrire} ces caractères, correctement et dans le bon ordre. En chinois, l'ordre des traits n'est pas un détail de calligraphie : il conditionne la lisibilité, la rapidité d'écriture et même la recherche dans les dictionnaires classés par clés. Cette section vous donne les outils pour écrire proprement cinq caractères essentiels de notre thème : \zh{网}, \zh{络}, \zh{信}, \zh{电} et \zh{脑}.

\courssub{Observer avant d'écrire : la logique des caractères chinois}

\begin{exemplebox}[Observer d'abord : deux familles de caractères]
Regardez attentivement ces caractères du glossaire avant de lire la suite :
\begin{itemize}
\item \zh{网} wǎng « filet, réseau » : un seul bloc, avec un cadre extérieur et des éléments à l'intérieur.
\item \zh{信} xìn « message, lettre » : deux parties côte à côte, une à gauche, une à droite.
\item \zh{脑} nǎo « cerveau » : deux parties côte à côte, une étroite à gauche, une large à droite.
\end{itemize}
Question : dans quel sens traceriez-vous chaque caractère ? De haut en bas ? De gauche à droite ? De l'extérieur vers l'intérieur ? La réponse dépend de la \emph{structure} du caractère.
\end{exemplebox}

Un caractère chinois n'est pas un dessin arbitraire : c'est une construction organisée. La plupart des caractères sont composés d'une \cle{clé} (ou \cle{radical}, 部首 bùshǒu), qui donne une indication de sens, et d'une partie \emph{phonétique}, qui donne une indication de prononciation. Identifier la clé, c'est déjà deviner le sens général du caractère ; connaître l'ordre des traits, c'est pouvoir l'écrire de mémoire, sans modèle sous les yeux.

\begin{definition}[Clé (radical) et ordre des traits]
La \cle{clé} (部首 bùshǒu) est le composant graphique sous lequel un caractère est classé dans les dictionnaires ; elle indique le champ sémantique (la soie, l'homme, la chair, le filet…). L'\cle{ordre des traits} (笔顺 bǐshùn) est la succession réglée des mouvements du pinceau ou du stylo pour tracer un caractère. Chaque trait se trace dans une direction précise, et les traits s'enchaînent selon des règles générales.
\end{definition}

\courssub{Rappel des règles générales d'écriture}

Avant d'étudier nos cinq caractères, rappelons les quatre grandes règles qui gouvernent l'écriture de tous les caractères chinois.

\begin{propriete}[Les quatre règles d'or de l'ordre des traits]
\begin{enumerate}
\item \textbf{De haut en bas} (从上到下 cóng shàng dào xià) : on trace d'abord les éléments du haut, puis ceux du bas. Exemple : \zh{三} sān « trois » : trait supérieur, trait médian, trait inférieur.
\item \textbf{De gauche à droite} (从左到右 cóng zuǒ dào yòu) : dans un caractère composé de deux parties latérales, on écrit d'abord la partie gauche, puis la partie droite. Exemple : \zh{你} nǐ « tu » : 亻 puis 尔.
\item \textbf{L'extérieur avant l'intérieur} (先外后内 xiān wài hòu nèi) : dans un caractère à enclosure, on trace d'abord le cadre, puis le contenu ; la fermeture éventuelle du bas vient en dernier. Exemple : \zh{同} tóng « ensemble ».
\item \textbf{L'horizontal avant le vertical} (先横后竖 xiān héng hòu shù) : quand un trait horizontal et un trait vertical se croisent, l'horizontal passe d'abord. Exemple : \zh{十} shí « dix » : 一 puis 丨.
\end{enumerate}
\end{propriete}

\begin{attention}[Erreur fréquente des francophones]
En français, nous écrivons les lettres dans le sens de la lecture, sans règle de construction interne : chacun trace son « a » comme il veut. En chinois, \textbf{l'ordre des traits est obligatoire}, pas facultatif. Un caractère tracé dans le désordre est souvent déformé (les proportions se règlent trait après trait), plus lent à écrire, et illisible en écriture cursive. Prenez dès maintenant l'habitude de compter les traits et de verbaliser l'ordre en écrivant : « horizontal, puis vertical… ».
\end{attention}

\courssub{Étude caractère par caractère}

\textbf{1. 网 wǎng « filet, réseau » — clé 冂, 6 traits}\par

\zh{网} est à l'origine le dessin d'un filet de pêche : un cadre avec des mailles croisées à l'intérieur. Sa clé est \zh{冂} (jiōng), la clé du « cadre » ou de l'« enceinte », que l'on retrouve aussi dans \zh{同} tóng ou \zh{内} nèi « intérieur ». C'est un caractère à enclosure : la règle « extérieur avant intérieur » s'applique pleinement.

\begin{methode}[Comment écrire 网 en 6 traits]
\begin{enumerate}
\item Trait 1 : le vertical gauche du cadre 冂 (丨), de haut en bas.
\item Trait 2 : l'horizontal-vertical à droite (trait coudé (trait brisé)), qui forme le haut et le côté droit du cadre, en un seul mouvement.
\item Traits 3 et 4 : le premier 乂 à l'intérieur : le trait tombant à gauche (丿), puis le trait tombant à droite (㇏).
\item Traits 5 et 6 : le deuxième 乂 à l'intérieur, en dessous du premier : 丿 puis ㇏.
\end{enumerate}
Résumé : \textbf{le cadre d'abord, puis les deux 乂 de haut en bas}. Ne fermez rien en bas : le cadre 冂 reste ouvert.
\end{methode}

\textbf{2. 络 luò « relier, filet » — clé 纟, 9 traits}\par

\zh{络} est un caractère composé gauche-droite : à gauche la clé de la soie \zh{纟} (sī), à droite la partie phonétique \zh{各} (gè), qui rapproche la prononciation (gè → luò, la ressemblance était plus nette en chinois ancien). On écrit donc d'abord 纟 (3 traits), puis 各 (6 traits), de gauche à droite et, à l'intérieur de chaque partie, de haut en bas.

\begin{propriete}[La clé de la soie 纟, clé des fils et des réseaux]
Le radical 纟 (forme simplifiée de 糸, « soie ») apparaît dans de très nombreux caractères liés aux fils, aux tissus et, par extension, aux réseaux et aux liens :
\begin{itemize}
\item \zh{络} luò : relier, enfiler → \zh{网络} wǎngluò « réseau » ;
\item \zh{线} xiàn : fil, ligne → \zh{网线} wǎngxiàn « câble réseau », \zh{在线} zàixiàn « en ligne » ;
\item \zh{纸} zhǐ : papier (autrefois fabriqué à partir de fibres de soie) ;
\item \zh{红} hóng : rouge (la teinture des fils de soie) ; \zh{绿} lǜ : vert.
\end{itemize}
Retenir la clé 纟, c'est donc débloquer le sens de toute une famille de caractères. Quand vous rencontrez un caractère inconnu avec 纟 à gauche, pensez : « fil, lien, réseau ».
\end{propriete}

\textbf{3. 信 xìn « message, lettre, confiance » — clé 亻, 9 traits}\par

\zh{信} est composé de la clé de l'homme \zh{亻} (rén, forme latérale de 人) à gauche, et de \zh{言} (yán, « parole ») à droite : « la parole d'un homme », c'est-à-dire une parole à laquelle on peut se fier — d'où les deux sens de « message » et de « confiance ». Ordre d'écriture : d'abord 亻 (2 traits : 丿 puis 丨), puis 言 (7 traits : le point 丶, trois horizontaux 一 de haut en bas, puis la bouche 口 en bas). Règles combinées : gauche avant droite, haut avant bas.

\begin{exemplebox}[信息 : « la nouvelle du cœur »]
Le mot \zh{信息} xìnxī « information » est un composé éclairant :
\begin{itemize}
\item \zh{信} xìn : message, nouvelle ;
\item \zh{息} xī : souffle, respiration, et par extension « nouvelle » (composé de 自 zì « soi-même » et de 心 xīn « cœur » : le souffle qui vient du cœur).
\end{itemize}
Exemple : \zh{我给你发一条信息。} \emph{Wǒ gěi nǐ fā yì tiáo xìnxī.} « Je t'envoie un message. »\\
Autre exemple : \zh{网上的信息很多。} \emph{Wǎngshang de xìnxī hěn duō.} « Il y a beaucoup d'informations sur internet. »
\end{exemplebox}

\textbf{4. 电 diàn « électricité » — clé 丨 (ou 田), 5 traits}\par

\zh{电} représentait à l'origine un éclair dans un champ (田). On le classe sous la clé 丨 (trait vertical) ou, dans certains dictionnaires, sous 田. Ses 5 traits s'écrivent de haut en bas, et le trait de fermeture vient en dernier :
\begin{enumerate}
\item Trait 1 : le vertical gauche du 田 (丨).
\item Trait 2 : l'horizontal-vertical (trait coudé (trait brisé)), haut et côté droit.
\item Trait 3 : l'horizontal médian (一).
\item Trait 4 : le vertical médian (丨), qui descend sous la boîte.
\item Trait 5 : la fermeture horizontale du bas (一) — \textbf{toujours en dernier}, règle de l'enclosure.
\end{enumerate}
On retrouve 电 dans tout le vocabulaire moderne : \zh{电脑} diànnǎo « ordinateur » (litt. « cerveau électrique »), \zh{电话} diànhuà « téléphone » (« parole électrique »), \zh{电影} diànyǐng « cinéma » (« ombre électrique »).

\textbf{5. 脑 nǎo « cerveau » — clé 月, 10 traits}\par

\zh{脑} est un caractère gauche-droite : à gauche la clé \zh{月} (ici la « chair », et non la lune : c'est la forme latérale du radical de la viande 肉, présente dans presque tous les caractères du corps humain), à droite la partie phonétique composée de 亠 (tóu, « couvercle ») et de 凶 (xiōng). On écrit d'abord 月 (4 traits), puis la partie droite (6 traits), de haut en bas.

\begin{attention}[月 étroit à gauche : ne pas confondre avec 日]
Quand 月 est placé \textbf{à gauche} d'un caractère (comme dans 脑, 脸 liǎn « visage », 脚 jiǎo « pied »), il se trace \textbf{étroit et allongé}, environ deux fois moins large que sa forme autonome, pour laisser la place à la partie droite. Deux pièges pour les francophones :
\begin{itemize}
\item Ne pas écrire un 月 large comme dans 月 yuè « lune, mois » isolé : le caractère devient déséquilibré.
\item Ne pas confondre 月 (chair) avec 日 rì (soleil) : 日 est plus court et plus carré, 月 est allongé et son premier trait est un 丿 (tombant), pas un 丨 (vertical). Comparez : \zh{明} míng (日 + 月, « clair ») et \zh{脑} nǎo (月-chair + phonétique).
\end{itemize}
\end{attention}

\courssub{Schémas de décomposition}

\begin{popfigure}
\begin{tikzpicture}[node distance=0.8cm, every node/.style={font=\normalsize}]
\node[draw=popblue, fill=popblueL, rounded corners, minimum width=1.8cm, minimum height=1.1cm] (xin) {\zh{信} xìn};
\node[draw=popgreen, fill=popgreenL, rounded corners, below left=0.8cm and 0.5cm of xin, minimum width=1.4cm] (ren) {亻 homme};
\node[draw=poporange, fill=poporangeL, rounded corners, below right=0.8cm and 0.5cm of xin, minimum width=1.4cm] (yan) {言 parole};
\draw[-{Stealth[length=3mm]}, thick, popdark] (xin.south west) -- (ren.north east);
\draw[-{Stealth[length=3mm]}, thick, popdark] (xin.south east) -- (yan.north west);
\node[below=2.2cm of xin, text=popmuted, align=center] {« la parole d'un homme »\\[-2pt]→ message, confiance};

\node[draw=popblue, fill=popblueL, rounded corners, minimum width=1.8cm, minimum height=1.1cm, right=6cm of xin] (nao) {\zh{脑} nǎo};
\node[draw=popgreen, fill=popgreenL, rounded corners, below left=0.8cm and 0.5cm of nao, minimum width=1.4cm] (yue) {月 chair};
\node[draw=poporange, fill=poporangeL, rounded corners, below right=0.8cm and 0.5cm of nao, minimum width=1.6cm] (you) {亠 + 凶};
\draw[-{Stealth[length=3mm]}, thick, popdark] (nao.south west) -- (yue.north east);
\draw[-{Stealth[length=3mm]}, thick, popdark] (nao.south east) -- (you.north west);
\node[below=2.2cm of nao, text=popmuted, align=center] {partie du corps + phonétique\\[-2pt]→ cerveau};
\end{tikzpicture}
\legende{Décomposition de 信 (亻 + 言) et de 脑 (月 + 亠 + 凶) : la clé donne le sens, l'autre partie guide le sens ou le son.}
\end{popfigure}

\courssub{Tableau récapitulatif}

\begin{center}
\begin{tabularx}{\linewidth}{|X|X|X|X|}
\hline
\rowcolor{popblueL} Caractère & Clé (radical) & Nombre de traits & Règle dominante \\
\hline
\zh{网} wǎng « réseau » & 冂 (cadre) & 6 & extérieur avant intérieur \\
\hline
\zh{络} luò « relier » & 纟 (soie) & 9 & gauche avant droite \\
\hline
\zh{信} xìn « message » & 亻 (homme) & 9 & gauche avant droite, haut avant bas \\
\hline
\zh{电} diàn « électricité » & 丨 / 田 & 5 & haut avant bas, fermeture en dernier \\
\hline
\zh{脑} nǎo « cerveau » & 月 (chair) & 10 & gauche avant droite, 月 étroit \\
\hline
\end{tabularx}
\end{center}

\begin{aretenir}[L'essentiel]
\begin{itemize}
\item Chaque caractère possède une \cle{clé} qui indique son champ de sens : 纟 → fils et réseaux ; 亻 → la personne ; 月 à gauche → le corps ; 冂 → l'enclosure.
\item L'ordre des traits obéit à quatre règles : haut avant bas, gauche avant droite, extérieur avant intérieur, horizontal avant vertical ; la fermeture d'un cadre vient en dernier.
\item Comprendre la composition d'un caractère (信 = homme + parole ; 电脑 = cerveau électrique) est le meilleur moyen de mémoriser durablement le vocabulaire.
\end{itemize}
\end{aretenir}

\courssub{Exercice résolu et entraînement au tracé}

\begin{exoresolu}[Remettre les traits dans l'ordre]
Énoncé. Voici les traits du caractère \zh{信} dans le désordre : 口 — 一 (×3) — 亻 — 丶. Remettez les groupes de traits dans l'ordre correct d'écriture et justifiez avec les règles générales.
\tcblower
Solution détaillée. Ordre correct : 亻 → 丶 → 一 一 一 → 口.
\begin{enumerate}
\item \zh{信} est un caractère gauche-droite : la clé 亻 (homme) à gauche s'écrit d'abord (règle « gauche avant droite »).
\item Dans la partie droite 言, on écrit de haut en bas : le point 丶 en haut, puis les trois horizontaux 一, puis la bouche 口 en bas (règle « haut avant bas »).
\end{enumerate}
Vérification : 2 traits (亻) + 7 traits (言) = 9 traits au total, conforme au dictionnaire.
\end{exoresolu}

\begin{exercice}[Tracé guidé sur ardoise ou cahier]
Sur votre ardoise ou dans votre cahier, écrivez chaque caractère cinq fois en comptant les traits à voix haute :
\begin{enumerate}
\item \zh{网} : cadre 冂 d'abord, puis les deux 乂 de haut en bas (6 traits).
\item \zh{电} : la boîte 田, puis la fermeture en dernier (5 traits).
\item \zh{脑} : 月 étroit à gauche, puis la partie droite (10 traits).
\item Recopiez ensuite le mot \zh{电脑} diànnǎo « ordinateur » et la phrase \zh{我用电脑上网。} \emph{Wǒ yòng diànnǎo shàngwǎng.} « Je me connecte à internet avec mon ordinateur. »
\end{enumerate}
Critères de réussite : traits comptés correctement, 月 étroit dans 脑, fermeture de 电 en dernier, proportions équilibrées.
\end{exercice}

\begin{corrige}[Exercice de tracé guidé]
\begin{enumerate}
\item \zh{网} : 丨, (trait brisé), 丿, ㇏, 丿, ㇏ — 6 traits. Le cadre reste ouvert en bas.
\item \zh{电} : 丨, (trait brisé), 一, 丨, 一 (fermeture) — 5 traits.
\item \zh{脑} : 月 en 4 traits (丿, (trait brisé), 一, 一) tracé étroit, puis 亠 (丶, 一) et 凶 (丿, 丶, ㇄, 丨) — 10 traits au total.
\item Phrase modèle : \zh{我用电脑上网。} \emph{Wǒ yòng diànnǎo shàngwǎng.} « Je me connecte à internet avec mon ordinateur. » Vérifiez que 我 (7 traits, clé 戈) est tracé de haut en bas et que 网 garde son cadre ouvert.
\end{enumerate}
\end{corrige}

\begin{experience}[Atelier d'écriture en binômes]
Par binômes : l'élève A dicte un caractère du thème (\zh{网}, \zh{络}, \zh{信}, \zh{电}, \zh{脑}) en annonçant sa clé et son nombre de traits ; l'élève B l'écrit sur l'ardoise en verbalisant chaque trait (« horizontal, puis vertical… »). L'élève A corrige à l'aide du tableau récapitulatif. Échangez les rôles après trois caractères. Pour finir, chaque binôme écrit au tableau le mot \zh{网络} ou \zh{电脑} et la classe vérifie l'ordre des traits collectivement.
\end{experience}

\begin{savaistu}[Pourquoi les clés sont-elles si utiles ?]
Les dictionnaires chinois traditionnels classent les caractères sous 214 clés historiques (système du dictionnaire Kangxi). Aujourd'hui encore, pour chercher un caractère inconnu dont on ignore la prononciation, on identifie sa clé, on compte les traits restants, et on le retrouve dans l'index. Par exemple, pour chercher \zh{脑} : clé 月, puis 6 traits restants. C'est la même logique que notre dictionnaire alphabétique… mais fondée sur le dessin du caractère ! Au Cameroun comme en Chine, les élèves qui maîtrisent les clés apprennent les nouveaux caractères deux fois plus vite, car chaque caractère devient une petite histoire logique au lieu d'un dessin à mémoriser par cœur.
\end{savaistu}

\coursec{Compréhension du texte}

Après avoir appris à écrire les caractères clés du thème (\zh{网}, \zh{电}, \zh{脑}, \zh{机}…) grâce aux clés et à l'ordre des traits, nous revenons maintenant au texte support. L'objectif de cette section est de vérifier que vous le \cle{comprenez} vraiment : savoir lire, c'est savoir retrouver une information précise et la reformuler avec ses propres mots, en chinois correct.

\courssub{Rappel du texte support}

Avant de répondre aux questions, relisons attentivement le texte étudié dans les sections précédentes.

\begin{textebox}[互联网和我们的生活 — Hùliánwǎng hé wǒmen de shēnghuó]
现在，互联网对我们的生活非常重要。\\
Xiànzài, hùliánwǎng duì wǒmen de shēnghuó fēicháng zhòngyào.\\
« Aujourd'hui, internet est très important pour notre vie. »\\[4pt]
我们用电脑和手机上网。在网上，我们可以看新闻、听音乐、买东西，也可以和朋友聊天。\\
Wǒmen yòng diànnǎo hé shǒujī shàngwǎng. Zài wǎngshàng, wǒmen kěyǐ kàn xīnwén, tīng yīnyuè, mǎi dōngxi, yě kěyǐ hé péngyou liáotiān.\\
« Nous utilisons l'ordinateur et le téléphone pour nous connecter. Sur internet, nous pouvons lire les informations, écouter de la musique, acheter des choses, et aussi discuter avec des amis. »\\[4pt]
网上购物很方便：你不用出门，就可以买很多东西。在雅温得和杜阿拉，很多年轻人喜欢在网上买东西。\\
Wǎngshàng gòuwù hěn fāngbiàn: nǐ bú yòng chūmén, jiù kěyǐ mǎi hěn duō dōngxi. Zài Yǎwēndé hé Dùālā, hěn duō niánqīngrén xǐhuan zài wǎngshàng mǎi dōngxi.\\
« Les achats en ligne sont très pratiques : tu n'as pas besoin de sortir, tu peux acheter beaucoup de choses. À Yaoundé et à Douala, beaucoup de jeunes aiment acheter sur internet. »\\[4pt]
但是，我们不能花太多时间在网上。上网太多对眼睛不好，也对学习不好。我们应该每天上网一两个小时，不应该上五六个小时。\\
Dànshì, wǒmen bù néng huā tài duō shíjiān zài wǎngshàng. Shàngwǎng tài duō duì yǎnjing bù hǎo, yě duì xuéxí bù hǎo. Wǒmen yīnggāi měi tiān shàngwǎng yī liǎng ge xiǎoshí, bù yīnggāi shàng wǔ liù ge xiǎoshí.\\
« Mais nous ne devons pas passer trop de temps sur internet. Trop de connexion est mauvaise pour les yeux et aussi pour les études. Nous devrions nous connecter une ou deux heures par jour, et non cinq ou six heures. »\\[4pt]
互联网很有用，但是我们要好好用它。\\
Hùliánwǎng hěn yǒuyòng, dànshì wǒmen yào hǎohǎo yòng tā.\\
« Internet est très utile, mais nous devons bien l'utiliser. »
\end{textebox}

\courssub{Méthode : comment répondre à une question de compréhension}

En chinois comme en français, une bonne réponse de compréhension n'est jamais un mot isolé : c'est une \cle{phrase complète} construite à partir de la question et du texte.

\begin{methode}[Répondre à une question de compréhension en trois étapes]
\begin{enumerate}
\item \textbf{Repérer} dans le texte la phrase qui contient l'information demandée (chercher les mots-clés de la question : \zh{互联网}, \zh{电脑}, \zh{方便}…).
\item \textbf{Recopier la structure de la question} et remplacer le point d'interrogation par l'information du texte : la question donne souvent le squelette de la réponse.
\item \textbf{Reformuler} en phrase complète Sujet + Verbe + Complément, avec la bonne ponctuation chinoise (。), puis vérifier les tons et les mots-outils.
\end{enumerate}
\end{methode}

\begin{popfigure}
\begin{tikzpicture}[
  node distance=0.55cm and 1.1cm,
  qbox/.style={rectangle, rounded corners, draw=popblue, fill=popblueL, align=center, minimum width=3.2cm, minimum height=1cm, font=\small},
  sbox/.style={rectangle, rounded corners, draw=poporange, fill=poporangeL, align=center, minimum width=3.2cm, minimum height=1cm, font=\small},
  rbox/.style={rectangle, rounded corners, draw=popgreen, fill=popgreenL, align=center, minimum width=3.2cm, minimum height=1cm, font=\small},
  arr/.style={-{Stealth[length=3mm]}, thick, draw=popdark}
]
\node[qbox, align=center] (q){1. Question\\ \zh{网上购物方便吗？}};
\node[sbox, right=of q, align=center] (s){2. Repérage dans le texte\\ \zh{网上购物很方便}};
\node[rbox, right=of s, align=center] (r){3. Réponse complète\\ Sujet + Verbe + Complément\\ \zh{网上购物很方便。}};
\draw[arr] (q) -- (s);
\draw[arr] (s) -- (r);
\end{tikzpicture}
\legende{De la question à la réponse : repérer, recopier la structure, compléter.}
\end{popfigure}

\courssub{Observation : la grammaire cachée dans les questions}

Avant de répondre, observons deux points de grammaire présents dans les questions elles-mêmes.

\begin{propriete}[La question en 吗 (ma) et la réponse complète]
La particule \zh{吗} transforme une phrase affirmative en question fermée (oui/non) : on l'ajoute à la fin, sans changer l'ordre des mots.\\
Structure : \textbf{Phrase affirmative + 吗？}\\
Pour répondre, on \textbf{reprend le verbe ou l'adjectif de la question} (éventuellement nié), puis on développe en phrase complète. En chinois, on ne dit pas « oui » tout seul comme en français : on répète le verbe ou l'adjectif.
\begin{itemize}
\item \zh{网上购物方便吗？} (Wǎngshàng gòuwù fāngbiàn ma?) — Les achats en ligne sont-ils pratiques ?
\item \zh{方便。/ 网上购物很方便。} (Fāngbiàn. / Wǎngshàng gòuwù hěn fāngbiàn.) — Oui, ils sont pratiques.
\item \zh{不方便。} (Bù fāngbiàn.) — Non, ils ne sont pas pratiques.
\end{itemize}
\end{propriete}

\begin{propriete}[不能 (bù néng) : « ne pas pouvoir / ne pas devoir »]
\zh{不能} exprime l'\cle{impossibilité} ou l'\cle{interdiction} : « ne pas pouvoir, il ne faut pas ».\\
Structure : \textbf{Sujet + 不能 + Verbe (+ Complément)}\\
\zh{我们不能花太多时间在网上。} (Wǒmen bù néng huā tài duō shíjiān zài wǎngshàng.) — « Nous ne devons pas passer trop de temps sur internet. »\\
Négation de \zh{能} (pouvoir) ; à ne pas confondre avec \zh{不应该} (bù yīnggāi, « on ne devrait pas », conseil moral) et \zh{不用} (bú yòng, « pas besoin de »).
\end{propriete}

\begin{exemplebox}[Exemples traduits autour du texte]
\begin{itemize}
\item \zh{网上购物很方便。} (Wǎngshàng gòuwù hěn fāngbiàn.) — Les achats en ligne sont très pratiques.
\item \zh{你不用出门，就可以买东西。} (Nǐ bú yòng chūmén, jiù kěyǐ mǎi dōngxi.) — Tu n'as pas besoin de sortir pour acheter des choses.
\item \zh{我们不能花太多时间在网上。} (Wǒmen bù néng huā tài duō shíjiān zài wǎngshàng.) — Il ne faut pas passer trop de temps sur internet.
\item \zh{上网太多对眼睛不好。} (Shàngwǎng tài duō duì yǎnjing bù hǎo.) — Trop de connexion est mauvaise pour les yeux.
\end{itemize}
\end{exemplebox}

\begin{attention}[Pièges fréquents des francophones]
\begin{itemize}
\item À une question en \zh{吗}, ne répondez jamais seulement par un mot traduisant « oui » : le chinois n'a pas de mot unique pour « oui ». On répète le verbe ou l'adjectif : \zh{方便。} ou la phrase complète \zh{很方便。}
\item Ne dites pas \zh{是} (shì) pour répondre « oui » à une question dont le verbe n'est pas \zh{是} : \zh{方便吗？} → \zh{方便。}, pas \zh{是。}
\item Attention à l'ordre des mots : \zh{花太多时间在网上} (passer trop de temps sur internet) — le complément de lieu \zh{在网上} se place après le complément d'objet \zh{时间}, jamais avant le verbe comme en français (« passer sur internet trop de temps »).
\end{itemize}
\end{attention}

\courssub{Questions de compréhension : correction collective}

Nous répondons maintenant aux cinq questions. Pour chacune, nous citons d'abord la phrase du texte qui justifie la réponse, puis nous reformulons.

\begin{exoresolu}[Question 1 : 互联网对我们的生活怎么样？]
Énoncé : \zh{互联网对我们的生活怎么样？} (Hùliánwǎng duì wǒmen de shēnghuó zěnmeyàng?) — Que représente internet pour notre vie ?
\tcblower
\textbf{Repérage.} Le mot-clé est \zh{怎么样} (« comment ») ; la première phrase du texte contient la réponse : \zh{互联网对我们的生活非常重要。}\\
\textbf{Réponse.} \zh{互联网对我们的生活非常重要。} (Hùliánwǎng duì wǒmen de shēnghuó fēicháng zhòngyào.) — « Internet est très important pour notre vie. »\\
\textbf{Justification.} Le texte le dit dès la première phrase avec l'adverbe \zh{非常} (très, extrêmement) devant l'adjectif \zh{重要} (important). La structure \zh{对……怎么样} demande une appréciation : on répond avec \zh{对……很/非常 + adjectif}.
\end{exoresolu}

\begin{exoresolu}[Question 2 : 我们用电脑和手机做什么？]
Énoncé : \zh{我们用电脑和手机做什么？} (Wǒmen yòng diànnǎo hé shǒujī zuò shénme?) — Que faisons-nous avec l'ordinateur et le téléphone ?
\tcblower
\textbf{Repérage.} Mots-clés : \zh{电脑} et \zh{手机}. Le texte dit : \zh{我们用电脑和手机上网。在网上，我们可以看新闻、听音乐、买东西，也可以和朋友聊天。}\\
\textbf{Réponse.} \zh{我们用电脑和手机上网。在网上，我们可以看新闻、听音乐、买东西，也可以和朋友聊天。} (Wǒmen yòng diànnǎo hé shǒujī shàngwǎng. Zài wǎngshàng, wǒmen kěyǐ kàn xīnwén, tīng yīnyuè, mǎi dōngxi, yě kěyǐ hé péngyou liáotiān.) — « Nous utilisons l'ordinateur et le téléphone pour nous connecter. Sur internet, nous pouvons lire les informations, écouter de la musique, acheter des choses et discuter avec des amis. »\\
\textbf{Justification.} La question utilise \zh{用……做什么} (« utiliser… pour faire quoi ») ; la réponse reprend la même structure \zh{用 + outil + verbe}. Remarquez la virgule chinoise d'énumération \zh{、} entre les activités.
\end{exoresolu}

\begin{exoresolu}[Question 3 : 网上购物方便吗？]
Énoncé : \zh{网上购物方便吗？} (Wǎngshàng gòuwù fāngbiàn ma?) — Les achats en ligne sont-ils pratiques ?
\tcblower
\textbf{Repérage.} Le texte affirme : \zh{网上购物很方便：你不用出门，就可以买很多东西。}\\
\textbf{Réponse.} \zh{方便。网上购物很方便。你不用出门，就可以买很多东西。} (Fāngbiàn. Wǎngshàng gòuwù hěn fāngbiàn. Nǐ bú yòng chūmén, jiù kěyǐ mǎi hěn duō dōngxi.) — « Oui, c'est pratique. Les achats en ligne sont très pratiques : on n'a pas besoin de sortir pour acheter beaucoup de choses. »\\
\textbf{Justification.} Question fermée en \zh{吗} : on répond d'abord en répétant l'adjectif (\zh{方便}), puis on développe avec la phrase complète et la justification du texte (\zh{不用出门}). Le texte ajoute un exemple camerounais : à Yaoundé et à Douala, beaucoup de jeunes achètent en ligne.
\end{exoresolu}

\begin{exoresolu}[Question 4 : 我们能不能花太多时间在网上？]
Énoncé : \zh{我们能不能花太多时间在网上？} (Wǒmen néng bù néng huā tài duō shíjiān zài wǎngshàng?) — Pouvons-nous passer trop de temps sur internet ?
\tcblower
\textbf{Repérage.} Le texte dit : \zh{我们不能花太多时间在网上。上网太多对眼睛不好，也对学习不好。}\\
\textbf{Réponse.} \zh{不能。我们不能花太多时间在网上，因为上网太多对眼睛不好，也对学习不好。} (Bù néng. Wǒmen bù néng huā tài duō shíjiān zài wǎngshàng, yīnwèi shàngwǎng tài duō duì yǎnjing bù hǎo, yě duì xuéxí bù hǎo.) — « Non. Nous ne devons pas passer trop de temps sur internet, parce que c'est mauvais pour les yeux et pour les études. »\\
\textbf{Justification.} La question est de la forme affirmative-négative \zh{能不能} (verbe + négation + verbe), équivalente à \zh{能……吗}. On répond par \zh{不能} puis par la phrase complète. On ajoute la conjonction \zh{因为} (parce que) pour lier la raison, structure attendue à l'oral et à l'écrit. Le texte précise même la limite conseillée : une à deux heures par jour (\zh{一两个小时}), pas cinq ou six.
\end{exoresolu}

\begin{exoresolu}[Question 5 (ouverte) : 你常常上网吗？你上网做什么？]
Énoncé : \zh{你常常上网吗？你上网做什么？} (Nǐ chángcháng shàngwǎng ma? Nǐ shàngwǎng zuò shénme?) — Te connectes-tu souvent ? Pour quoi faire ?
\tcblower
\textbf{Modèle de réponse.} \zh{我常常上网。我上网看新闻、听音乐，也和朋友聊天。我不花太多时间在网上。} (Wǒ chángcháng shàngwǎng. Wǒ shàngwǎng kàn xīnwén, tīng yīnyuè, yě hé péngyou liáotiān. Wǒ bù huā tài duō shíjiān zài wǎngshàng.) — « Je me connecte souvent. Je me connecte pour lire les informations, écouter de la musique et discuter avec des amis. Je ne passe pas trop de temps en ligne. »\\
\textbf{Variante.} \zh{我不常常上网。我每天上网一个小时，学习汉语。} (Wǒ bù chángcháng shàngwǎng. Wǒ měi tiān shàngwǎng yī ge xiǎoshí, xuéxí Hànyǔ.) — « Je ne me connecte pas souvent. Je me connecte une heure par jour pour apprendre le chinois. »\\
\textbf{Conseil.} Question personnelle : le texte ne donne pas la réponse, mais il fournit tout le vocabulaire nécessaire (\zh{看新闻}, \zh{听音乐}, \zh{买东西}, \zh{聊天}). Réutilisez ces collocations et adaptez-les à votre situation.
\end{exoresolu}

\courssub{Tableau récapitulatif : de la question à la réponse}

\begin{center}
\begin{tabularx}{\linewidth}{|X|X|X|}
\hline
\rowcolor{popblueL} Type de question & Structure de la réponse & Exemple (question → réponse) \\
\hline
Question en \zh{吗} (oui/non) & Répéter le verbe/adjectif + phrase complète & \zh{方便吗？} → \zh{方便。/ 很方便。} \\
\hline
Question en \zh{能不能} & \zh{能。} ou \zh{不能。} + justification & \zh{能不能花太多时间？} → \zh{不能。} \\
\hline
Question en \zh{什么} (quoi) & Reprendre la structure + remplacer \zh{什么} par l'information & \zh{做什么？} → \zh{看新闻、听音乐。} \\
\hline
Question en \zh{怎么样} (comment) & \zh{很/非常 + adjectif} & \zh{怎么样？} → \zh{非常重要。} \\
\hline
\end{tabularx}
\end{center}

\begin{aretenir}[L'essentiel de la section]
\begin{itemize}
\item Une réponse de compréhension = phrase complète Sujet + Verbe + Complément, jamais un mot isolé.
\item La question donne le squelette de la réponse : on recopie sa structure et on remplace le mot interrogatif par l'information du texte.
\item \zh{吗} → réponse en répétant le verbe (\zh{方便。}) ; \zh{不能} = interdiction ou impossibilité (\zh{不能花太多时间}).
\item Toujours justifier sa réponse par une citation du texte.
\end{itemize}
\end{aretenir}

\begin{exercice}[À vous de répondre]
Répondez par des phrases complètes en chinois, en citant le texte.
\begin{enumerate}
\item \zh{在网上，我们可以做什么？} (citez au moins trois activités)
\item \zh{上网太多对什么不好？}
\item \zh{我们应该每天上网几个小时？}
\item \zh{你常常在网上买东西吗？} (réponse personnelle)
\end{enumerate}
\end{exercice}

\begin{corrige}[Exercice « À vous de répondre »]
\begin{enumerate}
\item \zh{在网上，我们可以看新闻、听音乐、买东西，也可以和朋友聊天。} — Le texte énumère quatre activités séparées par \zh{、}.
\item \zh{上网太多对眼睛不好，也对学习不好。} — Structure \zh{对……不好} reprise deux fois avec \zh{也}.
\item \zh{我们应该每天上网一两个小时，不应该上五六个小时。} — \zh{应该} (devoir, conseil) s'oppose à \zh{不应该}.
\item Réponse libre, par exemple : \zh{我常常在网上买东西，很方便。} ou \zh{我不常常在网上买东西。} — Toute phrase complète avec le vocabulaire du texte est acceptée.
\end{enumerate}
\end{corrige}

\coursec{Collocations et production guidée}

Après avoir analysé le texte de compréhension et répondu aux questions, nous passons maintenant à l'appropriation active du lexique et des structures. Cette section vise à transformer la compréhension passive en compétence de production écrite et orale, en ancrant les collocations les plus fréquentes et en guidant la rédaction d'un paragraphe personnel cohérent.

\courssub{Collocations clés du thème « Internet »}

En chinois, les verbes d'action sur internet forment des unités lexicales fixes (collocations verbe-objet, notées VO) qu'il faut mémoriser par blocs, et non mot à mot. Le tableau ci-dessous réunit les six collocations cibles de cette leçon, classées par fonction communicative.

\begin{center}
\begin{tabularx}{\linewidth}{|X|X|X|X|}
\hline
\rowcolor{popblueL} \textbf{Caractères} & \textbf{Pinyin} & \textbf{Nature} & \textbf{Traduction} \\
\hline
上网 & shàngwǎng & VO & se connecter (à internet), aller en ligne \\
\hline
发邮件 & fā yóujiàn & VO & envoyer un courriel \\
\hline
查资料 & chá zīliào & VO & chercher des informations, faire des recherches \\
\hline
聊天 & liáotiān & VO & discuter, chatter (oral ou écrit) \\
\hline
下载音乐 & xiàzǎi yīnyuè & VO & télécharger de la musique \\
\hline
网上购物 & wǎngshàng gòuwù & VO & faire des achats en ligne, acheter sur internet \\
\hline
\end{tabularx}
\end{center}

\begin{exemplebox}[Exemples contextuels]
\zh{我每天晚上上网一个小时。} (Wǒ měitiān wǎnshang shàngwǎng yí ge xiǎoshí.) — Je me connecte une heure chaque soir.

\zh{老师让我们查资料写作业。} (Lǎoshī ràng wǒmen chá zīliào xiě zuòyè.) — Le professeur nous demande de chercher des informations pour faire le devoir.

\zh{我用手机发邮件给同学。} (Wǒ yòng shǒujī fā yóujiàn gěi tóngxué.) — J'envoie un courriel à mon camarade avec mon téléphone.

\zh{周末我喜欢网上购物，很方便。} (Zhōumò wǒ xǐhuan wǎngshàng gòuwù, hěn fāngbiàn.) — Le week-end, j'aime faire des achats en ligne, c'est pratique.
\end{exemplebox}

\begin{attention}[Piège francophone : « surfer sur internet »]
En français, on utilise souvent le verbe « surfer » ou l'expression « aller sur internet ». En chinois, \textbf{上网} (shàngwǎng, litt. « monter sur le filet / le réseau ») est le verbe unique, neutre et standard pour « se connecter / aller en ligne ». N'essayez pas de traduire « surfer » par \zh{冲浪} (chōnglàng), qui désigne le sport (le surf) ou un usage très familier et vieilli. \textbf{上网} s'emploie seul ou suivi d'une durée (\zh{上网两个小时}), jamais avec une préposition du type « sur » : on ne dit pas \zh{*在网上上网}.
\end{attention}

\courssub{Structures utiles : exprimer l'outil (\cle{用}) et la fréquence}

Deux structures grammaticales sont indispensables pour parler de ses habitudes numériques : la préposition \cle{用} (yòng) pour marquer l'instrument, et le placement des adverbes de fréquence avant le verbe.

\begin{propriete}[Structure instrumentale : Sujet + 用 + Outil + Verbe (Objet)]
La préposition \textbf{用} (yòng, « utiliser, au moyen de ») introduit l'instrument ou le moyen par lequel l'action est réalisée. Le groupe \cle{用 + Outil} se place \textbf{avant} le verbe principal (ou la collocation verbale). L'outil est souvent \zh{电脑} (diànnǎo, ordinateur) ou \zh{手机} (shǒujī, téléphone portable / smartphone).
\begin{center}
\begin{tabularx}{\linewidth}{|X|X|X|}
\hline
\rowcolor{poporangeL} \textbf{Structure} & \textbf{Exemple (Caractères + Pinyin)} & \textbf{Traduction} \\
\hline
S + 用 + Outil + V(O) & \zh{我用电脑查资料。} (Wǒ yòng diànnǎo chá zīliào.) & Je cherche des informations sur mon ordinateur. \\
\hline
S + 用 + Outil + V(O) & \zh{他用手机聊天。} (Tā yòng shǒujī liáotiān.) & Il discute sur son téléphone. \\
\hline
S + 用 + Outil + V(O) & \zh{我们用微信聊天。} (Wǒmen yòng Wēixìn liáotiān.) & Nous discutons avec WeChat. \\
\hline
\end{tabularx}
\end{center}
\emph{Remarque :} en chinois, on « utilise » une application (\zh{微信} WeChat, \zh{QQ}) comme un outil avec \cle{用}, alors qu'en français on dit « via » ou « sur ».
\end{propriete}

\begin{propriete}[Adverbes de fréquence : Sujet + Adverbe de fréquence + Verbe]
Les adverbes de fréquence (\zh{每天} měitiān « chaque jour », \zh{常常} chángcháng « souvent », \zh{经常} jīngcháng « régulièrement », \zh{有时} yǒushí « parfois », \zh{从不} cóngbù « jamais ») se placent \textbf{immédiatement après le sujet et avant le verbe} (ou avant le groupe \cle{用 + Outil + Verbe}). C'est une différence majeure avec le français, où « souvent », « parfois » peuvent aller en fin de phrase.
\begin{center}
\begin{tabularx}{\linewidth}{|X|X|X|}
\hline
\rowcolor{popgreenL} \textbf{Structure} & \textbf{Exemple} & \textbf{Traduction} \\
\hline
S + 每天 + V & \zh{我每天上网。} (Wǒ měitiān shàngwǎng.) & Je me connecte chaque jour. \\
\hline
S + 常常 + 用 + Outil + V & \zh{我常常用手机聊天。} (Wǒ chángcháng yòng shǒujī liáotiān.) & Je discute souvent sur mon téléphone. \\
\hline
S + 有时 + V & \zh{我有时下载音乐。} (Wǒ yǒushí xiàzǎi yīnyuè.) & Je télécharge parfois de la musique. \\
\hline
\end{tabularx}
\end{center}
\end{propriete}

\begin{exemplebox}[Combinaison Fréquence + Outil + Action]
\zh{我每天用电脑查资料，常常用手机聊天。} (Wǒ měitiān yòng diànnǎo chá zīliào, chángcháng yòng shǒujī liáotiān.) — Chaque jour, je cherche des informations sur l'ordinateur, et je discute souvent sur le téléphone.

\zh{我姐姐从不网上购物，她觉得不安全。} (Wǒ jiějie cóngbù wǎngshàng gòuwù, tā juéde bù ānquán.) — Ma grande sœur n'achète jamais en ligne, elle trouve que ce n'est pas sûr.
\end{exemplebox}

\courssub{La conjonction \cle{因为} (yīnwèi) pour justifier son opinion}

Pour exprimer une opinion nuancée (avantages et inconvénients d'internet), la structure cause-conséquence est essentielle.

\begin{propriete}[Cause et conséquence : 因为 (yīnwèi) ... ，(所以 suǒyǐ) ...]
\textbf{因为} (yīnwèi, « parce que ») introduit la cause ; \textbf{所以} (suǒyǐ, « donc ») introduit la conséquence. Deux emplois sont corrects :
\begin{itemize}
\item \textbf{因为 + Cause ，所以 + Conséquence} : les deux conjonctions peuvent se combiner dans la même phrase (contrairement au français, où « parce que... donc... » est incorrect !).
\item \textbf{Conséquence ，因为 + Cause} : la proposition avec \cle{因为} peut aussi suivre l'opinion qu'elle justifie ; dans ce cas, on n'utilise pas \cle{所以}.
\end{itemize}
\begin{center}
\begin{tabularx}{\linewidth}{|X|X|X|}
\hline
\rowcolor{poppurpleL} \textbf{Structure} & \textbf{Exemple} & \textbf{Traduction} \\
\hline
Conséquence, 因为 + Cause & \zh{我觉得互联网很方便，因为可以查资料。} (Wǒ juéde hùliánwǎng hěn fāngbiàn, yīnwèi kěyǐ chá zīliào.) & Je trouve internet pratique, parce qu'on peut chercher des informations. \\
\hline
因为 + Cause, 所以 + Conséquence & \zh{因为网速慢，所以我不能看视频。} (Yīnwèi wǎngsù màn, suǒyǐ wǒ bù néng kàn shìpín.) & Comme la connexion est lente, je ne peux pas regarder de vidéos. \\
\hline
\end{tabularx}
\end{center}
\end{propriete}

\begin{attention}[Deux erreurs fréquentes avec 因为 / 所以]
\begin{itemize}
\item \textbf{Ne traduisez pas mot à mot le français « parce que... donc... »} : en français cette combinaison est fautive, mais en chinois \zh{因为……所以……} est la structure standard et tout à fait correcte. Exemple : \zh{因为上网太久对眼睛不好，所以我每天只上网两个小时。}
\item \textbf{Ne commencez pas une phrase par \cle{所以} sans cause exprimée} : \cle{所以} (« donc ») renvoie toujours à une cause déjà mentionnée (avec \cle{因为} ou dans le contexte). Une phrase isolée comme \zh{*所以我不上网。} en début de texte est incorrecte.
\end{itemize}
\end{attention}

\courssub{Modèle de texte : « Internet et moi »}

Voici un texte original, écrit par un élève camerounais fictif, qui mobilise le vocabulaire, la structure \cle{用}, les adverbes de fréquence et \cle{因为} / \cle{所以}. Il servira de modèle pour votre production.

\begin{textebox}[Modèle de production : 互联网和我 (Hùliánwǎng hé wǒ)]
\zh{我叫保罗，住在雅温得。我每天上网，因为功课需要查资料。} \\
\zh{我常常用电脑发邮件给老师，也用手机和同学聊天。} \\
\zh{我觉得互联网很方便，因为可以网上购物，不用去市场。} \\
\zh{周末我有时下载音乐，听非洲的歌和中国的歌。} \\
\zh{但是，上网太久对眼睛不好，所以我每天只上网两个小时。}

\textbf{Pinyin :} \\
Wǒ jiào Bǎoluó, zhù zài Yǎwēndé. Wǒ měitiān shàngwǎng, yīnwèi gōngkè xūyào chá zīliào. \\
Wǒ chángcháng yòng diànnǎo fā yóujiàn gěi lǎoshī, yě yòng shǒujī hé tóngxué liáotiān. \\
Wǒ juéde hùliánwǎng hěn fāngbiàn, yīnwèi kěyǐ wǎngshàng gòuwù, bùyòng qù shìchǎng. \\
Zhōumò wǒ yǒushí xiàzǎi yīnyuè, tīng Fēizhōu de gē hé Zhōngguó de gē. \\
Dànshì, shàngwǎng tài jiǔ duì yǎnjing bù hǎo, suǒyǐ wǒ měitiān zhǐ shàngwǎng liǎng ge xiǎoshí.

\textbf{Traduction :} \\
Je m'appelle Paul, j'habite à Yaoundé. Je me connecte chaque jour, parce que les devoirs exigent de chercher des informations. \\
J'utilise souvent l'ordinateur pour envoyer des courriels au professeur, et aussi le téléphone pour discuter avec mes camarades. \\
Je trouve internet très pratique, parce qu'on peut acheter en ligne, sans avoir besoin d'aller au marché. \\
Le week-end, je télécharge parfois de la musique, j'écoute des chansons africaines et chinoises. \\
Mais rester en ligne trop longtemps n'est pas bon pour les yeux, donc je ne me connecte que deux heures par jour.
\end{textebox}

\begin{aretenir}[Glossaire du texte modèle]
\begin{itemize}
\item \zh{雅温得} (Yǎwēndé) : Yaoundé (capitale du Cameroun).
\item \zh{功课} (gōngkè) : devoirs, leçons.
\item \zh{需要} (xūyào) : avoir besoin de, nécessiter.
\item \zh{不用} (bùyòng) : ne pas avoir besoin de, ne pas devoir (absence de nécessité).
\item \zh{非洲} (Fēizhōu) : Afrique.
\item \zh{眼睛} (yǎnjing) : yeux.
\item \zh{只} (zhǐ) : seulement (restriction, placé avant le verbe).
\item \zh{两个小时} (liǎng ge xiǎoshí) : deux heures (classificateur \zh{个} ; devant un classificateur, « deux » se dit \zh{两} et non \zh{二}).
\end{itemize}
\end{aretenir}

\courssub{Méthode de production guidée : 4 étapes}

\begin{methode}[Rédiger son paragraphe « Internet et moi »]
Suivez cette trame en 4 phrases types pour construire un texte cohérent, varié et correct.
\begin{enumerate}
\item \textbf{Phrase d'habitude (Fréquence + Action principale) :} utilisez \zh{每天} ou \zh{常常} + \zh{上网} (+ raison avec \zh{因为} si possible). \emph{Ex. :} \zh{我常常上网，因为喜欢查资料。}
\item \textbf{Phrase d'usage (Outil + Action précise) :} utilisez \zh{用} + \zh{电脑/手机} + \zh{发邮件/聊天/查资料}. \emph{Ex. :} \zh{我用手机和朋友聊天。}
\item \textbf{Phrase d'opinion (Sentiment + Justification) :} \zh{我觉得互联网……，因为……。} (Adjectifs utiles : \zh{方便} fāngbiàn « pratique », \zh{有意思} yǒu yìsi « intéressant », \zh{快} kuài « rapide », \zh{危险} wēixiǎn « dangereux ».) \emph{Ex. :} \zh{我觉得互联网很有意思，因为可以看视频。}
\item \textbf{Phrase de mise en garde (Mais + Limite) :} \zh{但是……，所以……。} (Vocabulaire : \zh{对眼睛不好} duì yǎnjing bù hǎo « mauvais pour les yeux », \zh{浪费时间} làngfèi shíjiān « faire perdre du temps », \zh{网速慢} wǎngsù màn « connexion lente ».) \emph{Ex. :} \zh{但是上网太久对眼睛不好，所以我只上网一个小时。}
\end{enumerate}
\textbf{Conseil :} écrivez d'abord les caractères, puis le pinyin au-dessus, puis vérifiez l'ordre des mots : Sujet — Fréquence — \zh{用} + Outil — Verbe — Complément.
\end{methode}

\begin{exoresolu}[Application de la méthode : rédaction guidée]
\textbf{Consigne :} complétez le paragraphe ci-dessous en choisissant le bon élément entre parenthèses pour obtenir un texte correct et cohérent.

\zh{我\_\_\_\_\_\_（叫/姓）马克，住在\_\_\_\_\_\_（杜阿拉/北京）。我\_\_\_\_\_\_（每天/从不）上网，\_\_\_\_\_\_（因为/所以）要查资料做功课。我\_\_\_\_\_\_（常常/从不）用\_\_\_\_\_\_（书/手机）发邮件，\_\_\_\_\_\_（也/但）用电脑聊天。\_\_\_\_\_\_（我觉得/因为）互联网很\_\_\_\_\_\_（方便/慢），\_\_\_\_\_\_（因为/所以）可以网上购物。\_\_\_\_\_\_（但是/和）上网太久\_\_\_\_\_\_（对/给）眼睛不好，\_\_\_\_\_\_（因为/所以）我控制时间。}

\tcblower
\textbf{Solution détaillée :}
\begin{enumerate}
\item \zh{叫} (jiào) — « Je m'appelle ». \zh{姓} (xìng) ne s'emploie qu'avec le nom de famille (\zh{我姓马}).
\item \zh{杜阿拉} (Dùālā) — contexte camerounais (Douala).
\item \zh{每天} (měitiān) — habitude positive ; \zh{从不} contredirait la raison « pour faire les devoirs ».
\item \zh{因为} (yīnwèi) — introduit la cause de l'action principale.
\item \zh{常常} (chángcháng) — fréquence pour l'action secondaire.
\item \zh{手机} (shǒujī) — outil numérique pour \zh{发邮件} ; \zh{书} (livre) n'est pas un outil numérique.
\item \zh{也} (yě) — ajout d'une deuxième action ; \zh{但} ne peut pas s'employer seul comme conjonction de phrase.
\item \zh{我觉得} (wǒ juéde) — introduction de l'opinion personnelle.
\item \zh{方便} (fāngbiàn) — adjectif positif attendu ; \zh{慢} est négatif et contredirait la suite.
\item \zh{因为} (yīnwèi) — justification de l'opinion (« parce qu'on peut... »).
\item \zh{但是} (dànshì) — transition adversative (inconvénient).
\item \zh{对} (duì) — structure \zh{对……不好} (« mauvais pour... ») ; \zh{给} est incorrect ici.
\item \zh{所以} (suǒyǐ) — conséquence logique : je limite mon temps.
\end{enumerate}
\textbf{Texte final :} \zh{我叫马克，住在杜阿拉。我每天上网，因为要查资料做功课。我常常用手机发邮件，也用电脑聊天。我觉得互联网很方便，因为可以网上购物。但是上网太久对眼睛不好，所以我控制时间。}

(Wǒ jiào Mǎkè, zhù zài Dùālā. Wǒ měitiān shàngwǎng, yīnwèi yào chá zīliào zuò gōngkè. Wǒ chángcháng yòng shǒujī fā yóujiàn, yě yòng diànnǎo liáotiān. Wǒ juéde hùliánwǎng hěn fāngbiàn, yīnwèi kěyǐ wǎngshàng gòuwù. Dànshì shàngwǎng tài jiǔ duì yǎnjing bù hǎo, suǒyǐ wǒ kòngzhì shíjiān.)
\end{exoresolu}

\courssub{Exercice de production personnelle}

\begin{exercice}[À vous de jouer : « Internet et moi »]
Rédigez un paragraphe de 5 à 6 phrases en chinois simplifié (avec le pinyin au-dessus si vous le souhaitez) en suivant la \textbf{méthode en 4 étapes} et la trame ci-dessous. Utilisez au minimum \textbf{4 collocations} du tableau de vocabulaire et les structures \cle{用} et \cle{因为}.

\textbf{Trame obligatoire :}
\begin{enumerate}
\item \zh{我\_\_\_\_\_\_上网。} (fréquence)
\item \zh{我用\_\_\_\_\_\_ \_\_\_\_\_\_。} (outil + collocation)
\item \zh{我觉得互联网\_\_\_\_\_\_，因为\_\_\_\_\_\_。} (opinion + cause)
\item \zh{但是\_\_\_\_\_\_，所以\_\_\_\_\_\_。} (inconvénient + conséquence)
\item (Optionnel) \zh{周末我\_\_\_\_\_\_。} (activité du week-end)
\end{enumerate}

\textbf{Contraintes :}
\begin{itemize}
\item Contexte : vous êtes élève à Yaoundé, Douala, Bafoussam, Buea... (citez votre ville).
\item Vocabulaire obligatoire : \zh{上网}, \zh{用}, \zh{因为}, \zh{但是}, \zh{所以}, \zh{觉得}.
\item Minimum 4 collocations distinctes parmi les 6 du tableau.
\end{itemize}
\end{exercice}

\begin{corrige}[Exercice production personnelle]
\textbf{Exemple de production attendue (élève à Bafoussam) :}

\zh{我叫玛丽，住在巴富萨姆。} (Wǒ jiào Mǎlì, zhù zài Bāfùsàmǔ.) \\
\zh{我每天上网，因为要查资料。} (Wǒ měitiān shàngwǎng, yīnwèi yào chá zīliào.) \\
\zh{我常常用手机聊天，也用电脑发邮件。} (Wǒ chángcháng yòng shǒujī liáotiān, yě yòng diànnǎo fā yóujiàn.) \\
\zh{我觉得互联网很有意思，因为可以看中国电影。} (Wǒ juéde hùliánwǎng hěn yǒu yìsi, yīnwèi kěyǐ kàn Zhōngguó diànyǐng.) \\
\zh{但是上网太久浪费时间，所以我每天只上网一个小时。} (Dànshì shàngwǎng tài jiǔ làngfèi shíjiān, suǒyǐ wǒ měitiān zhǐ shàngwǎng yí ge xiǎoshí.) \\
\zh{周末我有时下载音乐。} (Zhōumò wǒ yǒushí xiàzǎi yīnyuè.)

\textbf{Critères de réussite (auto-évaluation) :}
\begin{itemize}
\item 5 phrases minimum.
\item Ville camerounaise citée.
\item 4 collocations cibles utilisées correctement (blocs VO intacts).
\item Structure \zh{用 + Outil + Verbe} correcte (ordre des mots).
\item \zh{因为} place la cause au bon endroit.
\item \zh{但是 / 所以} articulent la concession.
\item Pinyin correct (tons) si fourni.
\end{itemize}
\end{corrige}

\courssub{Activité orale : relecture par les pairs et correction collective}

\begin{experience}[Atelier « Correcteur du net »]
\textbf{Objectif :} développer l'esprit critique linguistique et la correction phonologique (tons).

\textbf{Déroulement (15 min) :}
\begin{enumerate}
\item \textbf{Échange (5 min) :} en binôme, échangez votre production écrite. Lisez le texte de votre camarade à voix haute (entraînement oral).
\item \textbf{Grille de relecture (5 min) :} cochez les points suivants sur la copie de votre pair :
    \begin{itemize}
    \item Les 4 collocations sont présentes et bien écrites (caractères).
    \item L'ordre \zh{Sujet + 用 + Outil + Verbe} est respecté.
    \item \zh{因为} introduit bien la cause ; aucune phrase ne commence par \zh{所以} sans cause exprimée.
    \item Les tons du pinyin (s'il est écrit) semblent corrects : surveillez \zh{shàngwǎng} (4-3), \zh{yīnwèi} (1-4), \zh{duì} (4).
    \item Le texte mentionne une ville du Cameroun.
    \end{itemize}
\item \textbf{Mise en commun (5 min) :} le professeur relève 3 erreurs fréquentes au tableau (ex. : \zh{*我上网用手机} au lieu de \zh{我用手机上网} ; \zh{*我觉得方便因为……} sans sujet \zh{互联网} dans la proposition d'opinion ; confusion entre \zh{发} fā « envoyer » et \zh{查} chá « chercher »). Correction collective et explication rapide.
\end{enumerate}
\end{experience}

\begin{savaistu}[Le clavier pinyin et les méthodes de saisie]
Au Cameroun comme en Chine, on tape le chinois sur ordinateur ou smartphone grâce aux \textbf{méthodes de saisie par pinyin} (\zh{拼音输入法} pīnyīn shūrùfǎ). On tape les lettres de l'alphabet (ex. : \texttt{shangwang}), et le logiciel propose les caractères correspondants (\zh{上网}...) ; il faut alors choisir le bon caractère (touche 1, 2, 3...).
\begin{itemize}
\item \textbf{Sogou Pinyin} (\zh{搜狗拼音}) et \textbf{Microsoft Pinyin} (\zh{微软拼音}) sont les plus utilisés en Chine.
\item Sur Android et iOS : claviers \textbf{Gboard}, \textbf{Google Pinyin} ou clavier natif « Chinois (simplifié) — Pinyin ».
\item \textbf{Astuce :} pour taper la voyelle \zh{ü}, on tape la lettre \texttt{v} (ex. : \texttt{nv} pour \zh{女} nǚ, \texttt{lv} pour \zh{绿} lǜ).
\item \textbf{Contexte camerounais :} dans les cybercafés de Yaoundé ou de Douala, les claviers sont souvent AZERTY français. Il faut ajouter la langue « Chinois (simplifié) » dans les paramètres Windows, puis basculer d'une langue à l'autre avec la combinaison Win + Espace.
\end{itemize}
\end{savaistu}

\begin{popfigure}
\begin{tikzpicture}[
    grow cyclic,
    every node/.style={align=center},
    root/.style={circle, fill=popblue, text=white, font=\bfseries\large, minimum size=2.2cm},
    outil/.style={rectangle, rounded corners, fill=poporange, text=white, font=\bfseries, minimum width=2.4cm, minimum height=1cm},
    action/.style={rectangle, rounded corners, fill=popgreen, text=white, font=\small, minimum width=2.6cm, minimum height=0.9cm},
    lien/.style={thick, popmuted}
]
\node[root, align=center] (yong) at (0,0){用 \\ (yòng)};
\node[outil, align=center] (pc) at (-4.5,2.5){电脑 \\ (diànnǎo)};
\node[outil, align=center] (tel) at (4.5,2.5){手机 \\ (shǒujī)};
\node[action, align=center] (a1) at (-7.5,5){查资料 \\ (chá zīliào)};
\node[action, align=center] (a2) at (-4.5,5.5){发邮件 \\ (fā yóujiàn)};
\node[action, align=center] (a3) at (-1.5,5){网上购物 \\ (wǎngshàng gòuwù)};
\node[action, align=center] (b1) at (1.5,5){聊天 \\ (liáotiān)};
\node[action, align=center] (b2) at (4.5,5.5){发邮件 \\ (fā yóujiàn)};
\node[action, align=center] (b3) at (7.5,5){下载音乐 \\ (xiàzǎi yīnyuè)};
\draw[lien] (yong) -- (pc);
\draw[lien] (yong) -- (tel);
\draw[lien] (pc) -- (a1);
\draw[lien] (pc) -- (a2);
\draw[lien] (pc) -- (a3);
\draw[lien] (tel) -- (b1);
\draw[lien] (tel) -- (b2);
\draw[lien] (tel) -- (b3);
\end{tikzpicture}
\legende{Carte mentale : collocations verbales classées par outil (\zh{用电脑} / \zh{用手机}). L'action commune \zh{发邮件} apparaît sur les deux branches.}
\end{popfigure}

\newpage

\coursec{Activité d'intégration}

\courssub{Objectifs de l'activité}

À la fin de cette activité d'intégration, tu seras capable de :

\begin{itemize}
\item mobiliser le \cle{vocabulaire du thème « internet »} (上网, 网站, 查资料, 聊天, 发邮件, 方便, 重要…) dans un texte personnel ;
\item employer correctement les \cle{structures de fréquence} (常常, 每天, 有时候, 很少) et les \cle{structures d'opinion} (我觉得…, …很方便, …很重要) ;
\item rédiger un court texte cohérent de 5 à 7 phrases intitulé \zh{我和互联网} ;
\item lire ton texte à voix haute avec une prononciation correcte (pinyin et tons) et \cle{comprendre les questions} d'un camarade pour y répondre oralement ;
\item écouter le portrait d'un binôme et poser deux questions de compréhension pertinentes.
\end{itemize}

\courssub{Situation}

Le club informatique de ton lycée à Yaoundé prépare la « Semaine du numérique ». La professeure de chinois propose au club d'ouvrir un \cle{mur des portraits numériques} : chaque élève de Première rédige en chinois un petit texte intitulé \zh{我和互联网} (« Internet et moi ») pour raconter sa vie numérique. Ces portraits seront affichés dans la salle multimédia, et les délégués d'une classe de chinois d'un lycée partenaire de Douala viendront les lire. Pour être affiché, ton portrait doit être clair, correct et personnel : il doit dire \textbf{à quelle fréquence} tu te connectes, \textbf{avec quel appareil}, \textbf{ce que tu fais} sur internet, \textbf{ton opinion} sur internet, et contenir \textbf{un conseil ou une mise en garde} personnelle.

Pour t'inspirer, voici le portrait rédigé par l'annonce officielle du club :

\begin{textebox}[L'annonce du club informatique]
同学们好！\\
Tóngxuemen hǎo !\\
« Bonjour chers camarades ! »\\[4pt]
我们俱乐部要做一个“我和互联网”的墙。\\
Wǒmen jùlèbù yào zuò yí ge « Wǒ hé hùliánwǎng » de qiáng.\\
« Notre club va créer un mur “Internet et moi”. »\\[4pt]
请你写五到七个句子：你常常上网吗？你用什么上网？你在网上做什么？你觉得互联网怎么样？\\
Qǐng nǐ xiě wǔ dào qī ge jùzi : nǐ chángcháng shàngwǎng ma ? Nǐ yòng shénme shàngwǎng ? Nǐ zài wǎngshang zuò shénme ? Nǐ juéde hùliánwǎng zěnmeyàng ?\\
« Écris cinq à sept phrases : est-ce que tu te connectes souvent ? Avec quoi te connectes-tu ? Que fais-tu sur internet ? Que penses-tu d'internet ? »\\[4pt]
最好的文章会贴在多媒体教室的墙上。\\
Zuì hǎo de wénzhāng huì tiē zài duōméitǐ jiàoshì de qiángshang.\\
« Les meilleurs textes seront affichés au mur de la salle multimédia. »
\end{textebox}

\courssub{Consignes}

\textbf{Tâche 1 — Comprendre la consigne du club.} Lis l'annonce du club informatique et réponds en français : combien de phrases faut-il écrire ? Quelles sont les quatre questions auxquelles ton portrait doit répondre ? Où seront affichés les meilleurs textes ?

\textbf{Tâche 2 — Préparer ton lexique.} Dans ton cahier, complète la grille suivante avec le vocabulaire de la leçon (en caractères et en pinyin) :

\begin{center}
\begin{tabularx}{\linewidth}{|X|X|}
\hline
\rowcolor{popblueL} Question & Ma réponse (mots chinois) \\
\hline
Fréquence de connexion & 每天 / 常常 / 有时候 / 很少 \\
\hline
Appareil utilisé & 手机 / 电脑 / 平板电脑 \\
\hline
Mes activités préférées (2 ou 3) & 查资料 / 聊天 / 发邮件 / 看视频 / 听音乐 / 玩游戏 \\
\hline
Mon opinion (2 adjectifs) & 方便 / 重要 / 有意思 / 有用 \\
\hline
Ma mise en garde & 不能花太多时间 / 要好好学习 \\
\hline
\end{tabularx}
\end{center}

\textbf{Tâche 3 — Rédiger ton portrait.} Rédige ton texte \zh{我和互联网} de 5 à 7 phrases en suivant ce plan : (1) phrase d'ouverture avec la fréquence ; (2) l'appareil ; (3) deux ou trois activités avec 喜欢 ; (4) ton opinion avec 我觉得 ; (5) une mise en garde avec 但是 ou 不过. Écris en caractères, avec le pinyin au-dessus si ton professeur le demande.

\textbf{Tâche 4 — Relecture guidée.} Vérifie ton texte avec la grille de contrôle : chaque phrase a-t-elle un sujet ? Les adverbes de fréquence sont-ils placés \textbf{avant} le verbe ? As-tu utilisé 在 + lieu + verbe pour « sur internet » ? As-tu bien écrit les caractères 网, 络, 便 ?

\textbf{Tâche 5 — Présentation orale en binôme.} Lis ton portrait à voix haute à ton binôme, lentement et en respectant les tons. Ton binôme écoute, puis te pose deux questions de compréhension, par exemple \zh{你常常上网吗？} ou \zh{你上网做什么？} Réponds en phrases complètes. Échangez ensuite les rôles.

\textbf{Tâche 6 — Le mur des portraits.} Recopie ton texte corrigé sur une fiche. La classe vote pour les trois portraits les plus clairs et les plus originaux : ils seront affichés au mur de la salle multimédia.

\courssub{Ressources à mobiliser}

\begin{aretenir}[Boîte à outils de la leçon]
\begin{itemize}
\item \cle{Vocabulaire} : 互联网 (hùliánwǎng, internet), 上网 (shàngwǎng, se connecter), 手机 (shǒujī, téléphone portable), 电脑 (diànnǎo, ordinateur), 网站 (wǎngzhàn, site web), 查资料 (chá zīliào, chercher des informations), 聊天 (liáotiān, bavarder en ligne), 发邮件 (fā yóujiàn, envoyer un courriel), 方便 (fāngbiàn, pratique), 重要 (zhòngyào, important).
\item \cle{Fréquence} : Sujet + 每天 / 常常 / 有时候 / 很少 + verbe. L'adverbe se place toujours \textbf{avant} le verbe : \zh{我每天上网。} (Wǒ měitiān shàngwǎng. « Je me connecte tous les jours. »)
\item \cle{Lieu de l'action} : Sujet + 在 + lieu + verbe : \zh{我在网上查资料。} (Wǒ zài wǎngshang chá zīliào. « Je cherche des informations sur internet. »)
\item \cle{Opinion} : 我觉得 + phrase ; adjectif seul comme attribut : \zh{互联网很方便。} (Hùliánwǎng hěn fāngbiàn. « Internet est très pratique. »)
\item \cle{Opposition} : 但是 / 不过 + phrase : \zh{但是，我们不能花太多时间上网。} (Dànshì, wǒmen bù néng huā tài duō shíjiān shàngwǎng. « Mais nous ne devons pas passer trop de temps en ligne. »)
\end{itemize}
\end{aretenir}

\begin{attention}[Pièges à éviter dans ton portrait]
\begin{itemize}
\item Ne dis pas \zh{我上网每天} : en chinois, l'adverbe de fréquence précède le verbe → \zh{我每天上网。}
\item Ne dis pas \zh{我喜欢网上查资料} : le lieu s'introduit avec 在 → \zh{我喜欢在网上查资料。}
\item 上网 est un verbe séparable : on dit \zh{上一次网} (« se connecter une fois »), pas \zh{上网上一次}.
\item Attention aux tons : 上网 est shàngwǎng (4\textsuperscript{e} + 3\textsuperscript{e} ton) ; 方便 est fāngbiàn (1\textsuperscript{er} + 4\textsuperscript{e} ton).
\item N'oublie pas le 很 de liaison devant l'adjectif attribut : \zh{互联网很重要}, et non \zh{互联网重要} (phrase correcte mais qui sonne comparative).
\end{itemize}
\end{attention}

\courssub{Proposition de production et corrigé commenté}

\begin{exoresolu}[Portrait numérique modèle : 我和互联网]
Rédige ton portrait numérique de 5 à 7 phrases en suivant le plan de la tâche 3, puis compare-le au modèle ci-dessous.
\tcblower
\textbf{Production modèle :}\\[4pt]
我和互联网\\
Wǒ hé hùliánwǎng\\
« Internet et moi »\\[4pt]
我每天用我的手机上网。\\
Wǒ měitiān yòng wǒ de shǒujī shàngwǎng.\\
« Je me connecte tous les jours avec mon téléphone portable. »\\[4pt]
我喜欢在网上查资料，也喜欢和朋友聊天。\\
Wǒ xǐhuan zài wǎngshang chá zīliào, yě xǐhuan hé péngyou liáotiān.\\
« J'aime chercher des informations sur internet, et j'aime aussi bavarder avec mes amis. »\\[4pt]
有时候，我给老师发邮件。\\
Yǒu shíhou, wǒ gěi lǎoshī fā yóujiàn.\\
« Parfois, j'envoie un courriel à mon professeur. »\\[4pt]
我觉得互联网很方便，也很重要。\\
Wǒ juéde hùliánwǎng hěn fāngbiàn, yě hěn zhòngyào.\\
« Je trouve qu'internet est très pratique, et aussi très important. »\\[4pt]
但是，我们不能花太多时间上网。\\
Dànshì, wǒmen bù néng huā tài duō shíjiān shàngwǎng.\\
« Mais nous ne devons pas passer trop de temps en ligne. »\\[4pt]
学习是最重要的！\\
Xuéxí shì zuì zhòngyào de !\\
« Les études sont le plus important ! »\\[6pt]
\textbf{Commentaire des choix de langue :}
\begin{itemize}
\item \textbf{Fréquence} : 每天 est placé juste après le sujet et avant le verbe 上网, conformément à la règle ; 有时候 ouvre la troisième phrase, position possible en tête de phrase pour varier le rythme.
\item \textbf{Appareil} : la structure 用 + objet + verbe (« utiliser… pour… ») introduit naturellement 手机 : 用我的手机上网.
\item \textbf{Activités} : 喜欢 + verbe permet d'enchaîner deux activités ; 也 (« aussi ») évite la répétition de 我喜欢 ; 在网上 respecte la structure 在 + lieu + verbe.
\item \textbf{Opinion} : 我觉得 + phrase complète, avec les deux adjectifs du thème 方便 et 重要, précédés du 很 de liaison.
\item \textbf{Mise en garde} : 但是 marque l'opposition ; 不能 + verbe exprime l'interdiction morale (« il ne faut pas ») ; 花太多时间 (« passer trop de temps ») est une collocation très utile du thème.
\item \textbf{Chute} : le superlatif 最 + adjectif + 的 (学习是最重要的) conclut le portrait sur une phrase courte et mémorable, adaptée à un affichage mural.
\end{itemize}
\end{exoresolu}

\courssub{Bilan de l'activité}

Dans cette activité d'intégration, tu as mobilisé l'ensemble des savoir-faire de la leçon : tu as \textbf{compris} une annonce authentique en chinois, \textbf{réutilisé} le vocabulaire d'internet dans un texte personnel, \textbf{appliqué} les structures de fréquence (常常, 每天, 有时候), de lieu (在网上) et d'opinion (我觉得…很方便), puis tu as \textbf{produit oralement} ton portrait devant un camarade en respectant les tons, et tu as \textbf{écouté} pour poser des questions de compréhension. Le mur des portraits de la classe constitue désormais une banque de modèles : relis les portraits affichés, repère les structures que tu n'avais pas utilisées, et garde ton texte dans ton portfolio — il te servira de base pour l'expression écrite du module « Mass médias et télécommunication » et pour l'évaluation de fin de séquence.

\newpage

\coursec{Méthodes et exercices résolus}

\courssub{Savoir-faire 1 : Lire et comprendre un court texte chinois sur internet}

\begin{methode}[Comprendre un texte sur internet en 5 étapes]
\begin{enumerate}
\item \textbf{Première lecture globale} : lisez le texte une première fois sans dictionnaire. Repérez les mots déjà connus (\zh{网络} wǎngluò « internet », \zh{手机} shǒujī « téléphone portable », \zh{信息} xìnxī « information ») pour deviner le thème général.
\item \textbf{Repérage des mots-clés du thème} : soulignez le vocabulaire d'internet : \zh{上网} shàngwǎng « se connecter », \zh{网站} wǎngzhàn « site web », \zh{发电子邮件} fā diànzǐ yóujiàn « envoyer un courriel », \zh{聊天} liáotiān « bavarder en ligne ».
\item \textbf{Lecture analytique} : phrase par phrase, identifiez le sujet, le verbe, le complément. En chinois, l'ordre de base est Sujet + Verbe + Complément d'objet ; les compléments de temps, de lieu et d'outil se placent avant le verbe : \zh{我每天上网。} Wǒ měitiān shàngwǎng. « Je me connecte chaque jour. »
\item \textbf{Reformulation} : résumez chaque phrase en français simple pour vérifier que vous avez compris (qui ? fait quoi ? quand ? où ? pourquoi ?).
\item \textbf{Vérification} : relisez le texte en entier et contrôlez que votre résumé correspond à toutes les informations données.
\end{enumerate}
Réflexe d'examen : les réponses aux questions de compréhension se trouvent presque toujours dans le texte, dans le même ordre que les questions. Cherchez la phrase qui contient le mot-clé de la question, puis recopiez ou adaptez la réponse en chinois correct.
\end{methode}

\begin{textebox}[小王的网络生活 — texte support]
小王是雅温得一所中学的学生。他很喜欢上网。每天晚上，他用手机上网，看新闻，听音乐，还和朋友聊天。周末，他常常在网上学习汉语。他觉得网络很有用，但是不能花太多时间在网上。\\
Xiǎo Wáng shì Yǎwēndé yì suǒ zhōngxué de xuésheng. Tā hěn xǐhuan shàngwǎng. Měitiān wǎnshang, tā yòng shǒujī shàngwǎng, kàn xīnwén, tīng yīnyuè, hái hé péngyou liáotiān. Zhōumò, tā chángcháng zài wǎngshang xuéxí Hànyǔ. Tā juéde wǎngluò hěn yǒuyòng, dànshì bù néng huā tài duō shíjiān zài wǎngshang.\\
« Petit Wang est élève d'un lycée de Yaoundé. Il aime beaucoup se connecter à internet. Chaque soir, il se connecte avec son téléphone portable, lit les informations, écoute de la musique et bavarde aussi avec ses amis. Le week-end, il apprend souvent le chinois en ligne. Il trouve qu'internet est très utile, mais qu'on ne doit pas passer trop de temps en ligne. »
\end{textebox}

\begin{exoresolu}[Compréhension de texte : questions sur 小王的网络生活]
\textbf{Énoncé.} Relisez le texte support ci-dessus, puis répondez aux questions en chinois, par des phrases complètes.
\begin{enumerate}
\item 小王是哪里的学生？Xiǎo Wáng shì nǎlǐ de xuésheng ?
\item 他用什么上网？Tā yòng shénme shàngwǎng ?
\item 周末他常常在网上做什么？Zhōumò tā chángcháng zài wǎngshang zuò shénme ?
\item 他觉得网络怎么样？Tā juéde wǎngluò zěnmeyàng ?
\end{enumerate}
\tcblower
\textbf{Solution détaillée.}
\begin{enumerate}
\item La question porte sur le lieu (« où » : \zh{哪里}). La première phrase donne la réponse : \zh{雅温得一所中学}. Réponse rédigée : \zh{小王是雅温得一所中学的学生。} Xiǎo Wáng shì Yǎwēndé yì suǒ zhōngxué de xuésheng. « Petit Wang est élève d'un lycée de Yaoundé. » (On reprend la structure complète Sujet + 是 + complément.)
\item Le mot-clé est \zh{用什么} « avec quoi ». Le texte dit \zh{他用手机上网}. Réponse : \zh{他用手机上网。} Tā yòng shǒujī shàngwǎng. « Il se connecte avec son téléphone portable. » La structure est Sujet + 用 + outil + verbe.
\item Le mot-clé est \zh{周末} « le week-end ». Le texte dit \zh{周末，他常常在网上学习汉语}. Réponse : \zh{周末他常常在网上学习汉语。} Zhōumò tā chángcháng zài wǎngshang xuéxí Hànyǔ. « Le week-end, il apprend souvent le chinois en ligne. »
\item La question \zh{怎么样} demande une opinion. Le texte dit \zh{他觉得网络很有用，但是不能花太多时间在网上}. Réponse : \zh{他觉得网络很有用，但是不能花太多时间在网上。} Tā juéde wǎngluò hěn yǒuyòng, dànshì bù néng huā tài duō shíjiān zài wǎngshang. « Il trouve qu'internet est très utile, mais qu'on ne doit pas y passer trop de temps. »
\end{enumerate}
\textbf{Conclusion méthodologique} : chaque réponse reprend un élément exact du texte, inséré dans une phrase complète avec sujet et verbe. Une réponse d'un seul mot (\zh{手机}) est juste mais rapporte moins de points qu'une phrase complète.
\end{exoresolu}

\courssub{Savoir-faire 2 : Reconnaître et prononcer le vocabulaire du thème}

\textbf{Glossaire du thème : internet}\par

\begin{center}
\begin{tabularx}{\linewidth}{|X|X|X|X|}
\hline
\rowcolor{popblueL} Caractères & Pinyin & Nature & Traduction \\
\hline
网络 & wǎngluò & nom & internet, réseau \\
\hline
上网 & shàngwǎng & verbe & se connecter à internet \\
\hline
网上 & wǎngshang & nom & sur internet, en ligne \\
\hline
网站 & wǎngzhàn & nom & site web \\
\hline
网友 & wǎngyǒu & nom & ami(e) connu(e) sur internet \\
\hline
电脑 & diànnǎo & nom & ordinateur \\
\hline
手机 & shǒujī & nom & téléphone portable \\
\hline
电子邮件 & diànzǐ yóujiàn & nom & courriel, e-mail \\
\hline
发 & fā & verbe & envoyer, émettre \\
\hline
收 & shōu & verbe & recevoir \\
\hline
下载 & xiàzài & verbe & télécharger \\
\hline
搜索 & sōusuǒ & verbe & rechercher (sur internet) \\
\hline
密码 & mìmǎ & nom & mot de passe \\
\hline
聊天 & liáotiān & verbe & bavarder, discuter (en ligne) \\
\hline
新闻 & xīnwén & nom & informations, actualités \\
\hline
信息 & xìnxī & nom & information, message \\
\hline
资料 & zīliào & nom & documents, données \\
\hline
软件 & ruǎnjiàn & nom & logiciel \\
\hline
有用 & yǒuyòng & adjectif & utile \\
\hline
方便 & fāngbiàn & adjectif & pratique, commode \\
\hline
花 & huā & verbe & dépenser, passer (du temps) \\
\hline
\end{tabularx}
\end{center}

\begin{methode}[Mémoriser le vocabulaire d'internet]
\begin{enumerate}
\item \textbf{Classer par familles} : regroupez les mots par sens : le matériel (\zh{电脑} diànnǎo « ordinateur », \zh{手机} shǒujī « portable »), les actions (\zh{上网} shàngwǎng « se connecter », \zh{发} fā « envoyer », \zh{下载} xiàzài « télécharger », \zh{搜索} sōusuǒ « rechercher »), les outils (\zh{网站} wǎngzhàn « site », \zh{电子邮件} diànzǐ yóujiàn « courriel », \zh{密码} mìmǎ « mot de passe »).
\item \textbf{Lire à voix haute avec les tons} : prononcez chaque mot trois fois en marquant les tons : \zh{wǎngluò} (3\textsuperscript{e} + 4\textsuperscript{e} ton), \zh{shàngwǎng} (4\textsuperscript{e} + 3\textsuperscript{e} ton).
\item \textbf{Utiliser les caractères comme indices} : le caractère \zh{网} wǎng « filet/réseau » revient dans \zh{网络}, \zh{网站}, \zh{上网}, \zh{网友} wǎngyǒu « ami en ligne ». Le caractère \zh{电} diàn « électricité » revient dans \zh{电脑}, \zh{电话}, \zh{电子邮件}.
\item \textbf{Tester par traduction croisée} : faites-vous dicter la liste français vers chinois, puis chinois vers français.
\item \textbf{Employer chaque mot dans une phrase} : un mot n'est vraiment acquis que s'il est utilisé dans un contexte : \zh{我用电脑发电子邮件。} Wǒ yòng diànnǎo fā diànzǐ yóujiàn. « J'envoie des courriels avec l'ordinateur. »
\end{enumerate}
\end{methode}

\begin{exoresolu}[Vocabulaire : retrouver les mots du thème]
\textbf{Énoncé.} Complétez les phrases avec le mot qui convient : \zh{上网}、\zh{网站}、\zh{密码}、\zh{下载}、\zh{电子邮件}。
\begin{enumerate}
\item 我每天晚上用手机\_\_\_\_\_\_，看新闻。
\item 这个\_\_\_\_\_\_有很多汉语学习资料。
\item 我想给你发一个\_\_\_\_\_\_。
\item 请不要告诉别人你的\_\_\_\_\_\_。
\item 你可以在网上\_\_\_\_\_\_这首歌。
\end{enumerate}
\tcblower
\textbf{Solution détaillée.}
\begin{enumerate}
\item \zh{我每天晚上用手机上网，看新闻。} Wǒ měitiān wǎnshang yòng shǒujī shàngwǎng, kàn xīnwén. « Chaque soir, je me connecte avec mon portable et je lis les informations. » Le contexte « téléphone + regarder les infos » exige le verbe \zh{上网} « se connecter ».
\item \zh{这个网站有很多汉语学习资料。} Zhège wǎngzhàn yǒu hěn duō Hànyǔ xuéxí zīliào. « Ce site web contient beaucoup de documents pour apprendre le chinois. » Le classificateur \zh{个} et le contenu « documents » désignent un \zh{网站} « site ».
\item \zh{我想给你发一个电子邮件。} Wǒ xiǎng gěi nǐ fā yí ge diànzǐ yóujiàn. « Je veux t'envoyer un courriel. » Le verbe \zh{发} « envoyer » appelle le complément \zh{电子邮件}.
\item \zh{请不要告诉别人你的密码。} Qǐng bú yào gàosu biérén nǐ de mìmǎ. « S'il te plaît, ne dis pas ton mot de passe aux autres. » Ce qu'on ne doit pas révéler, c'est le \zh{密码} « mot de passe ».
\item \zh{你可以在网上下载这首歌。} Nǐ kěyǐ zài wǎngshang xiàzài zhè shǒu gē. « Tu peux télécharger cette chanson sur internet. » Attention au classificateur des chansons : \zh{首} shǒu, pas \zh{个}.
\end{enumerate}
\textbf{Conclusion} : pour choisir le bon mot, repérez le verbe ou le classificateur de la phrase : ce sont eux qui imposent le complément.
\end{exoresolu}

\courssub{Savoir-faire 3 : Employer les structures avec 用, 在 et 花}

\begin{propriete}[Les trois structures clés pour parler d'internet]
\begin{enumerate}
\item \textbf{L'outil} : Sujet + \zh{用} + outil + verbe (+ objet). \zh{我用电脑学习。} Wǒ yòng diànnǎo xuéxí. « J'étudie avec l'ordinateur. » Le complément d'outil se place AVANT le verbe, contrairement au français.
\item \textbf{Le lieu (réel ou virtuel)} : Sujet + \zh{在} + lieu + verbe (+ objet). \zh{我在网上买东西。} Wǒ zài wǎngshang mǎi dōngxi. « J'achète des choses sur internet. » Le complément de lieu précède aussi le verbe.
\item \textbf{Le temps passé} : Sujet + \zh{花} + durée + verbe. \zh{他每天花两个小时上网。} Tā měitiān huā liǎng ge xiǎoshí shàngwǎng. « Il passe deux heures par jour sur internet. »
\item \textbf{La comparaison} (révision) : A + \zh{比} + B + adjectif. \zh{手机比电脑方便。} Shǒujī bǐ diànnǎo fāngbiàn. « Le portable est plus pratique que l'ordinateur. »
\end{enumerate}
Réflexe : en chinois, tout ce qui encadre l'action (temps, lieu, outil, manière) se place AVANT le verbe ; seul le complément d'objet suit.
\end{propriete}

\begin{center}
\begin{tabularx}{\linewidth}{|X|X|X|}
\hline
\rowcolor{popblueL} Structure & Exemple (caractères + pinyin) & Traduction \\
\hline
Sujet + 用 + outil + verbe & 我用手机上网。Wǒ yòng shǒujī shàngwǎng. & Je me connecte avec mon portable. \\
\hline
Sujet + 在 + lieu + verbe & 她在网上学习汉语。Tā zài wǎngshang xuéxí Hànyǔ. & Elle apprend le chinois en ligne. \\
\hline
Sujet + 花 + durée + verbe & 我们每天花一个小时聊天。Wǒmen měitiān huā yí ge xiǎoshí liáotiān. & Nous passons une heure par jour à discuter. \\
\hline
A + 比 + B + adjectif & 电脑比手机快。Diànnǎo bǐ shǒujī kuài. & L'ordinateur est plus rapide que le téléphone. \\
\hline
\end{tabularx}
\end{center}

\begin{exoresolu}[Grammaire : traduction thème]
\textbf{Énoncé.} Traduisez en chinois :
\begin{enumerate}
\item J'apprends le chinois sur internet avec mon téléphone portable.
\item Ma grande sœur passe trois heures par jour à bavarder en ligne.
\item L'ordinateur est plus rapide que le téléphone.
\end{enumerate}
\tcblower
\textbf{Solution détaillée.}
\begin{enumerate}
\item Analyse : sujet « je » \zh{我} ; lieu « sur internet » \zh{在网上} ; outil « avec mon portable » \zh{用手机} ; verbe « apprendre » \zh{学习} ; objet « le chinois » \zh{汉语}. Ordre chinois : Sujet + lieu + outil + verbe + objet. Traduction : \zh{我在网上用手机学习汉语。} Wǒ zài wǎngshang yòng shǒujī xuéxí Hànyǔ. (On peut aussi inverser lieu et outil : \zh{我用手机在网上学习汉语。} Les deux ordres sont corrects.)
\item Analyse : « passer du temps à faire » = \zh{花} + durée + verbe ; « par jour » = \zh{每天} placé avant le verbe ; « bavarder en ligne » = \zh{聊天} (ou \zh{在网上聊天}). Traduction : \zh{我姐姐每天花三个小时聊天。} Wǒ jiějie měitiān huā sān ge xiǎoshí liáotiān. « Ma grande sœur passe trois heures par jour à bavarder en ligne. »
\item Structure comparative : A + \zh{比} + B + adjectif. Traduction : \zh{电脑比手机快。} Diànnǎo bǐ shǒujī kuài. « L'ordinateur est plus rapide que le téléphone. » Ne traduisez jamais « plus... que » par \zh{更...比} : \zh{比} suffit à lui seul.
\end{enumerate}
\textbf{Conclusion} : avant de traduire, réécrivez mentalement la phrase française dans l'ordre chinois : Sujet + temps + lieu + outil + verbe + objet.
\end{exoresolu}

\courssub{Savoir-faire 4 : Écrire les caractères du thème (clés et ordre des traits)}

\begin{methode}[Analyser et tracer les caractères d'internet]
\begin{enumerate}
\item \textbf{Identifier la clé (radical)} : elle donne le sens général. Exemples : \zh{网} figure un filet ; \zh{话} huà « parole » a la clé de la parole \zh{讠} ; \zh{信} xìn « lettre/message » a la clé de l'homme \zh{亻} + \zh{言} « parole » (un homme et sa parole : la confiance, le message).
\item \textbf{Respecter l'ordre des traits} : de haut en bas, de gauche à droite, l'horizontal avant le vertical, l'extérieur avant l'intérieur, on ferme la boîte en dernier.
\item \textbf{Décomposer avant d'écrire} : \zh{聊} liáo = \zh{耳} « oreille » à gauche + \zh{卯} à droite ; \zh{电} diàn = le champ \zh{田} traversé par un crochet vertical qui dépasse en bas.
\item \textbf{Écrire en répétant le pinyin} : à chaque caractère tracé, prononcez le mot avec ses tons pour lier forme, son et sens.
\end{enumerate}
\end{methode}

\begin{exoresolu}[Écriture : clés et composition]
\textbf{Énoncé.} 1) Donnez la clé (radical) de chaque caractère : \zh{信}、\zh{话}、\zh{网}。2) Décrivez en français l'ordre des traits de \zh{电}。3) Recopiez cinq fois le mot \zh{网络} en respectant l'ordre des traits.
\tcblower
\textbf{Solution détaillée.}
\begin{enumerate}
\item \zh{信} xìn : clé de l'homme \zh{亻} (à gauche) ; la partie droite \zh{言} signifie « parole ». \zh{话} huà : clé de la parole simplifiée \zh{讠} (à gauche), avec \zh{舌} « langue » à droite. \zh{网} wǎng : l'enceinte \zh{冂} contient deux croisements qui figurent les mailles du filet ; c'est un caractère pictographique.
\item \zh{电} diàn (5 traits) : on trace d'abord le cadre du champ : trait vertical gauche, puis horizontal-crochet formant le haut et le côté droit ; ensuite l'horizontal du milieu, puis le vertical du milieu qui se prolonge en bas en crochet \zh{乚} ; enfin l'horizontal du bas qui ferme. Règles appliquées : l'extérieur avant l'intérieur, on ferme en dernier ; le trait final (le crochet) dépasse du cadre, ce qui distingue \zh{电} de \zh{田}.
\item \zh{网络} wǎngluò : \zh{网} se trace de l'extérieur vers l'intérieur (l'enceinte d'abord, les deux croisements ensuite) ; \zh{络} se trace de gauche à droite : la clé de la soie \zh{纟} d'abord (trois traits), puis la partie droite \zh{各} de haut en bas. Répétez cinq fois en prononçant « wǎngluò » à chaque écriture.
\end{enumerate}
\textbf{Conclusion} : connaître la clé permet de deviner le sens d'un caractère inconnu et de le retrouver dans un dictionnaire ; respecter l'ordre des traits rend l'écriture plus rapide et plus lisible.
\end{exoresolu}

\courssub{Savoir-faire 5 : Produire des phrases avec les collocations du thème}

\begin{methode}[Construire un court paragraphe sur internet]
\begin{enumerate}
\item \textbf{Mémoriser les collocations} (verbe + complément figé) : \zh{上网} « se connecter », \zh{发电子邮件} « envoyer un courriel », \zh{看新闻} « lire les infos », \zh{听音乐} « écouter de la musique », \zh{下载软件} « télécharger un logiciel », \zh{搜索资料} « rechercher des documents », \zh{和朋友聊天} « bavarder avec des amis ».
\item \textbf{Plan en trois phrases} : (a) qui et avec quel outil ; (b) quelles activités (2 ou 3 collocations reliées par \zh{还} « et aussi ») ; (c) une opinion avec \zh{觉得} ou \zh{认为}.
\item \textbf{Relier} avec \zh{每天}、\zh{常常}、\zh{但是} pour un texte fluide.
\item \textbf{Se relire} : vérifier les tons du pinyin, la place des compléments avant le verbe, et l'accord verbe-complément des collocations.
\end{enumerate}
\end{methode}

\begin{exoresolu}[Production guidée : mon usage d'internet]
\textbf{Énoncé.} En vous aidant des collocations apprises, écrivez un court paragraphe (4 à 5 phrases) sur votre usage d'internet. Utilisez au moins : \zh{用}、\zh{上网}、\zh{发电子邮件}、\zh{觉得}。
\tcblower
\textbf{Solution détaillée (modèle de production).}

\zh{我每天都上网。我用手机看新闻，听音乐，还和同学聊天。周末，我常常用电脑发电子邮件，也在网上学习汉语。有时候我下载学习资料。我觉得网络很方便，但是我们不能花太多时间在网上。}

Wǒ měitiān dōu shàngwǎng. Wǒ yòng shǒujī kàn xīnwén, tīng yīnyuè, hái hé tóngxué liáotiān. Zhōumò, wǒ chángcháng yòng diànnǎo fā diànzǐ yóujiàn, yě zài wǎngshang xuéxí Hànyǔ. Yǒu shíhou wǒ xiàzài xuéxí zīliào. Wǒ juéde wǎngluò hěn fāngbiàn, dànshì wǒmen bù néng huā tài duō shíjiān zài wǎngshang.

« Je me connecte tous les jours. Avec mon portable, je lis les infos, j'écoute de la musique et je bavarde aussi avec mes camarades. Le week-end, j'envoie souvent des courriels avec l'ordinateur et j'apprends aussi le chinois en ligne. Parfois, je télécharge des documents d'étude. Je trouve qu'internet est très pratique, mais nous ne devons pas y passer trop de temps. »

\textbf{Justifications grammaticales} : \zh{每天} et \zh{周末} (compléments de temps) précèdent le verbe ; \zh{用手机} et \zh{用电脑} (outils) précèdent le verbe ; \zh{在网上} (lieu) précède \zh{学习} ; \zh{还} et \zh{也} se placent avant le verbe ; l'opinion finale utilise \zh{觉得} + adjectif \zh{方便} et la nuance \zh{但是} + \zh{不能}.
\end{exoresolu}

\begin{attention}[Les 5 erreurs qui coûtent des points à l'examen]
\begin{enumerate}
\item \textbf{Calquer l'ordre français} : écrire \zh{我学习在网上} est faux. Le complément de lieu se place AVANT le verbe : \zh{我在网上学习。} Wǒ zài wǎngshang xuéxí. De même pour l'outil : \zh{我用手机上网。} et non \zh{我上网用手机}.
\item \textbf{Confondre les caractères proches} : \zh{网} wǎng « réseau » n'est pas \zh{门} mén « porte » ; \zh{电} diàn « électricité » n'est pas \zh{田} tián « champ » (le trait final dépasse !) ; \zh{信} xìn « message » n'est pas \zh{言} yán « parole ». Un trait oublié change le mot et fait perdre le point.
\item \textbf{Négliger les tons du pinyin} : \zh{wǎng} (3\textsuperscript{e} ton, « réseau ») et \zh{wàng} (4\textsuperscript{e} ton, « oublier ») n'ont rien à voir ; \zh{mǎi} « acheter » et \zh{mài} « vendre » ne se distinguent que par le ton. Écrire le pinyin sans tons est sanctionné.
\item \textbf{Mal employer les mots-outils} : \zh{了} ne marque pas le passé en général mais l'accomplissement ; ne dites pas \zh{我每天上网了}. Pour une habitude, pas de \zh{了} : \zh{我每天上网。} De même, \zh{的} est obligatoire entre le possesseur et le nom : \zh{我的密码} wǒ de mìmǎ « mon mot de passe ».
\item \textbf{Oublier ou mal choisir le classificateur} : une chanson se compte avec \zh{首} : \zh{一首歌} yì shǒu gē ; un courriel avec \zh{个} ou \zh{封} : \zh{一个电子邮件} / \zh{一封电子邮件} ; un ordinateur avec \zh{台} : \zh{一台电脑} yì tái diànnǎo. Le classificateur \zh{个} n'est pas universel !
\end{enumerate}
\end{attention}

\begin{exercice}[Pour s'entraîner : internet et moi]
\begin{enumerate}
\item Traduisez en français : \zh{我哥哥每天晚上花两个小时上网。} Wǒ gēge měitiān wǎnshang huā liǎng ge xiǎoshí shàngwǎng.
\item Traduisez en chinois : « Mon ami télécharge souvent de la musique sur internet. »
\item Répondez en chinois par une phrase complète : 你常常在网上做什么？Nǐ chángcháng zài wǎngshang zuò shénme ?
\item Donnez la clé (radical) de \zh{聊} et de \zh{站}, puis recopiez trois fois le mot \zh{网站}.
\end{enumerate}
\end{exercice}

\begin{corrige}[Exercice « Pour s'entraîner »]
\begin{enumerate}
\item « Mon grand frère passe deux heures sur internet chaque soir. » Structure : Sujet + temps (\zh{每天晚上}) + \zh{花} + durée + verbe.
\item \zh{我的朋友常常在网上下载音乐。} Wǒ de péngyou chángcháng zài wǎngshang xiàzài yīnyuè. Le complément de lieu \zh{在网上} précède le verbe \zh{下载}.
\item Réponse possible : \zh{我常常在网上看新闻，听音乐。} Wǒ chángcháng zài wǎngshang kàn xīnwén, tīng yīnyuè. « Je lis souvent les infos et j'écoute de la musique en ligne. » Toute phrase complète avec \zh{在网上} + verbe est acceptée.
\item \zh{聊} liáo : clé de l'oreille \zh{耳} (à gauche). \zh{站} zhàn : clé \zh{立} « se tenir debout » (à gauche). Recopiez \zh{网站} trois fois : \zh{网} de l'extérieur vers l'intérieur ; \zh{站} de gauche à droite (\zh{立} puis \zh{占}).
\end{enumerate}
\end{corrige}

\newpage

\coursec{Exercices}

\coursec{Leçon 35 — Vocabulaire et compréhension de texte : Internet}

\courssub{Mise en route et découverte du thème}

\begin{experience}[Brainstorming : Internet et moi]
En binôme, discutez en chinois (ou en français si nécessaire) pendant 3 minutes. Notez 5 mots chinois que vous associez à \zh{网络} (Internet).
\begin{itemize}
    \item Outils : \zh{手机、电脑、平板电脑、网线}...
    \item Actions : \zh{上网、聊天、看视频、发邮件、查资料、下载、网购}...
    \item Applications : \zh{微信、抖音、QQ、浏览器}...
\end{itemize}
Mise en commun au tableau : l'enseignant classe les mots par catégorie (Noms d'outils, Verbes d'action, Noms de plateformes) et introduit le vocabulaire de la leçon.
\end{experience}

\begin{textebox}[Dialogue d'introduction : \zh{周末上网}]
\zh{—— 小王，你周末通常做什么？} \\
\emph{—— Xiǎo Wáng, nǐ zhōumò tōngcháng zuò shénme?} \\
« Petit Wang, que fais-tu d'habitude le week-end ? » \\

\zh{—— 我通常在家上网。我用电脑查资料，也用手机看视频、和朋友聊天。} \\
\emph{—— Wǒ tōngcháng zài jiā shàngwǎng. Wǒ yòng diànnǎo chá zīliào, yě yòng shǒujī kàn shìpín, hé péngyou liáotiān.} \\
« D'habitude je suis à la maison sur Internet. J'utilise l'ordinateur pour chercher des infos, et le téléphone pour regarder des vidéos et chatter avec des amis. » \\

\zh{—— 网络真方便！但是要注意保护眼睛，别上太久。} \\
\emph{—— Wǎngluò zhēn fāngbiàn! Dànshì yào zhùyì bǎohù yǎnjing, bié shàng tài jiǔ.} \\
« Internet est vraiment pratique ! Mais fais attention à protéger tes yeux, ne reste pas trop longtemps en ligne. »
\end{textebox}

\courssub{Texte support : Lecture et écoute}

\begin{textebox}[Texte original : \zh{网络生活} (La vie sur Internet)]
\zh{现在，网络已经成为我们生活的一部分。在喀麦隆，无论是在雅温得还是杜阿拉，越来越多的中学生有自己的手机，他们每天用手机上网。} \\
\emph{Xiànzài, wǎngluò yǐjīng chéngwéi wǒmen shēnghuó de yībùfèn. Zài Kāmàilóng, wúlùn shì zài Yǎwēndé háishì Dù'ālā, yuè lái yuè duō de zhōngxuéshēng yǒu zìjǐ de shǒujī, tāmen měitiān yòng shǒujī shàngwǎng.} \\
« Aujourd'hui, Internet est déjà devenu une partie de notre vie. Au Cameroun, que ce soit à Yaoundé ou à Douala, de plus en plus de lycéens ont leur propre téléphone portable, ils utilisent leur téléphone pour aller sur Internet chaque jour. » \\

\zh{小李是雅温得一所高中的学生。他每天放学后，先回家吃饭，然后用电脑查资料做作业。周末的时候，他喜欢用微信和朋友聊天，看中国的短视频，有时候也下载中文歌曲听。} \\
\emph{Xiǎo Lǐ shì Yǎwēndé yī suǒ gāozhōng de xuéshēng. Tā měitiān fàngxué hòu, xiān huí jiā chīfàn, ránhòu yòng diànnǎo chá zīliào zuò zuòyè. Zhōumò de shíhou, tā xǐhuan yòng Wēixìn hé péngyou liáotiān, kàn Zhōngguó de duǎn shìpín, yǒushíhou yě xiàzǎi Zhōngwén gēqǔ tīng.} \\
« Petit Li est un élève d'un lycée de Yaoundé. Chaque jour après la classe, il rentre d'abord manger, puis utilise l'ordinateur pour chercher des infos et faire ses devoirs. Le week-end, il aime utiliser WeChat pour chatter avec des amis, regarder des vidéos courtes chinoises, parfois aussi télécharger des chansons chinoises pour les écouter. » \\

\zh{他说：「网络很方便，让我知道中国的文化。但是上网时间太长，对眼睛不好，所以我每天只上网两个小时。」} \\
\emph{Tā shuō: « Wǎngluò hěn fāngbiàn, ràng wǒ zhīdào Zhōngguó de wénhuà. Dànshì shàngwǎng shíjiān tài cháng, duì yǎnjing bù hǎo, suǒyǐ wǒ měitiān zhǐ shàngwǎng liǎng gè xiǎoshí. »} \\
« Il dit : "Internet est très pratique, il me permet de connaître la culture chinoise. Mais le temps passé en ligne trop long est mauvais pour les yeux, donc je ne vais sur Internet que deux heures par jour." » \\

\zh{在中国，网络发展得更快。人们不用现金，用手机扫码就能买东西、坐地铁、付账单。网购、网课、外卖，已经是中国人普通的生活方式。} \\
\emph{Zài Zhōngguó, wǎngluò fāzhǎn de gèng kuài. Rénmen bù yòng xiànjīn, yòng shǒujī sǎomǎ jiù néng mǎi dōngxi, zuò dìtiě, fù zhàngdān. Wǎnggòu, wǎngkè, wàimài, yǐjīng shì Zhōngguórén pǔtōng de shēnghuó fāngshì.} \\
« En Chine, Internet se développe encore plus vite. Les gens n'utilisent pas d'argent liquide, avec le téléphone ils scannent un code pour acheter des choses, prendre le métro, payer les factures. Achats en ligne, cours en ligne, livraison de repas, sont déjà le mode de vie ordinaire des Chinois. » \\

\zh{小李想：「以后我要去中国留学，现在多练习中文打字，将来一定很有用。」} \\
\emph{Xiǎo Lǐ xiǎng: « Yǐhòu wǒ yào qù Zhōngguó liúxué, xiànzài duō liànxí Zhōngwén dǎzì, jiānglái yídìng hěn yǒuyòng. »} \\
« Petit Li pense : "Plus tard je veux aller en Chine pour étudier, maintenant je m'entraîne plus à la frappe chinoise, à l'avenir ce sera sûrement très utile." »
\end{textebox}

\courssub{Glossaire du thème : Vocabulaire de base (HSK 2--3)}

\begin{aretenir}[Vocabulaire thématique : Internet et télécommunications]
\begin{center}
\begin{tabularx}{\linewidth}{|X|X|X|X|}
\hline
\rowcolor{popblueL} \textbf{Caractères} & \textbf{Pinyin} & \textbf{Nature} & \textbf{Traduction} \\
\hline
\zh{网络} & wǎngluò & n. & Internet, réseau, réseau social \\
\hline
\zh{上网} & shàngwǎng & VO (v. sep.) & Aller sur Internet, se connecter \\
\hline
\zh{网站} & wǎngzhàn & n. & Site web \\
\hline
\zh{电脑} & diànnǎo & n. & Ordinateur \\
\hline
\zh{手机} & shǒujī & n. & Téléphone portable, smartphone \\
\hline
\zh{平板电脑} & píngbǎn diànnǎo & n. & Tablette tactile \\
\hline
\zh{登录} & dēnglù & v. & Se connecter (sur un compte), se logger \\
\hline
\zh{注册} & zhùcè & v. & S'inscrire, créer un compte \\
\hline
\zh{密码} & mìmǎ & n. & Mot de passe \\
\hline
\zh{账号} & zhànghào & n. & Compte (utilisateur) \\
\hline
\zh{邮件} & yóujiàn & n. & Courriel, e-mail \\
\hline
\zh{发邮件} & fā yóujiàn & VO & Envoyer un e-mail \\
\hline
\zh{查资料} & chá zīliào & VO & Faire des recherches, chercher des informations \\
\hline
\zh{聊天} & liáotiān & VO (v. sep.) & Discuter, chatter (en ligne) \\
\hline
\zh{看视频} & kàn shìpín & VO & Regarder une vidéo \\
\hline
\zh{下载} & xiàzǎi & v. & Télécharger \\
\hline
\zh{上传} & shàngchuán & v. & Télécharger (vers un serveur), uploader \\
\hline
\zh{微信} & Wēixìn & n. & WeChat (application chinoise omniprésente) \\
\hline
\zh{视频} & shìpín & n. & Vidéo \\
\hline
\zh{方便} & fāngbiàn & adj. & Pratique, commode \\
\hline
\zh{扫码} & sǎomǎ & VO & Scanner un code QR \\
\hline
\zh{网购} & wǎnggòu & VO & Faire des achats en ligne \\
\hline
\zh{外卖} & wàimài & n./VO & Plat à emporter / Commander à manger \\
\hline
\zh{保护} & bǎohù & v. & Protéger \\
\hline
\zh{隐私} & yǐnsī & n. & Vie privée, données privées \\
\hline
\zh{眼睛} & yǎnjing & n. & Yeux \\
\hline
\end{tabularx}
\end{center}
\end{aretenir}

\courssub{Points de grammaire clés}

\begin{propriete}[Structure de l'instrument : \zh{用} (yòng)]
Pour exprimer \textbf{l'outil ou le moyen} utilisé pour accomplir une action, on place la préposition \zh{用} (ou \zh{用着}) \textbf{avant le verbe}, selon la structure :
\begin{center}
\textbf{Sujet + \zh{用} + Outil/Moyen + Verbe (+ Objet)}
\end{center}
\end{propriete}

\begin{exemplebox}[Exemples]
\zh{我用电脑查资料。} \emph{(Wǒ yòng diànnǎo chá zīliào.)} --- J'utilise l'ordinateur pour faire des recherches. \\
\zh{他用手机发邮件。} \emph{(Tā yòng shǒujī fā yóujiàn.)} --- Il utilise son téléphone pour envoyer un mail. \\
\zh{请用中文写。} \emph{(Qǐng yòng Zhōngwén xiě.)} --- Écrivez en chinois (utilisez le chinois pour écrire).
\end{exemplebox}
\begin{attention}[Piège fréquent : Ordre des mots]
\emph{Erreur :} \zh{*我查资料用电脑。} / \zh{*我电脑用查资料。} \\
\emph{Règle :} En chinois, le complément d'instrument (\zh{用} + Outil) se place \textbf{avant} le verbe, jamais après. C'est une différence majeure avec le français (« Je cherche des infos \textbf{avec l'ordinateur} »).
\end{attention}


\begin{propriete}[Les verbes séparables (离合词 \zh{líhécí}) du thème]
Beaucoup de verbes chinois sont formés de \textbf{Verbe + Objet} (structure VO). Ils fonctionnent comme un seul verbe mais peuvent être séparés : l'objet peut suivre le verbe, ou être déplacé (ex: après un indicateur d'aspect, un complément de durée, ou pour l'objet topicalisé).
\begin{center}
\begin{tabularx}{\linewidth}{|X|X|X|}
\hline
\rowcolor{popblueL} \textbf{Verbe séparable (VO)} & \textbf{Sens} & \textbf{Exemple de séparation} \\
\hline
\zh{上网} (shàng wǎng) & Aller sur Internet & \zh{上了两个小时网} (été en ligne 2h) \\
\hline
\zh{聊天} (liáo tiān) & Chatter & \zh{聊了半天天} (bavardé toute la matinée) \\
\hline
\zh{发邮件} (fā yóujiàn) & Envoyer un mail & \zh{发了一封邮件} (envoyé un mail) \\
\hline
\zh{查资料} (chá zīliào) & Chercher des infos & \zh{查了一会儿资料} (cherché un moment) \\
\hline
\zh{下载} (xià zǎi) & Télécharger & \zh{下载了一个软件} (téléchargé un logiciel) \\
\hline
\end{tabularx}
\end{center}
\end{propriete}

\begin{attention}[Attention : \zh{上网} vs \zh{网络}]
\zh{网络} (wǎngluò) = Nom « Internet / Le réseau ». \zh{上网} (shàngwǎng) = Verbe « Aller sur Internet / Se connecter ».
\emph{Erreur :} \zh{*我网络。} \rightarrow \emph{Correct :} \zh{我上网。} / \zh{我用网络。}
\end{attention}


\begin{propriete}[Position des adverbes de fréquence (\zh{每天、经常、有时、从不})]
Les adverbes de fréquence se placent \textbf{avant le verbe} (ou avant \zh{用} s'il y a un complément d'instrument), généralement après le sujet.
\begin{center}
\textbf{Sujet + Adverbe de fréquence + (\zh{用} Outil) + Verbe}
\end{center}
\end{propriete}

\begin{exemplebox}[Exemples]
\zh{我\cle{每天}用手机上网。} \emph{(Wǒ \cle{měitiān} yòng shǒujī shàngwǎng.)} --- Je vais sur Internet avec mon téléphone \textbf{tous les jours}. \\
\zh{他\cle{经常}查资料。} \emph{(Tā \cle{jīngcháng} chá zīliào.)} --- Il fait des recherches \textbf{souvent}. \\
\zh{我们\cle{有时}看视频。} \emph{(Wǒmen \cle{yǒushí} kàn shìpín.)} --- Nous regardons des vidéos \textbf{parfois}.
\end{exemplebox}


\courssub{Écriture : Clés (radicaux) et ordre des traits}

\begin{aretenir}[Analyse de caractères clés de la leçon]
\begin{center}
\begin{tabularx}{\linewidth}{|X|X|X|X|}
\hline
\rowcolor{popblueL} \textbf{Caractère} & \textbf{Radical (Clé)} & \textbf{Nb traits} & \textbf{Ordre des traits (description synthétique)} \\
\hline
\zh{网} & \zh{纟} (sī, soie) & 6 & 1. \zh{纟} (3 traits : point, pli, horizontal) ; 2. \zh{亡} (4 traits : horizontal, horizontal, vertical, point) — \emph{Note : total 6 traits en simplifié (纟=3, 亡=4, mais fusion). Ordre standard : 纟 puis 亡.} \\
\hline
\zh{信} & \zh{亻} (rén, homme) & 9 & 1. \zh{亻} (2 traits : oblique, verticale) ; 2. \zh{言} (7 traits : points, horizontaux, verticale, horizontaux de base) \\
\hline
\zh{电} & \zh{田} (tián, champ) / \zh{电} (caractère unique) & 5 & 1. Horizontal ; 2. Vertical ; 3. Horizontal (cadre \zh{田}) ; 4. Horizontal milieu ; 5. Grand trait oblique final (comme \zh{乙}) \\
\hline
\zh{脑} & \zh{月} (yuè, viande/lune) & 10 & 1. \zh{月} (4 traits : deux obliques, horizontal, horizontal) ; 2. \zh{乃} (6 traits : horizontal, pli, horizontal, vertical, horizontal, point) \\
\hline
\end{tabularx}
\end{center}
\end{aretenir}

\begin{methode}[Méthode : Mémoriser l'écriture]
\begin{enumerate}
    \item Identifier le radical (sémantique souvent) : \zh{纟} (lien, filet), \zh{亻} (humain, message), \zh{月} (organe/corps).
    \item Respecter l'ordre : \textbf{Haut $\to$ Bas, Gauche $\to$ Droite, Horizontal $\to$ Vertical, Extérieur $\to$ Intérieur $\to$ Fermeture}.
    \item S'entraîner sur papier quadrillé (田字格) : 5 répétitions par caractère.
\end{enumerate}
\end{methode}


\courssub{Compréhension du texte et collocations}

\begin{exercice}[Questions de compréhension sur le texte \zh{网络生活}]
Répondez \textbf{en chinois} (caractères + pinyin) en vous basant strictement sur le texte.
\begin{enumerate}
    \item 根据课文，小李住在哪里？他是什么身份？
    \item 小李每天放学后先做什么？然后做什么？(Citez le texte)
    \item 小李周末喜欢做什么？请列举三件事。
    \item 小李为什么控制上网时间？他每天上网多久？
    \item 课文介绍了中国网络生活的哪三个特点？(网购、网课、外卖)
\end{enumerate}
\end{exercice}

\begin{exercice}[Vocabulaire : Associez le chinois, le pinyin et le français]
Complétez le tableau suivant. (Version élève : cases vides ; Version prof : rempli ci-dessous).
\begin{center}
\begin{tabularx}{\linewidth}{|X|X|X|X|}
\hline
\rowcolor{popblueL} \textbf{Caractères} & \textbf{Pinyin} & \textbf{Nature} & \textbf{Traduction} \\
\hline
\zh{邮件} & yóujiàn & n. & Courriel / E-mail \\
\hline
\zh{密码} & mìmǎ & n. & Mot de passe \\
\hline
\zh{网站} & wǎngzhàn & n. & Site web \\
\hline
\zh{下载} & xiàzǎi & v. & Télécharger \\
\hline
\zh{聊天} & liáotiān & VO (v. sep.) & Discuter en ligne / Chatter \\
\hline
\zh{查资料} & chá zīliào & VO & Faire des recherches / Chercher des infos \\
\hline
\zh{发邮件} & fā yóujiàn & VO & Envoyer un e-mail \\
\hline
\zh{登录} & dēnglù & v. & Se connecter / Se logger \\
\hline
\zh{上传} & shàngchuán & v. & Uploader / Télécharger vers serveur \\
\hline
\zh{注册} & zhùcè & v. & S'inscrire / Créer un compte \\
\hline
\zh{方便} & fāngbiàn & adj. & Pratique / Commode \\
\hline
\zh{微信} & Wēixìn & n. & WeChat \\
\hline
\end{tabularx}
\end{center}
\end{exercice}

\begin{exercice}[Pinyin et Tons : Reliez les caractères à la bonne prononciation]
Écrivez le numéro du pinyin correct devant chaque mot. Attention aux tons !
\begin{enumerate}
    \item \zh{电脑} \hspace{1.5cm} \dotfill
    \item \zh{手机} \hspace{1.5cm} \dotfill
    \item \zh{视频} \hspace{1.5cm} \dotfill
    \item \zh{注册} \hspace{1.5cm} \dotfill
    \item \zh{方便} \hspace{1.5cm} \dotfill
    \item \zh{上传} \hspace{1.5cm} \dotfill
\end{enumerate}
\textbf{Liste des pinyin (dans le désordre) :}
\begin{enumerate}
    \item shǒujī
    \item diànnǎo
    \item shàngchuán
    \item zhùcè
    \item fāngbiàn
    \item shìpín
\end{enumerate}
\end{exercice}

\courssub{Je m'entraîne : Production guidée}

\begin{exercice}[Traduction : Français $\rightarrow$ Chinois (Caractères + Pinyin)]
Traduis les phrases suivantes.
\begin{enumerate}
    \item Je me connecte chaque jour avec mon téléphone pour chatter avec mes amis.
    \item S'il vous plaît, envoyez-moi un e-mail demain.
    \item Mon ordinateur est très pratique, je l'utilise souvent pour faire des recherches.
    \item Quel est ton mot de passe ? J'ai oublié le mien.
    \item Nous regardons une vidéo en ligne sur le site web de l'école.
\end{enumerate}
\end{exercice}

\begin{exercice}[Transformation de phrases : Substitution et restructuration]
Modèle : \zh{我用电脑查资料。} (Wǒ yòng diànnǎo chá zīliào.)
\begin{enumerate}
    \item Remplace \zh{电脑} par \zh{手机} et \zh{查资料} par \zh{发邮件}.
    \item Remplace \zh{我} par \zh{他} et \zh{用} par \zh{想用}.
    \item Remplace \zh{查资料} par \zh{看视频} et \zh{电脑} par \zh{平板电脑}.
    \item Transforme à la forme négative (ne pas utiliser... pour...).
    \item Transforme en question avec \zh{吗}.
\end{enumerate}
\end{exercice}

\begin{exercice}[Texte à trous : Collocations et mots-outils]
Complétez avec les mots de la liste. Chaque mot s'utilise \textbf{une seule fois}. Écrivez les caractères.
\begin{center}
\zh{上网 \quad 登录 \quad 密码 \quad 发 \quad 查 \quad 下载 \quad 注册 \quad 微信}
\end{center}
\begin{enumerate}
    \item 我想 \zh{\_\_\_\_\_\_} 一个新账号，但是我忘记了 \zh{\_\_\_\_\_\_}。
    \item 每天晚上，我用电脑 \zh{\_\_\_\_\_\_}，在 \zh{\_\_\_\_\_\_} 上和朋友聊天。
    \item 请 \zh{\_\_\_\_\_\_} 这个软件，它很 \zh{方便}。
    \item 老师让我们 \zh{\_\_\_\_\_\_} 资料，然后写报告。
    \item 我要 \zh{\_\_\_\_\_\_} 一封邮件给校长。
    \item D'abord, il faut \zh{\_\_\_\_\_\_} sur le \zh{网站} de l'école.
\end{enumerate}
\end{exercice}

\begin{exercice}[Caractères : Radicaux et ordre des traits -- Entraînement]
Complétez le tableau pour les caractères ci-dessous.
\begin{center}
\begin{tabularx}{\linewidth}{|X|X|X|X|}
\hline
\rowcolor{popblueL} \textbf{Caractère} & \textbf{Radical (Clé)} & \textbf{Nb traits} & \textbf{Ordre des traits (description)} \\
\hline
\zh{网} & \dotfill & \dotfill & \dotfill \\
\hline
\zh{信} & \dotfill & \dotfill & \dotfill \\
\hline
\zh{电} & \dotfill & \dotfill & \dotfill \\
\hline
\zh{脑} & \dotfill & \dotfill & \dotfill \\
\hline
\end{tabularx}
\end{center}
\end{exercice}

\courssub{Je me prépare au Bac : Compréhension, Traduction, Expression}

\begin{exercice}[Compréhension de texte : L'internet au Cameroun et en Chine]
Lisez le texte suivant (version enrichie HSK 3-4), puis répondez \textbf{en chinois} aux questions. Citez le texte.

\begin{textebox}[Texte Bac : \zh{5G时代的教育机遇} (Les opportunités éducatives à l'ère 5G)]
\zh{随着5G技术的普及，喀麦隆的网速越来越快。现在的学生不仅能在网上查阅百科全书，还能参加线上课程，和北京、上海的老师实时互动。} \\
\emph{Suízhe 5G jìshù de pǔjí, Kāmàilóng de wǎngsù yuè lái yuè kuài. Xiànzài de xuéshēng bùjǐn néng zài wǎngshàng cháyuè bǎikēquánshū, hái néng cānjiā xiànshàng kèchéng, hé Běijīng, Shànghǎi de lǎoshī shíshí hùdòng.} \\
« Avec la popularisation de la technologie 5G, la vitesse du net au Cameroun devient de plus en plus rapide. Les élèves actuels peuvent non seulement consulter des encyclopédies en ligne, mais aussi suivre des cours en ligne, et interagir en temps réel avec des professeurs de Pékin et Shanghai. » \\

\zh{这对于偏远地区的孩子来说，是一个改变命运的机会。但是，网络安全教育也很重要，大家要保护好个人隐私，不轻易相信网上的陌生人。} \\
\emph{Zhè duìyú piānyuǎn dìqū de háizi lái shuō, shì yīgè gǎibiàn mìngyùn de jīhuì. Dànshì, wǎngluò ānquán jiàoyù yě hěn zhòngyào, dàjiā yào bǎohù hǎo gèrén yǐnsī, bù qīngyì xiāngxìn wǎngshàng de mòshēngrén.} \\
« Pour les enfants des zones reculées, c'est une opportunité de changer de destin. Mais l'éducation à la cybersécurité est aussi importante, tous doivent bien protéger leur vie privée, ne pas croire facilement les inconnus sur Internet. »
\end{textebox}

\begin{enumerate}
    \item 根据课文，5G技术给喀麦隆学生带来了哪两个主要好处？(用课文原句词汇回答)
    \item 为什么说这是「改变命运的机会」对于偏远地区的孩子？
    \item 课文提到了哪两点网络安全建议？
    \item Traduisez en français la phrase : \zh{大家要保护好个人隐私，不轻易相信网上的陌生人。}
\end{enumerate}
\end{exercice}

\begin{exercice}[Thème et Version : Traduction spécialisée]

\textbf{Partie A : Version (Chinois $\rightarrow$ Français)}
Traduis le paragraphe ci-dessus (\zh{5G时代的教育机遇}) en français naturel, niveau lycée. Soigne les connecteurs logiques (\zh{随着、不仅...还...、对于...来说、但是}).
\end{exercice}


\textbf{Partie B : Thème (Français $\rightarrow$ Chinois)}
Traduis en chinois (caractères + pinyin). Utilise le vocabulaire de la leçon : \zh{网络, 手机, 方便, 查资料, 视频, 微信, 登录, 密码, 保护, 上网, 聊天, 每天, 用}.
\begin{quote}
« Internet change l'éducation au Cameroun. Beaucoup d'élèves à Buea ou à Garoua ont un smartphone. Ils se connectent au Wi-Fi de l'école pour chercher des informations et regarder des cours en vidéo. C'est très pratique. Mais ils doivent protéger leur mot de passe et ne pas chatter trop longtemps. L'enseignant dit : "Internet est un bon outil, mais il faut l'utiliser correctement." »
\end{quote}


\begin{exercice}[Expression écrite : \zh{「我的网络生活」}]
\textbf{Consigne :} Rédigez un court article (environ \textbf{120 à 150 caractères chinois}) pour le journal du club de chinois, sur le thème \zh{「我的网络生活」}.

\textbf{Points obligatoires :}
\begin{itemize}
    \item Tes outils : \zh{手机} ou \zh{电脑} ? Fréquence (\zh{每天、经常、有时}) ?
    \item Tes activités : \zh{查资料、发邮件、聊天、看视频、下载、网购}...
    \item Ton avis : \zh{方便 / 快 / 有意思} ? Inconvénient ? (\zh{对眼睛不好、浪费时间}) ?
    \item Un conseil ou une résolution.
\end{itemize}

\textbf{Contraintes de langue :}
\begin{itemize}
    \item Au moins 3 structures \zh{用...} (outil).
    \item Au moins 2 collocations VO (\zh{上网、发邮件、查资料、看视频}...).
    \item Ponctuation chinoise correcte (。 ， ！).
    \item Écriture directe en caractères.
\end{itemize}

\begin{center}
\begin{tikzpicture}[scale=0.85, every node/.style={font=\small}]
% Grille 14x8 (approx 112 cases)
\draw[popline, thick] (0,0) rectangle (14,-8);
\foreach \y in {-1,-2,-3,-4,-5,-6,-7} {
    \draw[popline] (0,\y) -- (14,\y);
}
\foreach \x in {1,2,3,4,5,6,7,8,9,10,11,12,13} {
    \draw[popline] (\x,0) -- (\x,-8);
}
% Repères numérotation lignes
\foreach \y/\num in {-0.5/1, -1.5/2, -2.5/3, -3.5/4, -4.5/5, -5.5/6, -6.5/7, -7.5/8} {
    \node[popmuted, anchor=east] at (-0.2,\y) {\num};
}
\node[popmuted, align=center] at (7,-9) {Espace de rédaction : 14 cases $\times$ 8 lignes = 112 cases \\ (Écrire sur deux colonnes si nécessaire)};
\end{tikzpicture}
\end{center}
\end{exercice}

\begin{exemplebox}[Modèle de rédaction (indicatif)]
\zh{我的网络生活} \\
\zh{我每天用手机上网，也经常用电脑查资料做作业。网络很方便，我用微信和同学聊天，看中文视频学中文，有时网购买书。但是上网太久对眼睛不好，也浪费时间。所以我决定：每天只上网两小时，保护好密码，不轻易加陌生人。网络是好工具，要正确使用！} \\
\emph{Wǒ de wǎngluò shēnghuó. Wǒ měitiān yòng shǒujī shàngwǎng, yě jīngcháng yòng diànnǎo chá zīliào zuò zuòyè. Wǎngluò hěn fāngbiàn, wǒ yòng Wēixìn hé tóngxué liáotiān, kàn Zhōngwén shìpín xué Zhōngwén, yǒushí wǎnggòu mǎi shū. Dànshì shàngwǎng tài jiǔ duì yǎnjing bù hǎo, yě làngfèi shíjiān. Suǒyǐ wǒ juédìng: měitiān zhǐ shàngwǎng liǎng xiǎoshí, bǎohù hǎo mìmǎ, bù qīngyì jiā mòshēngrén. Wǎngluò shì hǎo gōngjù, yào zhèngquè shǐyòng!} \\
(~130 caractères. Contient : 3x \zh{用}, 4x VO, avis + inconvénient + résolution.)
\end{exemplebox}

\begin{savaistu}[Culture numérique : Cameroun vs Chine]
\begin{itemize}
    \item \textbf{Paiement mobile :} En Chine, \zh{微信支付} (WeChat Pay) et \zh{支付宝} (Alipay) ont quasi remplacé l'espèce (\zh{现金}) et la carte bancaire. On paie tout par \zh{扫码} (scan QR code) : marché, taxi, métro, factures. Au Cameroun, \zh{Mobile Money} (Orange Money, MTN MoMo) domine pour les transferts et paiements factures, mais le paiement marchand par QR code se développe (ex: \zh{Glotap}, \zh{KmerPay}).
    \item \textbf{Éducation :} La Chine a massivement basculé vers les \zh{网课} (cours en ligne) pendant le Covid et maintient des plateformes nationales (\zh{学习强国}, \zh{雨课堂}). Au Cameroun, l'enseignement à distance (\zh{在线教育}) a progressé mais l'accès inégal à Internet (\zh{数字鸿沟}) reste un défi majeur en zone rurale.
    \item \textbf{Réseaux sociaux :} \zh{微信} (WeChat) est un « super-app » (messagerie + paiement + réseau social + services admin). L'équivalent le plus proche au Cameroun est l'usage combiné de WhatsApp + Mobile Money + Facebook. \zh{抖音} (Douyin) est la version chinoise de TikTok (contenu distinct, censuré).
\end{itemize}
\end{savaistu}

\begin{aretenir}[Check-list de révision Leçon 35]
\begin{itemize}
    \item [ ] Je connais le vocabulaire du tableau (caractères, pinyin, sens, nature).
    \item [ ] Je sais écrire \zh{网、信、电、脑} (radical, nb traits, ordre).
    \item [ ] Je maîtrise la structure \zh{用} + Outil + Verbe (place avant le verbe).
    \item [ ] Je reconnais les verbes séparables : \zh{上网、聊天、发邮件、查资料、下载} et je sais les séparer (ex: \zh{上了网、聊天聊了很久}).
    \item [ ] Je place correctement \zh{每天、经常、有时} (Sujet + Adv + Verbe).
    \item [ ] Je distingue \zh{网络} (nom) et \zh{上网} (verbe).
    \item [ ] Je peux répondre aux questions de compréhension en chinois.
    \item [ ] Je peux traduire un paragraphe Version/Thème sur le thème.
    \item [ ] Je peux rédiger 120-150 caractères sur « Ma vie sur Internet » avec les contraintes demandées.
\end{itemize}
\end{aretenir}

\newpage

\coursec{Corrigés des exercices}

\begin{corrige}[Exercice 1 — Questions de compréhension sur le texte \zh{网络生活}]
\begin{enumerate}
    \item \zh{小李住在雅温得，他是一所高中的学生。} \\
    \emph{Xiǎo Lǐ zhù zài Yǎwēndé, tā shì yī suǒ gāozhōng de xuéshēng.} \\
    « Petit Li habite à Yaoundé, il est élève d'un lycée. » \\
    \emph{Justification :} le texte dit \zh{小李是雅温得一所高中的学生} ; le classificateur des écoles est \zh{所} (suǒ).
    \item \zh{他每天放学后，先回家吃饭，然后用电脑查资料做作业。} \\
    \emph{Tā měitiān fàngxué hòu, xiān huí jiā chīfàn, ránhòu yòng diànnǎo chá zīliào zuò zuòyè.} \\
    « Chaque jour après la classe, il rentre d'abord manger, puis utilise l'ordinateur pour chercher des informations et faire ses devoirs. » \\
    \emph{Point de grammaire :} la chaîne chronologique \zh{先……然后……} (d'abord... ensuite...) est exigée par la question ; l'instrument \zh{用电脑} se place avant le verbe.
    \item \zh{周末的时候，他喜欢：一、用微信和朋友聊天；二、看中国的短视频；三、下载中文歌曲听。} \\
    \emph{Zhōumò de shíhou, tā xǐhuan: yī, yòng Wēixìn hé péngyou liáotiān; èr, kàn Zhōngguó de duǎn shìpín; sān, xiàzǎi Zhōngwén gēqǔ tīng.} \\
    « Le week-end, il aime : 1. chatter avec ses amis sur WeChat ; 2. regarder des vidéos courtes chinoises ; 3. télécharger des chansons chinoises pour les écouter. » \\
    \emph{Point de grammaire :} l'énumération \zh{一、……二、……三、……} avec la virgule chinoise de liste \zh{、}.
    \item \zh{因为上网时间太长，对眼睛不好，所以他每天只上网两个小时。} \\
    \emph{Yīnwèi shàngwǎng shíjiān tài cháng, duì yǎnjing bù hǎo, suǒyǐ tā měitiān zhǐ shàngwǎng liǎng gè xiǎoshí.} \\
    « Parce que passer trop de temps en ligne est mauvais pour les yeux, il ne va sur Internet que deux heures par jour. » \\
    \emph{Point de grammaire :} la paire \zh{因为……所以……} (parce que... donc...) ; \zh{只} (zhǐ, seulement) se place avant le verbe ; \zh{两个小时} : « deux » vaut \zh{两} (et non \zh{二}) devant un classificateur.
    \item \zh{课文介绍了三个特点：网购、网课、外卖。} \\
    \emph{Kèwén jièshào le sān gè tèdiǎn: wǎnggòu, wǎngkè, wàimài.} \\
    « Le texte présente trois caractéristiques : les achats en ligne, les cours en ligne et la livraison de repas. »
\end{enumerate}
\textbf{Points de vigilance :} répondre par des phrases complètes (sujet + verbe), ne pas recopier le texte sans l'adapter à la question ; attention aux tons de \zh{雅温得} (Yǎwēndé, 3\ieme{}--1\ier{}--2\ieme{} ton : \emph{yǎ--wēn--dé}).
\end{corrige}

\begin{corrige}[Exercice 2 — Vocabulaire : association chinois / pinyin / français]
\begin{center}
\begin{tabularx}{\linewidth}{|X|X|X|X|}
\hline
\rowcolor{popblueL} \textbf{Caractères} & \textbf{Pinyin} & \textbf{Nature} & \textbf{Traduction} \\
\hline
\zh{邮件} & yóujiàn & n. & Courriel / e-mail \\
\hline
\zh{密码} & mìmǎ & n. & Mot de passe \\
\hline
\zh{网站} & wǎngzhàn & n. & Site web \\
\hline
\zh{下载} & xiàzǎi & v. & Télécharger (vers son appareil) \\
\hline
\zh{聊天} & liáotiān & VO (v. sep.) & Discuter en ligne / chatter \\
\hline
\zh{查资料} & chá zīliào & VO & Faire des recherches \\
\hline
\zh{发邮件} & fā yóujiàn & VO & Envoyer un e-mail \\
\hline
\zh{登录} & dēnglù & v. & Se connecter (à un compte) \\
\hline
\zh{上传} & shàngchuán & v. & Mettre en ligne, uploader \\
\hline
\zh{注册} & zhùcè & v. & S'inscrire, créer un compte \\
\hline
\zh{方便} & fāngbiàn & adj. & Pratique, commode \\
\hline
\zh{微信} & Wēixìn & n. propre & WeChat \\
\hline
\end{tabularx}
\end{center}
\textbf{Points de vigilance :} distinguer \zh{下载} (xiàzǎi, télécharger vers son appareil, litt. « descendre-charger ») et \zh{上传} (shàngchuán, envoyer vers un serveur, litt. « monter-transmettre ») — erreur classique des francophones, car « télécharger » couvre les deux sens en français. \zh{微信} prend une majuscule en pinyin (nom propre). \zh{聊天} est un verbe séparable : on peut dire \zh{聊了一个小时的天} (chatter pendant une heure).
\end{corrige}

\begin{corrige}[Exercice 3 — Pinyin et tons]
\begin{enumerate}
    \item \zh{电脑} $\to$ diànnǎo — 4\ieme{} ton + 3\ieme{} ton (attention : \zh{脑} n'est PAS au 4\ieme{} ton).
    \item \zh{手机} $\to$ shǒujī — 3\ieme{} ton + 1\ier{} ton.
    \item \zh{视频} $\to$ shìpín — 4\ieme{} + 2\ieme{} ton.
    \item \zh{注册} $\to$ zhùcè — deux 4\ieme{} tons.
    \item \zh{方便} $\to$ fāngbiàn — 1\ier{} + 4\ieme{} ton.
    \item \zh{上传} $\to$ shàngchuán — 4\ieme{} + 2\ieme{} ton.
\end{enumerate}
\textbf{Points de vigilance :} ne pas confondre \zh{视频} (shìpín) et \zh{上传} (shàngchuán) qui commencent tous deux par \emph{sh} ; attention à la différence \emph{sh} (rétroflexe) / \emph{s} dans \zh{手机} (shǒujī) et aux initiales \emph{zh} + \emph{c} de \zh{注册} (zhùcè). Deux 3\ieme{} tons qui se suivent ne posent pas de problème ici, mais rappel : \zh{很好} se prononce \emph{ní hǎo} $\to$ non, \emph{hén hǎo} avec sandhi : le premier 3\ieme{} ton devient un 2\ieme{} ton à l'oral (\emph{hénhǎo}).
\end{corrige}

\begin{corrige}[Exercice 4 — Traduction Français $\to$ Chinois]
\begin{enumerate}
    \item \zh{我每天用手机上网和朋友聊天。} \\
    \emph{Wǒ měitiān yòng shǒujī shàngwǎng hé péngyou liáotiān.} \\
    \emph{Justification :} ordre obligatoire : Sujet + fréquence (\zh{每天}) + instrument (\zh{用} + outil) + verbe. « Avec mon téléphone » ne peut pas se placer après le verbe.
    \item \zh{请明天给我发一封邮件。} \\
    \emph{Qǐng míngtiān gěi wǒ fā yī fēng yóujiàn.} \\
    \emph{Justification :} le complément de temps \zh{明天} se place avant le verbe ; \zh{给} + personne précède le verbe ; classificateur des lettres et e-mails : \zh{封} (fēng).
    \item \zh{我的电脑很方便，我经常用它查资料。} \\
    \emph{Wǒ de diànnǎo hěn fāngbiàn, wǒ jīngcháng yòng tā chá zīliào.} \\
    \emph{Justification :} l'adjectif \zh{方便} fonctionne comme verbe d'état : pas de \zh{是} devant, mais \zh{很} ; \zh{它} (tā, « il » pour les objets) reprend \zh{电脑}.
    \item \zh{你的密码是什么？我忘了。} \\
    \emph{Nǐ de mìmǎ shì shénme? Wǒ wàng le.} \\
    \emph{Justification :} \zh{什么} (quoi) reste à la place du mot interrogé, sans déplacement comme en français ; \zh{了} marque l'accompli de \zh{忘} (avoir oublié). On accepte aussi \zh{我忘了我的密码} (Wǒ wàng le wǒ de mìmǎ).
    \item \zh{我们在学校的网站上看视频。} \\
    \emph{Wǒmen zài xuéxiào de wǎngzhàn shàng kàn shìpín.} \\
    \emph{Justification :} le complément de lieu \zh{在……上} (sur...) se place avant le verbe.
\end{enumerate}
\textbf{Points de vigilance :} la grande difficulté de ce thème est la place des compléments : temps, lieu et instrument se placent tous AVANT le verbe en chinois, alors qu'ils suivent le verbe en français. Schéma à retenir : Sujet + temps + lieu (\zh{在}) + instrument (\zh{用}) + verbe + complément d'objet.
\end{corrige}

\begin{corrige}[Exercice 5 — Transformation de phrases]
\begin{enumerate}
    \item \zh{我用手机发邮件。} \\
    \emph{Wǒ yòng shǒujī fā yóujiàn.} — « J'utilise le téléphone pour envoyer un e-mail. »
    \item \zh{他想用电脑查资料。} \\
    \emph{Tā xiǎng yòng diànnǎo chá zīliào.} — « Il veut utiliser l'ordinateur pour faire des recherches. » \\
    \emph{Justification :} le verbe modal \zh{想} (vouloir) se place avant \zh{用}.
    \item \zh{我用平板电脑看视频。} \\
    \emph{Wǒ yòng píngbǎn diànnǎo kàn shìpín.} — « J'utilise la tablette pour regarder des vidéos. »
    \item \zh{我不用电脑查资料。} \\
    \emph{Wǒ bù yòng diànnǎo chá zīliào.} — « Je n'utilise pas l'ordinateur pour faire des recherches. » \\
    \emph{Justification :} la négation \zh{不} se place devant \zh{用} (c'est l'usage de l'outil qui est nié).
    \item \zh{你用电脑查资料吗？} \\
    \emph{Nǐ yòng diànnǎo chá zīliào ma?} — « Utilises-tu l'ordinateur pour faire des recherches ? » \\
    \emph{Justification :} question oui/non par \zh{吗} final, sans inversion sujet-verbe ; le point d'interrogation chinois \zh{？} est obligatoire.
\end{enumerate}
\textbf{Points de vigilance :} la structure \zh{用} + outil + verbe est invariable ; seuls les modaux (\zh{想、要、可以}), la négation (\zh{不}) et la particule interrogative (\zh{吗}) viennent s'y greffer, toujours en position préverbale.
\end{corrige}

\begin{corrige}[Exercice 6 — Texte à trous : collocations]
\begin{enumerate}
    \item 我想\zh{注册}一个新账号，但是我忘记了\zh{密码}。 \\
    \emph{Wǒ xiǎng zhùcè yī gè xīn zhànghào, dànshì wǒ wàngjì le mìmǎ.} \\
    « Je veux créer un nouveau compte, mais j'ai oublié le mot de passe. »
    \item 每天晚上，我用电脑\zh{上网}，在\zh{微信}上和朋友聊天。 \\
    \emph{Měitiān wǎnshang, wǒ yòng diànnǎo shàngwǎng, zài Wēixìn shàng hé péngyou liáotiān.} \\
    « Chaque soir, je vais sur Internet avec mon ordinateur, je chatte avec mes amis sur WeChat. »
    \item 请\zh{下载}这个软件，它很方便。 \\
    \emph{Qǐng xiàzǎi zhè ge ruǎnjiàn, tā hěn fāngbiàn.} \\
    « Téléchargez ce logiciel, il est très pratique. »
    \item 老师让我们\zh{查}资料，然后写报告。 \\
    \emph{Lǎoshī ràng wǒmen chá zīliào, ránhòu xiě bàogào.} \\
    « Le professeur nous demande de chercher des informations, puis de rédiger un rapport. »
    \item 我要\zh{发}一封邮件给校长。 \\
    \emph{Wǒ yào fā yī fēng yóujiàn gěi xiàozhǎng.} \\
    « Je vais envoyer un e-mail au proviseur. »
    \item 先\zh{登录}学校网站。 \\
    \emph{Xiān dēnglù xuéxiào wǎngzhàn.} \\
    « D'abord, il faut se connecter au site web de l'école. »
\end{enumerate}
\textbf{Points de vigilance :} \zh{查} et \zh{发} sont donnés seuls car ils forment les collocations \zh{查资料} et \zh{发邮件} ; chaque mot de la liste n'est utilisé qu'une fois, ce qui permet de vérifier ses réponses par élimination. Noter la structure \zh{让} + personne + verbe (\zh{老师让我们查资料}, « le professeur nous fait/demande de chercher »).
\end{corrige}

\begin{corrige}[Exercice 7 — Caractères : radicaux et ordre des traits]
\begin{center}
\begin{tabularx}{\linewidth}{|X|X|X|X|}
\hline
\rowcolor{popblueL} \textbf{Caractère} & \textbf{Radical (clé)} & \textbf{Nb traits} & \textbf{Ordre des traits} \\
\hline
\zh{网} & \zh{冂} (cadre) & 6 & 1. vertical gauche ; 2. horizontal-crochet (haut et droite du cadre) ; 3. oblique \zh{ノ} ; 4. point \zh{丶} ; 5. oblique \zh{ノ} ; 6. point \zh{丶}. Règle : extérieur $\to$ intérieur. \\
\hline
\zh{信} & \zh{亻} (rén, homme) & 9 & 1--2. \zh{亻} (oblique, verticale) ; 3. point ; 4--6. trois horizontaux ; 7--9. \zh{口} (vertical, horizontal-pli, horizontal de fermeture). Gauche $\to$ droite, haut $\to$ bas. \\
\hline
\zh{电} & \zh{田} (tián, champ) & 5 & 1. vertical gauche ; 2. horizontal-pli (cadre) ; 3. horizontal du milieu ; 4. vertical du milieu ; 5. vertical-courbe-crochet final \zh{乚} qui dépasse en bas. Extérieur $\to$ intérieur $\to$ trait final. \\
\hline
\zh{脑} & \zh{月} (yuè, chair/corps) & 10 & 1--4. \zh{月} (oblique descendante, horizontal-crochet, deux horizontaux) ; 5. point ; 6. horizontal ; 7--10. \zh{凶} (l'intérieur \zh{乂} puis le cadre \zh{凵}). Gauche $\to$ droite. \\
\hline
\end{tabularx}
\end{center}
\textbf{Points de vigilance :} dans \zh{脑}, la clé \zh{月} est la clé de la « chair » (présente dans les parties du corps : \zh{脸、腿、脚}), pas la « lune » ; dans \zh{信}, on retient l'étymologie : un homme (\zh{亻}) et sa parole (\zh{言}) = le message, la confiance. Dans \zh{电}, le dernier trait \zh{乚} doit dépasser le cadre vers le bas. S'entraîner 5 fois par caractère sur papier quadrillé, un caractère par case.
\end{corrige}

\begin{corrige}[Exercice 8 — Compréhension Bac : \zh{5G时代的教育机遇}]
\begin{enumerate}
    \item \zh{一、学生能在网上查阅百科全书；二、学生能参加线上课程，和北京、上海的老师实时互动。} \\
    \emph{Yī, xuéshēng néng zài wǎngshàng cháyuè bǎikēquánshū; èr, xuéshēng néng cānjiā xiànshàng kèchéng, hé Běijīng, Shànghǎi de lǎoshī shíshí hùdòng.} \\
    « 1. Les élèves peuvent consulter des encyclopédies en ligne ; 2. ils peuvent suivre des cours en ligne et interagir en temps réel avec des professeurs de Pékin et Shanghai. » \\
    \emph{Justification :} la structure parallèle \zh{不仅……还……} (non seulement... mais aussi...) du texte donne exactement les deux bénéfices.
    \item \zh{因为偏远地区的孩子也能上线上课程，和大城市的老师互动，得到好的教育。} \\
    \emph{Yīnwèi piānyuǎn dìqū de háizi yě néng shàng xiànshàng kèchéng, hé dà chéngshì de lǎoshī hùdòng, dédào hǎo de jiàoyù.} \\
    « Parce que les enfants des zones reculées peuvent eux aussi suivre des cours en ligne, interagir avec des professeurs des grandes villes et recevoir une bonne éducation. » \\
    \emph{Justification :} question d'inférence — il faut reformuler : l'accès à distance supprime l'obstacle géographique, d'où « changer de destin ».
    \item \zh{一、保护好个人隐私；二、不轻易相信网上的陌生人。} \\
    \emph{Yī, bǎohù hǎo gèrén yǐnsī; èr, bù qīngyì xiāngxìn wǎngshàng de mòshēngrén.} \\
    « 1. Bien protéger sa vie privée ; 2. ne pas croire facilement les inconnus sur Internet. »
    \item Traduction : « Tout le monde doit bien protéger ses données personnelles, et ne pas accorder facilement sa confiance aux inconnus rencontrés sur Internet. » \\
    \emph{Justification :} \zh{好} après \zh{保护} est un complément de résultat (« protéger correctement, bien ») ; \zh{轻易} (qīngyì) = facilement, à la légère.
\end{enumerate}
\textbf{Barème indicatif (sur 10) :} Q1 : 3 pts (1,5 par bénéfice cité) ; Q2 : 2 pts (inférence pertinente en chinois correct) ; Q3 : 2 pts (1 par conseil) ; Q4 : 3 pts (exactitude 2 pts, qualité du français 1 pt).\\
\textbf{Points de vigilance :} répondre en phrases complètes ; reprendre le vocabulaire du texte (\zh{查阅、线上课程、实时互动、隐私、陌生人}) ; ne pas traduire mot à mot en Q4 (\zh{不轻易相信} = « ne pas croire facilement », pas « ne pas léger croire »).
\end{corrige}

\begin{corrige}[Exercice 9 — Thème et Version]
\textbf{Partie A — Version (modèle) :}\\
« Avec la généralisation de la technologie 5G, la vitesse d'Internet au Cameroun devient de plus en plus rapide. Les élèves d'aujourd'hui peuvent non seulement consulter des encyclopédies en ligne, mais aussi participer à des cours en ligne et interagir en temps réel avec des enseignants de Pékin et de Shanghai. Pour les enfants des régions reculées, c'est une occasion de changer de destin. Cependant, l'éducation à la sécurité en ligne est également très importante : chacun doit bien protéger sa vie privée et ne pas croire trop facilement les inconnus sur Internet. »\\
\emph{Connecteurs rendus :} \zh{随着} = « avec » ; \zh{不仅……还……} = « non seulement... mais aussi... » ; \zh{对于……来说} = « pour... » ; \zh{但是} = « cependant ».

\textbf{Partie B — Thème (modèle) :}\\
\zh{网络改变了喀麦隆的教育。在布埃亚或加鲁阿，很多学生有手机。他们用学校的Wi-Fi上网查资料、看视频课，这很方便。但是他们要保护好自己的密码，不要聊天太久。老师说：「网络是好工具，但是要正确使用。」} \\
\emph{Wǎngluò gǎibiàn le Kāmàilóng de jiàoyù. Zài Bù'āiyà huò Jiālǔ'ā, hěn duō xuéshēng yǒu shǒujī. Tāmen yòng xuéxiào de Wi-Fi shàngwǎng chá zīliào, kàn shìpín kè, zhè hěn fāngbiàn. Dànshì tāmen yào bǎohù hǎo zìjǐ de mìmǎ, bù yào liáotiān tài jiǔ. Lǎoshī shuō: "Wǎngluò shì hǎo gōngjù, dànshì yào zhèngquè shǐyòng."}\\
\emph{Justifications :} « changer » + accompli = \zh{改变了} ; « pour chercher... » = instrument \zh{用} avant le verbe ; « il faut l'utiliser correctement » = \zh{要正确使用} (\zh{正确} adverbe devant \zh{使用}) ; \zh{不要} + verbe = négation de l'impératif (« ne... pas »).\\
\textbf{Barème indicatif (sur 20) :} Version 10 pts (sens global 4, connecteurs 3, qualité du français 3) ; Thème 10 pts (exactitude grammaticale 4, vocabulaire de la leçon 3, caractères corrects 2, ponctuation chinoise 1).\\
\textbf{Points de vigilance :} noms de villes transcrits phonétiquement (\zh{布埃亚} Bù'āiyà pour Buea, \zh{加鲁阿} Jiālǔ'ā pour Garoua) ; ne pas oublier \zh{了} après \zh{改变} ; guillemets chinois \zh{「」} pour la citation.
\end{corrige}

\begin{corrige}[Exercice 10 — Expression écrite : \zh{我的网络生活}]
\textbf{Proposition de rédaction modèle (environ 135 caractères) :}\\
\zh{我的网络生活} \\
\zh{我每天用手机上网，也经常用电脑查资料做作业。网络很方便：我用微信和同学聊天，用电脑看中文视频，有时候也下载中文歌曲。但是，上网太久对眼睛不好，也浪费时间。所以我决定每天只上网两个小时，保护好自己的密码，不轻易相信网上的陌生人。网络是好工具，我们要正确使用它！} \\
\emph{Wǒ de wǎngluò shēnghuó. Wǒ měitiān yòng shǒujī shàngwǎng, yě jīngcháng yòng diànnǎo chá zīliào zuò zuòyè. Wǎngluò hěn fāngbiàn: wǒ yòng Wēixìn hé tóngxué liáotiān, yòng diànnǎo kàn Zhōngwén shìpín, yǒushíhou yě xiàzǎi Zhōngwén gēqǔ. Dànshì, shàngwǎng tài jiǔ duì yǎnjing bù hǎo, yě làngfèi shíjiān. Suǒyǐ wǒ juédìng měitiān zhǐ shàngwǎng liǎng gè xiǎoshí, bǎohù hǎo zìjǐ de mìmǎ, bù qīngyì xiāngxìn wǎngshàng de mòshēngrén. Wǎngluò shì hǎo gōngjù, wǒmen yào zhèngquè shǐyòng tā!}\\
« Ma vie sur Internet. Chaque jour je vais sur Internet avec mon téléphone, et j'utilise souvent l'ordinateur pour chercher des informations et faire mes devoirs. Internet est très pratique : je chatte avec mes camarades sur WeChat, je regarde des vidéos chinoises sur l'ordinateur, et parfois je télécharge des chansons chinoises. Mais passer trop de temps en ligne est mauvais pour les yeux et fait perdre du temps. J'ai donc décidé de ne passer que deux heures par jour en ligne, de bien protéger mon mot de passe et de ne pas croire facilement les inconnus sur Internet. Internet est un bon outil, nous devons l'utiliser correctement ! »

\textbf{Vérification des contraintes :}
\begin{itemize}
    \item Structures \zh{用} + outil : \zh{用手机上网、用电脑查资料、用微信聊天、用电脑看视频} (4 occurrences, 3 exigées).
    \item Collocations VO : \zh{上网、查资料、做作业、聊天、看视频、下载歌曲}.
    \item Outils et fréquence : \zh{手机、电脑、每天、经常、有时候}.
    \item Avis + inconvénient + résolution : \zh{很方便} / \zh{对眼睛不好、浪费时间} / \zh{我决定……}.
    \item Ponctuation chinoise : \zh{。 ， ： ！} correctement employées.
\end{itemize}
\textbf{Barème indicatif (sur 20) :} respect des points obligatoires 6 pts ; structures grammaticales (dont \zh{用} et VO) 6 pts ; exactitude des caractères 4 pts ; ponctuation et présentation 2 pts ; longueur conforme (120--150 caractères) 2 pts.\\
\textbf{Points de vigilance :} ne pas écrire \zh{*我上网用手机} (instrument avant le verbe) ; ne pas confondre \zh{网络} (nom, « Internet ») et \zh{上网} (verbe, « aller sur Internet ») ; un caractère par case, la ponctuation occupe aussi une case.
\end{corrige}

\newpage

\coursec{Fiche bilan}

\begin{popfigure}
\begin{center}\tcbox[colback=popblue,colframe=popblue,coltext=white,arc=9pt,boxsep=4pt,left=6pt,right=6pt]{\bfseries\small 网络 Wǎngluò Internet}\end{center}
\vspace{-2pt}
\begin{tcbraster}[raster columns=3,raster equal height=rows,raster column skip=6pt,raster row skip=6pt,enhanced,blankest]
\begin{tcolorbox}[enhanced,colback=white,colframe=poporange,boxrule=0.9pt,arc=3pt,title={Lexique de base},fonttitle=\bfseries\scriptsize,coltitle=white,colbacktitle=poporange,left=4pt,right=4pt,top=2pt,bottom=2pt,boxsep=1pt,toptitle=1.5pt,bottomtitle=1.5pt,halign=left]
\begin{itemize}[nosep,leftmargin=*,itemsep=1pt,font=\scriptsize]
\item 上网 shàngwǎng aller sur Internet
\item 网站 wǎngzhàn site web
\end{itemize}
\end{tcolorbox}
\begin{tcolorbox}[enhanced,colback=white,colframe=popgreen,boxrule=0.9pt,arc=3pt,title={Matériel},fonttitle=\bfseries\scriptsize,coltitle=white,colbacktitle=popgreen,left=4pt,right=4pt,top=2pt,bottom=2pt,boxsep=1pt,toptitle=1.5pt,bottomtitle=1.5pt,halign=left]
\begin{itemize}[nosep,leftmargin=*,itemsep=1pt,font=\scriptsize]
\item 电脑 diànnǎo ordinateur
\item 手机 shǒujī portable
\end{itemize}
\end{tcolorbox}
\begin{tcolorbox}[enhanced,colback=white,colframe=poppurple,boxrule=0.9pt,arc=3pt,title={Actions},fonttitle=\bfseries\scriptsize,coltitle=white,colbacktitle=poppurple,left=4pt,right=4pt,top=2pt,bottom=2pt,boxsep=1pt,toptitle=1.5pt,bottomtitle=1.5pt,halign=left]
\begin{itemize}[nosep,leftmargin=*,itemsep=1pt,font=\scriptsize]
\item 发邮件 fā yóujiàn envoyer un mail
\item 聊天 liáotiān discuter
\item 查资料 chá zīliào chercher des infos
\end{itemize}
\end{tcolorbox}
\begin{tcolorbox}[enhanced,colback=white,colframe=poppink,boxrule=0.9pt,arc=3pt,title={Avantages},fonttitle=\bfseries\scriptsize,coltitle=white,colbacktitle=poppink,left=4pt,right=4pt,top=2pt,bottom=2pt,boxsep=1pt,toptitle=1.5pt,bottomtitle=1.5pt,halign=left]
\begin{itemize}[nosep,leftmargin=*,itemsep=1pt,font=\scriptsize]
\item 方便 fāngbiàn pratique
\item 快 kuài rapide
\end{itemize}
\end{tcolorbox}
\begin{tcolorbox}[enhanced,colback=white,colframe=popgold,boxrule=0.9pt,arc=3pt,title={Dangers},fonttitle=\bfseries\scriptsize,coltitle=white,colbacktitle=popgold,left=4pt,right=4pt,top=2pt,bottom=2pt,boxsep=1pt,toptitle=1.5pt,bottomtitle=1.5pt,halign=left]
\begin{itemize}[nosep,leftmargin=*,itemsep=1pt,font=\scriptsize]
\item 浪费时间 làngfèi shíjiān perdre du temps
\item 影响学习 yǐngxiǎng xuéxí nuire aux études
\end{itemize}
\end{tcolorbox}
\begin{tcolorbox}[enhanced,colback=white,colframe=popblue!60,boxrule=0.9pt,arc=3pt,title={Structures},fonttitle=\bfseries\scriptsize,coltitle=white,colbacktitle=popblue!60,left=4pt,right=4pt,top=2pt,bottom=2pt,boxsep=1pt,toptitle=1.5pt,bottomtitle=1.5pt,halign=left]
\begin{itemize}[nosep,leftmargin=*,itemsep=1pt,font=\scriptsize]
\item 用 + outil + 做
\item 越来越 + adj.
\end{itemize}
\end{tcolorbox}
\begin{tcolorbox}[enhanced,colback=white,colframe=popgreen!60,boxrule=0.9pt,arc=3pt,title={Clés},fonttitle=\bfseries\scriptsize,coltitle=white,colbacktitle=popgreen!60,left=4pt,right=4pt,top=2pt,bottom=2pt,boxsep=1pt,toptitle=1.5pt,bottomtitle=1.5pt,halign=left]
\begin{itemize}[nosep,leftmargin=*,itemsep=1pt,font=\scriptsize]
\item 纟soie → 网 络
\item 讠parole → 话 说
\end{itemize}
\end{tcolorbox}
\end{tcbraster}

\legende{Carte mentale de la leçon : Internet (网络)}
\end{popfigure}

\begin{bilan}
\end{bilan}

\courssub{Lexique clé à connaître par cœur}
\begin{itemize}
  \item \zh{网络} wǎngluò : Internet, le réseau ; \zh{上网} shàngwǎng : aller sur Internet ; \zh{网站} wǎngzhàn : site web.
  \item \zh{电脑} diànnǎo : ordinateur ; \zh{手机} shǒujī : téléphone portable ; \zh{电子邮件} diànzǐ yóujiàn : courriel.
  \item \zh{发} fā : envoyer ; \zh{收} shōu : recevoir ; \zh{聊天} liáotiān : bavarder (en ligne) ; \zh{查资料} chá zīliào : chercher des informations.
  \item \zh{方便} fāngbiàn : pratique ; \zh{有用} yǒuyòng : utile ; \zh{浪费} làngfèi : gaspiller ; \zh{影响} yǐngxiǎng : influencer, nuire à.
  \item \zh{越来越} yuè lái yuè : de plus en plus ; \zh{网民} wǎngmín : internaute.
\end{itemize}

\courssub{Structures de grammaire à maîtriser}
\begin{center}
\begin{tabularx}{\linewidth}{|X|X|X|}
\hline
\rowcolor{popblueL} Structure & Exemple & Traduction \\
\hline
用 + outil + verbe & \zh{我用手机上网。} \emph{Wǒ yòng shǒujī shàngwǎng.} & Je vais sur Internet avec mon portable. \\
\hline
越来越 + adjectif & \zh{网民越来越多。} \emph{Wǎngmín yuè lái yuè duō.} & Les internautes sont de plus en plus nombreux. \\
\hline
可以 + verbe (possibilité) & \zh{我们可以在网上查资料。} \emph{Wǒmen kěyǐ zài wǎngshang chá zīliào.} & On peut chercher des informations en ligne. \\
\hline
对 + nom + 有/没有 + 影响 & \zh{上网太多对学习有影响。} \emph{Shàngwǎng tài duō duì xuéxí yǒu yǐngxiǎng.} & Trop aller sur Internet nuit aux études. \\
\hline
\end{tabularx}
\end{center}

\courssub{Clés et écriture}
\begin{itemize}
  \item \zh{网} (wǎng) : clé \zh{冂} ; extérieur avant intérieur, on ferme en dernier.
  \item \zh{络} (luò) : clé \zh{纟} (soie, à gauche) ; gauche $\rightarrow$ droite.
  \item \zh{话} (huà) : clé \zh{讠} (parole) ; la clé parole signale un lien avec le langage.
  \item \zh{电} (diàn) : trait vertical central en dernier ; haut $\rightarrow$ bas.
\end{itemize}

\courssub{Méthodes express}
\begin{enumerate}
  \item Lire le texte une première fois sans s'arrêter : repérer le thème et les mots connus.
  \item Souligner les mots du champ lexical d'Internet ; deviner les inconnus grâce aux clés (讠, 纟).
  \item Répondre aux questions en reprenant la structure de la question : \zh{他用什么上网？} $\rightarrow$ \zh{他用手机上网。}
  \item Pour produire : outil (用…) + action + avantage (方便、快) + danger (浪费、影响).
\end{enumerate}


\begin{aretenir}[Checklist avant le Bac]
\begin{itemize}
  \item[$\square$] Je sais écrire de mémoire : \zh{网络、上网、电脑、手机、发、收}.
  \item[$\square$] Je prononce correctement les tons de \zh{wǎngluò} (3-4) et \zh{shǒujī} (3-1).
  \item[$\square$] Je connais le vocabulaire : site web, courriel, internaute, chercher des informations.
  \item[$\square$] J'emploie correctement la structure \zh{用} + outil + verbe.
  \item[$\square$] Je maîtrise \zh{越来越} + adjectif sans ajouter \zh{很}.
  \item[$\square$] Je sais dire un avantage (\zh{方便、有用}) et un danger (\zh{浪费时间、影响学习}) d'Internet.
  \item[$\square$] Je reconnais les clés \zh{讠} (parole) et \zh{纟} (soie) dans les caractères.
  \item[$\square$] Je respecte l'ordre des traits : haut $\rightarrow$ bas, gauche $\rightarrow$ droite.
  \item[$\square$] Je réponds aux questions de compréhension par une phrase complète en chinois.
  \item[$\square$] Je mets la ponctuation chinoise pleine chasse ：，。？
  \item[$\square$] Je sais comparer l'usage d'Internet au Cameroun et en Chine en deux phrases simples.
\end{itemize}
\end{aretenir}

\findecours

\end{document}
